% Produire un Curriculum Vitae (CV) — Français Tle Tech — Cours signé OnBuch+
% Programme officiel MINESEC. Compiler avec : tectonic N06-produire-curriculum-vitae-cv.tex
\newcommand{\DOCMATIERE}{Français}
\newcommand{\DOCNIVEAU}{Tle Tech}
\newcommand{\DOCMODULE}{Français – classe de Terminale technique}
\newcommand{\DOCLECON}{N06}
\newcommand{\DOCTITRE}{Produire un Curriculum Vitae (CV)}
% OnBuch+ — préambule « cours signés » ENRICHI (compilé avec tectonic / XeLaTeX)
% Système visuel « Pop » d'OnBuch+ : fond crème (jamais blanc), bordures encre
% épaisses, ombres « dures » décalées, titres Archivo Black en capitales.
% Version MATHÉMATIQUES : graphes (pgfplots/tikz), boîtes pédagogiques
% (définition, propriété, méthode, attention, le savais-tu, exercice résolu,
% exercice, TP, bilan), page de garde et signature OnBuch+ ancrée en bas.
%
% Macros attendues dans main.tex AVANT \input{preamble} :
%   \DOCMATIERE  \DOCNIVEAU  \DOCMODULE  \DOCLECON (numéro)  \DOCTITRE
\documentclass[11pt]{article}

\usepackage{fontspec}
\usepackage[french]{babel}
\usepackage[a4paper,margin=2cm,top=3.4cm,bottom=2.8cm,headheight=1.4cm,footskip=1.3cm]{geometry}
\usepackage{amsmath,amssymb,amsfonts}
\usepackage{mathtools} % \xmapsto, \coloneqq... souvent utilisés par les modèles
\usepackage{cancel} % simplifications barrées (bilans de forces)
% \abs{z} (module, valeur absolue) et \norm{u} (norme), utilisés par les modèles
\providecommand{\abs}{}\let\abs\relax\DeclarePairedDelimiter{\abs}{\lvert}{\rvert}
\providecommand{\norm}{}\let\norm\relax\DeclarePairedDelimiter{\norm}{\lVert}{\rVert}
\usepackage[table]{xcolor} % [table] : fournit aussi \rowcolors (alternance de lignes)
\usepackage{pagecolor}
\usepackage{tikz}
\usetikzlibrary{arrows.meta,positioning,calc,shapes.geometric,shapes.misc,shapes.arrows,decorations.pathmorphing,decorations.pathreplacing,decorations.markings,patterns,shadows,mindmap,trees,fit,backgrounds,matrix,intersections,angles,quotes}
\usepackage{pgfplots}
\usepackage[version=4]{mhchem} % équations nucléaires \ce{^{235}_{92}U}
\usepackage[european,siunitx]{circuitikz} % schémas électriques
\pgfplotsset{compat=1.18}
\usepgfplotslibrary{fillbetween,groupplots} % aires entre courbes (calcul intégral)
\usepackage{enumitem}
\usepackage{fancyhdr}
% [most] charge toutes les bibliothèques tcolorbox (listings, minted, raster,
% documentation...) et gonfle inutilement la mémoire TeX sur un document long
% (~300 boîtes) : on ne charge que ce qui est réellement utilisé.
\usepackage[skins,breakable]{tcolorbox}
\usepackage{graphicx}
\usepackage{colortbl}
\usepackage{booktabs}
\usepackage{tabularx}
\usepackage{multirow}
\usepackage{diagbox} % case d'en-tête diagonale des tableaux à double entrée
% Colonnes extensibles centrée / gauche / droite (C, L, R), souvent utilisées par les modèles
\newcolumntype{C}{>{\centering\arraybackslash}X}
\newcolumntype{L}{>{\raggedright\arraybackslash}X}
\newcolumntype{R}{>{\raggedleft\arraybackslash}X}
\usepackage{array}
\usepackage{multicol}
\usepackage{siunitx}
% parse-numbers=false : accepte aussi \SI{2,5 \times 10^{-3}}{...}
\sisetup{locale=FR,output-decimal-marker={,},group-minimum-digits=4,parse-numbers=false}
\DeclareSIUnit\ohm{\ensuremath{\symup{\Omega}}} % U+2126 absent de la police
\DeclareSIUnit\division{div} % réglages d'oscilloscope (ms/div, V/div)
\DeclareSIUnit\tr{tr} % tours (vitesse de rotation, ex. \SI{1500}{\tr\per\minute})
\DeclareSIUnit\molar{M} % concentration molaire (ex. \SI{3}{\molar}, \SI{1}{\micro\molar})
\DeclareSIUnit\an{an} % année (le modèle confond parfois avec \year, un compteur TeX) — ex. \SI{5}{\kilo\meter\per\an}
\DeclareSIUnit\FCFA{FCFA} % franc CFA (ex. \SI{500}{\FCFA\per\kilo\gram})
\DeclareSIUnit\degreeBrix{\textdegree Brix} % teneur en sucre d'un sirop/jus (réfractométrie)
\DeclareSIUnit\month{mois} % le modèle utilise parfois \month, un compteur TeX, comme unité de durée
\DeclareSIUnit\microliter{\micro\liter} % le modèle écrit parfois \microliter en un seul token
\DeclareSIUnit\million{millions} % dénombrement (ex. \SI{15}{\million\per\mL} spermatozoïdes)
\usepackage{qrcode}
% \degree : commande fréquemment utilisée par les modèles pour un angle ou une
% température en dehors de \SI{}{} (non fournie par les paquets chargés).
\providecommand{\degree}{^{\circ}}
% \qed (fin de démonstration), utilisé par les modèles sans amsthm
\providecommand{\qed}{\hfill$\square$}
% \barre : double barre des valeurs interdites dans un tableau de variations (tkz-tab)
\providecommand{\barre}{\ensuremath{\|}}
% environnement proof (amsthm non chargé) utilisé par les modèles
\ifdefined\proof\else\newenvironment{proof}[1][Démonstration]{\par\noindent\textit{#1.}\ }{\qed\par}\fi
\usepackage{lastpage}
\usepackage{hyperref}
\hypersetup{hidelinks}

% ============================================================
% Palette « Pop » OnBuch+ — mêmes hex que la classe P (plus_ui.dart)
% ============================================================
\definecolor{poporange}{HTML}{F59321}
\definecolor{popdark}{HTML}{E07A0C}
\definecolor{popcream}{HTML}{FFF6E8}
\definecolor{popink}{HTML}{1C1714}
\definecolor{poppurple}{HTML}{7A5AE0}
\definecolor{popgreen}{HTML}{2E9E6A}
\definecolor{poppink}{HTML}{E0457B}
\definecolor{popblue}{HTML}{2F7DE1}
\definecolor{popgold}{HTML}{C98A12}
\definecolor{popmuted}{HTML}{8A7F76}
\definecolor{popline}{HTML}{EADFD2}
\definecolor{popgray}{HTML}{8A7F76}
\colorlet{popgrey}{popgray}
% Alias tolérants (le modèle nomme parfois les couleurs différemment)
\colorlet{popwhite}{white}
\colorlet{popblack}{popink}
\colorlet{popred}{poppink}
\colorlet{popyellow}{popgold}
% Teintes claires (fonds de tableaux/encadrés) — le modèle nomme parfois
% les couleurs avec un suffixe L (ex. popblueL) sans les avoir définies.
\colorlet{poporangeL}{poporange!20}
\colorlet{popdarkL}{popdark!20}
\colorlet{poppurpleL}{poppurple!20}
\colorlet{popgreenL}{popgreen!20}
\colorlet{poppinkL}{poppink!20}
\colorlet{popblueL}{popblue!20}
\colorlet{popgoldL}{popgold!20}
\colorlet{popredL}{popred!20}
\colorlet{popgrayL}{popgray!20}
% Teintes claires pour fonds de tableaux / zones
\colorlet{popblueL}{popblue!10!white}
\colorlet{poporangeL}{poporange!14!white}
\colorlet{popgreenL}{popgreen!12!white}
\colorlet{poppurpleL}{poppurple!12!white}
\colorlet{poppinkL}{poppink!12!white}
\colorlet{popgoldL}{popgold!14!white}

\pagecolor{popcream}

% ============================================================
% Typographie
% ============================================================
\defaultfontfeatures{Ligatures=TeX}
\setmainfont{PlusJakartaSans-Medium.ttf}[Path=../../../fonts/,BoldFont=PlusJakartaSans-Bold.ttf]
% Maths en sans-serif (Fira Math) avec chiffres et lettres droites de Plus
% Jakarta, pour que les formules se fondent dans le texte.
\usepackage[math-style=ISO,bold-style=ISO]{unicode-math}
\setmathfont{FiraMath-Regular.otf}[Path=../../../fonts/]
\setmathfont{PlusJakartaSans-Medium.ttf}[Path=../../../fonts/,range={up/num,up/{Latin,latin},\mathup}]
\usepackage{icomma} % virgule décimale sans espace en mode math
% Caractères mathématiques Unicode absents des polices de texte (Plus Jakarta,
% Archivo Black) : rendus par leur équivalent mathématique, en texte comme en
% formule — sinon ils s'affichent en carré vide (ex. « Arithmétique dans ℤ »).
\usepackage{newunicodechar}
\newunicodechar{ℝ}{\ensuremath{\mathbb{R}}}
\newunicodechar{ℤ}{\ensuremath{\mathbb{Z}}}
\newunicodechar{ℂ}{\ensuremath{\mathbb{C}}}
\newunicodechar{ℕ}{\ensuremath{\mathbb{N}}}
\newunicodechar{ℚ}{\ensuremath{\mathbb{Q}}}
\newunicodechar{𝔻}{\ensuremath{\mathbb{D}}}
\newunicodechar{α}{\ensuremath{\alpha}}
\newunicodechar{β}{\ensuremath{\beta}}
\newunicodechar{γ}{\ensuremath{\gamma}}
\newunicodechar{δ}{\ensuremath{\delta}}
\newunicodechar{ε}{\ensuremath{\varepsilon}}
\newunicodechar{θ}{\ensuremath{\theta}}
\newunicodechar{λ}{\ensuremath{\lambda}}
\newunicodechar{μ}{\ensuremath{\mu}}
\newunicodechar{σ}{\ensuremath{\sigma}}
\newunicodechar{τ}{\ensuremath{\tau}}
\newunicodechar{φ}{\ensuremath{\varphi}}
\newunicodechar{ψ}{\ensuremath{\psi}}
\newunicodechar{ω}{\ensuremath{\omega}}
\newunicodechar{Σ}{\ensuremath{\Sigma}}
% majuscules produites par \MakeUppercase dans les titres (α-aminés -> Α)
\newunicodechar{Α}{\ensuremath{\alpha}}
\newunicodechar{Β}{\ensuremath{\beta}}
\newunicodechar{Γ}{\ensuremath{\Gamma}}
\newunicodechar{Δ}{\ensuremath{\Delta}}
\newunicodechar{Ε}{\ensuremath{\varepsilon}}
\newunicodechar{Θ}{\ensuremath{\Theta}}
\newunicodechar{Λ}{\ensuremath{\Lambda}}
\newunicodechar{Μ}{\ensuremath{\mu}}
\newunicodechar{Τ}{\ensuremath{\tau}}
\newunicodechar{Φ}{\ensuremath{\Phi}}
\newunicodechar{Ψ}{\ensuremath{\Psi}}
\newunicodechar{Ω}{\ensuremath{\Omega}}
\newunicodechar{↦}{\ensuremath{\mapsto}}
\newunicodechar{∈}{\ensuremath{\in}}
\newunicodechar{∉}{\ensuremath{\notin}}
\newunicodechar{∧}{\ensuremath{\wedge}}
\newunicodechar{⇔}{\ensuremath{\Leftrightarrow}}
\newunicodechar{⟺}{\ensuremath{\Longleftrightarrow}}
\newunicodechar{⇒}{\ensuremath{\Rightarrow}}
\newunicodechar{⟹}{\ensuremath{\Longrightarrow}}
\newunicodechar{∘}{\ensuremath{\circ}}
\newunicodechar{∪}{\ensuremath{\cup}}
\newunicodechar{∩}{\ensuremath{\cap}}
\newunicodechar{≡}{\ensuremath{\equiv}}
\newunicodechar{⊂}{\ensuremath{\subset}}
\newunicodechar{⊥}{\ensuremath{\perp}}
\newunicodechar{∃}{\ensuremath{\exists}}
\newunicodechar{∀}{\ensuremath{\forall}}
\newunicodechar{∛}{\ensuremath{\sqrt[3]{\,}}}
\newunicodechar{∜}{\ensuremath{\sqrt[4]{\,}}}
\newunicodechar{⃗}{\ensuremath{{}^{\to}}}
\newunicodechar{ⁿ}{\ifmmode^{n}\else\textsuperscript{\ensuremath{n}}\fi}
\newunicodechar{ˣ}{\ifmmode^{x}\else\textsuperscript{\ensuremath{x}}\fi}
\newunicodechar{ᵇ}{\ifmmode^{b}\else\textsuperscript{\ensuremath{b}}\fi}
\newunicodechar{ʳ}{\ifmmode^{r}\else\textsuperscript{\ensuremath{r}}\fi}
\newunicodechar{ᵘ}{\ifmmode^{u}\else\textsuperscript{\ensuremath{u}}\fi}
\newunicodechar{ʸ}{\ifmmode^{y}\else\textsuperscript{\ensuremath{y}}\fi}
\newunicodechar{ᵃ}{\ifmmode^{a}\else\textsuperscript{\ensuremath{a}}\fi}
\newunicodechar{ᵉ}{\ifmmode^{e}\else\textsuperscript{\ensuremath{e}}\fi}
\newunicodechar{ᵏ}{\ifmmode^{k}\else\textsuperscript{\ensuremath{k}}\fi}
\newunicodechar{ᵖ}{\ifmmode^{p}\else\textsuperscript{\ensuremath{p}}\fi}
\newunicodechar{ᴺ}{\ifmmode^{N}\else\textsuperscript{\ensuremath{N}}\fi}
\newunicodechar{ᶜ}{\ifmmode^{c}\else\textsuperscript{\ensuremath{c}}\fi}
\newunicodechar{ʰ}{\ifmmode^{h}\else\textsuperscript{\ensuremath{h}}\fi}
\newunicodechar{ᵠ}{\ifmmode^{q}\else\textsuperscript{\ensuremath{q}}\fi}
\newunicodechar{ₙ}{\ifmmode_{n}\else\textsubscript{\ensuremath{n}}\fi}
\newunicodechar{ₐ}{\ifmmode_{a}\else\textsubscript{\ensuremath{a}}\fi}
\newunicodechar{ᵢ}{\ifmmode_{i}\else\textsubscript{\ensuremath{i}}\fi}
\newunicodechar{ᵣ}{\ifmmode_{r}\else\textsubscript{\ensuremath{r}}\fi}
\newunicodechar{ₖ}{\ifmmode_{k}\else\textsubscript{\ensuremath{k}}\fi}
\newunicodechar{ᵦ}{\ifmmode_{\beta}\else\textsubscript{\ensuremath{\beta}}\fi}
\newunicodechar{ₓ}{\ifmmode_{x}\else\textsubscript{\ensuremath{x}}\fi}
\newfontfamily\popdisplay{ArchivoBlack-Regular.ttf}[Path=../../../fonts/]
\newfontfamily\poplogo{SpaceGrotesk-Bold.ttf}[Path=../../../fonts/]
\newfontfamily\popsemi{PlusJakartaSans-SemiBold.ttf}[Path=../../../fonts/]
\newfontfamily\popxbold{PlusJakartaSans-ExtraBold.ttf}[Path=../../../fonts/]
\newfontfamily\popxboldls{PlusJakartaSans-ExtraBold.ttf}[Path=../../../fonts/,LetterSpace=3.0]

\setlist[enumerate]{leftmargin=1.7em,itemsep=5pt,topsep=4pt}
\setlist[itemize]{leftmargin=1.7em,itemsep=4pt,topsep=4pt,label={\color{poporange}\textbullet}}
\setlength{\parindent}{0pt}
\setlength{\parskip}{5pt}
\renewcommand{\arraystretch}{1.25}

% pgfplots : style Pop par défaut
\pgfplotsset{
  every axis/.append style={axis line style={popink,line width=1pt},tick style={popink},
    grid style={popline,line width=0.5pt},label style={font=\small},tick label style={font=\footnotesize}},
  popcurve/.style={poporange,line width=1.8pt,no markers,smooth},
}
% Même style disponible dans un \draw TikZ simple (hors pgfplots)
\tikzset{popcurve/.style={poporange,line width=1.8pt}}
\tikzset{popgrid/.style={popline,very thin}}
\colorlet{popcurve}{poporange} % le nom du style est parfois utilisé comme couleur

% ============================================================
% Pill / squiggle
% ============================================================
\newcommand{\poppill}[3]{%
  \tikz[baseline=-0.55ex]\node[draw=popink,line width=1.2pt,rounded corners=2.6pt,%
    fill=#2,inner xsep=6.5pt,inner ysep=3pt]{%
    \popxboldls\fontsize{7}{7}\selectfont\color{#3}\MakeUppercase{#1}};%
}
\newcommand{\popsquiggle}[1]{%
  \tikz[baseline=-0.4ex]{
    \draw[#1,line width=2.6pt,line cap=round]
      plot [smooth] coordinates {(0,0.09) (0.34,0.01) (0.68,0.21) (1.02,0.01) (1.36,0.10)};
  }%
}
% Mot-clé surligné (marqueur orange sous le texte)
\newcommand{\cle}[1]{{\popxbold\color{popdark}#1}}

% ============================================================
% En-tête / pied
% ============================================================
\pagestyle{fancy}
\fancyhf{}
\renewcommand{\headrulewidth}{0pt}
\renewcommand{\footrulewidth}{0pt}
\lhead{\raisebox{-0.15ex}{\poplogo\fontsize{13}{13}\selectfont\color{popink}OnBuch\color{poporange}+}}
\rhead{\poppill{\DOCMATIERE\ \textbullet\ \DOCNIVEAU\ \textbullet\ Leçon \DOCLECON}{white}{popink}}
\lfoot{\popxbold\fontsize{7.6}{7.6}\selectfont\color{popmuted}\MakeUppercase{Cours signé OnBuch+}\ \textbullet\ {\popsemi\DOCTITRE}}
\rfoot{\tikz[baseline=-0.6ex]\node[rounded corners=8.5pt,fill=popink,minimum height=17pt,minimum width=17pt,inner xsep=5pt,inner sep=0]{\popxbold\fontsize{8}{8}\selectfont\color{popcream}\thepage\,/\,\pageref*{LastPage}};}

% ============================================================
% Titres
% ============================================================
\newcommand{\chaptitre}[2]{%
  \begin{tcolorbox}[enhanced,colback=poporange,colframe=popink,boxrule=2.6pt,arc=11pt,
      drop shadow=popink,left=15pt,right=15pt,top=12pt,bottom=11pt,before skip=3pt,after skip=18pt]
    \poppill{#2}{popink}{popcream}\par\vspace{9pt}
    {\raggedright\hyphenpenalty=10000\exhyphenpenalty=10000\popdisplay\fontsize{22}{24}\selectfont\color{popcream}\MakeUppercase{#1}\par}
  \end{tcolorbox}
}
% Section numérotée : pastille orange avec le numéro + titre Archivo
\newcounter{popsec}
\newcommand{\coursec}[1]{%
  \stepcounter{popsec}\setcounter{popsub}{0}%
  \par\vspace{16pt}\noindent
  \begin{minipage}{\linewidth}
  \tikz[baseline=-0.7ex]\node[draw=popink,line width=1.6pt,fill=poporange,rounded corners=4pt,minimum size=22pt,inner sep=0]{\popdisplay\fontsize{12}{12}\selectfont\color{popcream}\thepopsec};%
  \hspace{8pt}{\popdisplay\fontsize{15}{17}\selectfont\color{popink}\MakeUppercase{#1}}\par
  \vspace{4pt}\noindent{\color{poporange}\rule{36pt}{3.6pt}}
  \end{minipage}\par\nopagebreak\vspace{8pt}
}
\newcounter{popsub}
\newcommand{\courssub}[1]{%
  \stepcounter{popsub}%
  \par\vspace{9pt}\noindent{\popxbold\fontsize{11.5}{13}\selectfont\color{popink}{\color{poporange}\thepopsec.\thepopsub}\hspace{6pt}#1}\par\nopagebreak\vspace{4pt}
}

% ============================================================
% Boîtes pédagogiques — même recette : fond blanc, bordure épaisse colorée,
% ombre dure encre, badge plein. Titre optionnel [#1].
% ============================================================
\tcbset{popbase/.style={breakable,enhanced,colback=white,boxrule=2.3pt,arc=9pt,
  shadow={3pt}{-3pt}{0pt}{popink},left=14pt,right=14pt,top=11pt,bottom=11pt,before skip=11pt,after skip=15pt,
  fontupper=\popsemi\fontsize{10.6}{14.5}\selectfont\color{popink},
  fontlower=\popsemi\fontsize{10.6}{14.5}\selectfont\color{popink}}}
% Coche dessinée (le glyphe ✓ n'existe pas dans les polices du document)
\newcommand{\popcheck}[1]{\tikz[baseline=-0.5ex]\draw[#1,line width=1.6pt,line cap=round,line join=round](-0.13,0.0)--(-0.04,-0.09)--(0.13,0.1);}
\providecommand{\coche}{\popcheck{popgreen}}
\providecommand{\resultat}[1]{\par\smallskip\noindent\fcolorbox{popgreen}{popgreenL}{\parbox{\dimexpr\linewidth-2\fboxsep-2\fboxrule\relax}{\textbf{Résultat.} #1}}\par\smallskip}
\newcommand{\popboxhead}[3]{\poppill{#1}{#2}{white}\ifx\relax#3\relax\else\hspace{6pt}{\popxbold\fontsize{10.6}{12}\selectfont\color{popink}#3}\fi\par\vspace{7pt}}

\newtcolorbox{aretenir}[1][]{popbase,colframe=popgreen,before upper={\popboxhead{À retenir}{popgreen}{#1}}}
\newtcolorbox{exemplebox}[1][]{popbase,colframe=poporange,before upper={\popboxhead{Exemple}{poporange}{#1}}}
\newtcolorbox{definition}[1][]{popbase,colframe=popblue,before upper={\popboxhead{Définition}{popblue}{#1}}}
\newtcolorbox{propriete}[1][]{popbase,colframe=poppurple,before upper={\popboxhead{Propriété}{poppurple}{#1}}}
\newtcolorbox{methode}[1][]{popbase,colframe=popgold,before upper={\popboxhead{Méthode}{popgold}{#1}}}
\newtcolorbox{attention}[1][]{popbase,colframe=poppink,before upper={\popboxhead{Attention}{poppink}{#1}}}
\newtcolorbox{experience}[1][]{popbase,colframe=popblue,colback=popblueL,before upper={\popboxhead{Expérience}{popblue}{#1}}}
% « Le savais-tu ? » : carte violette pleine (contexte, culture, Cameroun)
\newtcolorbox{savaistu}[1][]{popbase,colback=poppurple,colframe=popink,
  fontupper=\popsemi\fontsize{10.4}{14.2}\selectfont\color{white},
  before upper={\poppill{Le savais-tu\,?}{white}{poppurple}\ifx\relax#1\relax\else\hspace{6pt}{\popxbold\color{white}#1}\fi\par\vspace{7pt}}}
% Exercice résolu : énoncé en haut, \tcblower puis la solution sur fond crème
\newcounter{popexo}\newcounter{popexr}
\newtcolorbox{exoresolu}[1][]{popbase,colframe=poporange,colbacklower=popcream,
  segmentation style={popink,line width=1.2pt,dashed},
  before upper={\stepcounter{popexr}\popboxhead{Exercice résolu \thepopexr}{poporange}{#1}},
  before lower={\poppill{Solution}{popink}{popcream}\par\vspace{6pt}}}
% Exercice (non résolu en place) — la correction est dans la partie Corrigés
\newtcolorbox{exercice}[1][]{popbase,colframe=popink,
  before upper={\stepcounter{popexo}\popboxhead{Exercice \thepopexo}{popink}{#1}}}
% Corrigé
\newtcolorbox{corrige}[1][]{popbase,colframe=popgreen,colback=popgreenL,
  before upper={\popboxhead{Corrigé}{popgreen}{#1}}}
% Objectifs / prérequis / bilan
\newtcolorbox{objectifs}{popbase,colframe=popink,colback=poporangeL,
  before upper={\popboxhead{Objectifs de la leçon}{popink}{}}}
\newtcolorbox{prerequis}{popbase,colframe=popink,colback=popblueL,
  before upper={\popboxhead{Prérequis}{popblue}{}}}
\newtcolorbox{bilan}{popbase,colframe=popink,colback=white,boxrule=2.8pt,
  before upper={\popboxhead{Fiche bilan}{poporange}{L'essentiel à connaître pour l'examen}}}
% Figure encadrée avec légende
\newtcolorbox{popfigure}[1][]{enhanced,breakable=false,colback=white,colframe=popline,boxrule=1.4pt,arc=8pt,
  left=10pt,right=10pt,top=10pt,bottom=8pt,before skip=10pt,after skip=12pt,halign=center,
  lower separated=false,halign lower=center,
  fontlower=\popsemi\fontsize{9}{11.5}\selectfont\color{popmuted},
  #1}
\newcommand{\legende}[1]{\par\vspace{6pt}{\popsemi\fontsize{9}{11.5}\selectfont\color{popmuted}#1}}

% ============================================================
% Signature OnBuch+
% ============================================================
\newcommand{\poplogoinline}[1]{{\poplogo\fontsize{#1}{#1}\selectfont\color{popink}OnBuch\color{poporange}+}}
\newcommand{\signatureonbuch}{%
  \begin{tcolorbox}[enhanced,colback=popink,colframe=popink,boxrule=2.2pt,arc=10pt,
      drop shadow=poporange,left=14pt,right=14pt,top=10pt,bottom=10pt,before skip=0pt,after skip=0pt]
    \tikz[baseline=-0.7ex]\node[circle,fill=popgreen,draw=popcream,line width=1.3pt,minimum size=22pt,inner sep=0]{\popcheck{white}};%
    \hspace{9pt}\begin{minipage}[c]{0.62\linewidth}
      {\popxbold\fontsize{12}{13}\selectfont\color{popcream}Signé\ \textbullet\ {\poplogo OnBuch\color{poporange}+}}\par\vspace{2pt}
      {\raggedright\popsemi\fontsize{8}{10.5}\selectfont\color{popline}\MakeUppercase{Équipe pédagogique OnBuch+}\\\MakeUppercase{Conforme au programme officiel MINESEC}\par}
    \end{minipage}\hfill
    {\poplogo\fontsize{22}{22}\selectfont\color{popcream}OnBuch\color{poporange}+}
  \end{tcolorbox}%
}

% ============================================================
% Page de garde
% #1 titre · #2 matière · #3 niveau · #4 module · #5 n° leçon · #6 durée ·
% #7 description / objectifs (texte ou itemize)
% ============================================================
\newcommand{\pagegarde}[7]{%
  \thispagestyle{empty}
  \newgeometry{a4paper,margin=1.7cm,top=1.6cm,bottom=1.4cm}%
  \begin{tikzpicture}[remember picture,overlay]
    \fill[poporange!18] ([xshift=-3.2cm,yshift=-3.6cm]current page.north east) circle (4.6cm);
    \fill[poppurple!14] ([xshift=2.4cm,yshift=7.6cm]current page.south west) circle (3.8cm);
  \end{tikzpicture}%
  {\poplogo\fontsize{30}{30}\selectfont\color{popink}OnBuch\color{poporange}+}\hfill
  \raisebox{0.6ex}{\poppill{Cours signé \textbullet\ Premium}{popink}{popcream}}\par
  \vspace{1pt}\popsquiggle{poppurple}\par
  \vspace{8pt}
  \begin{tcolorbox}[enhanced,colback=poporange,colframe=popink,boxrule=2.8pt,arc=13pt,
      drop shadow=popink,left=18pt,right=18pt,top=12pt,bottom=13pt,after skip=10pt]
    \poppill{#2 \textbullet\ #3}{popink}{popcream}\hspace{5pt}\poppill{#4}{white}{popink}\par\vspace{8pt}
    {\popdisplay\fontsize{14}{14}\selectfont\color{popink}LEÇON #5}\par\vspace{4pt}
    {\raggedright\hyphenpenalty=10000\exhyphenpenalty=10000\popdisplay\fontsize{25}{27}\selectfont\color{popcream}\MakeUppercase{#1}\par}
    \vspace{8pt}
    {\popxbold\fontsize{9.5}{9.5}\selectfont\color{popink}\MakeUppercase{Durée conseillée : #6}}
  \end{tcolorbox}
  \begin{tcolorbox}[enhanced,colback=white,colframe=popink,boxrule=2.2pt,arc=10pt,
      drop shadow=popink,left=16pt,right=16pt,top=10pt,bottom=10pt,after skip=8pt]
    \poppill{Dans cette leçon}{poppurple}{white}\par\vspace{5pt}
    \popsemi\fontsize{9.3}{12.6}\selectfont\color{popink}#7
  \end{tcolorbox}
  \noindent\begin{minipage}[t]{0.487\linewidth}
    \begin{tcolorbox}[enhanced,colback=popgreen,colframe=popink,boxrule=2.2pt,arc=10pt,
        drop shadow=popink,left=11pt,right=11pt,top=10pt,bottom=10pt,height=3.1cm,valign=center]
      \tcbox[colback=white,colframe=popink,boxrule=1.3pt,arc=5pt,boxsep=0pt,nobeforeafter,
        left=3pt,right=3pt,top=3pt,bottom=3pt,valign=center]{\qrcode[height=1.9cm]{https://play.google.com/store/apps/details?id=com.luvvix.onbuch&hl=fr}}%
      \hspace{8pt}\begin{minipage}[c]{0.5\linewidth}
        {\popdisplay\fontsize{11}{12.5}\selectfont\color{white}\MakeUppercase{Télécharge OnBuch}}\par\vspace{4pt}
        {\raggedright\popsemi\fontsize{8.4}{11}\selectfont\color{white}Gratuit \textbullet\ Google Play\\Tuteur IA, annales, concours, résultats\par}
      \end{minipage}
    \end{tcolorbox}
  \end{minipage}\hfill
  \begin{minipage}[t]{0.487\linewidth}
    \begin{tcolorbox}[enhanced,colback=poppurple,colframe=popink,boxrule=2.2pt,arc=10pt,
        drop shadow=popink,left=11pt,right=11pt,top=10pt,bottom=10pt,height=3.1cm,valign=center]
      \tikz[baseline=-0.7ex]\node[circle,fill=white,draw=popink,line width=1.3pt,minimum size=1.6cm,inner sep=0]{\popdisplay\fontsize{24}{24}\selectfont\color{poppurple}+};%
      \hspace{8pt}\begin{minipage}[c]{0.6\linewidth}
        {\popdisplay\fontsize{11}{12.5}\selectfont\color{white}\MakeUppercase{Passe à OnBuch+}}\par\vspace{4pt}
        {\raggedright\popsemi\fontsize{8.4}{11}\selectfont\color{white}Tous les cours signés, Tuteur IA illimité, fiches PDF — onglet OnBuch+ de l'app\par}
      \end{minipage}
    \end{tcolorbox}
  \end{minipage}
  \vspace{8pt}
  \popcontenu
  \vfill
  \signatureonbuch
  \restoregeometry
  \newpage
}

% Rangée « Ce cours contient » (page de garde)
\newcommand{\poptile}[3]{% #1 couleur #2 gros texte #3 légende
  \begin{tcolorbox}[enhanced,colback=#1,colframe=popink,boxrule=2pt,arc=9pt,shadow={2.5pt}{-2.5pt}{0pt}{popink},
      left=6pt,right=6pt,top=9pt,bottom=9pt,halign=center,valign=center,height=1.75cm,width=\linewidth,nobeforeafter]
    {\popdisplay\fontsize{12.5}{13}\selectfont\color{white}\MakeUppercase{#2}}\par\vspace{3pt}
    {\centering\popsemi\fontsize{7.8}{9.5}\selectfont\color{white}#3\par}
  \end{tcolorbox}}
\newcommand{\popcontenu}{%
  \par\noindent{\popxboldls\fontsize{7.5}{7.5}\selectfont\color{popmuted}CE COURS SIGNÉ CONTIENT}\par\vspace{7pt}
  \noindent\begin{minipage}[t]{0.235\linewidth}\poptile{popblue}{Cours}{complet, illustré, pas à pas}\end{minipage}\hfill
  \begin{minipage}[t]{0.235\linewidth}\poptile{poporange}{Méthodes}{et exercices résolus}\end{minipage}\hfill
  \begin{minipage}[t]{0.235\linewidth}\poptile{poppink}{Exercices}{type examen + corrigés}\end{minipage}\hfill
  \begin{minipage}[t]{0.235\linewidth}\poptile{popgreen}{Fiche bilan}{l'essentiel à retenir}\end{minipage}\par}

% Sommaire
\newtcolorbox{sommaire}{enhanced,colback=white,colframe=popink,boxrule=2.2pt,arc=10pt,
  drop shadow=popink,left=18pt,right=18pt,top=13pt,bottom=13pt,before skip=6pt,after skip=20pt,
  before upper={\poppill{Sommaire}{poppurple}{white}\par\vspace{9pt}\popsemi\fontsize{10.6}{14.5}\selectfont\color{popink}}}

% Signature de fin de cours — poussée tout en bas de la dernière page
\newcommand{\findecours}{%
  \par\vfill\nopagebreak
  \begin{center}\popsquiggle{poporange}\end{center}\vspace{4pt}
  \signatureonbuch
}
\AtBeginDocument{\def\llbracket{\mathopen{[\mkern-3.5mu[}}\def\rrbracket{\mathclose{]\mkern-3.5mu]}}} % intervalles d'entiers (glyphe ⟦ absent de la police)
% Glyphes absents de Fira Math : redéfinis à partir de glyphes disponibles
\AtBeginDocument{%
  \def\setminus{\mathbin{\backslash}}\let\smallsetminus\setminus
  \def\vdots{\mathord{\vbox{\baselineskip3.2pt\lineskiplimit0pt\kern4pt\hbox{.}\hbox{.}\hbox{.}}}}%
  \def\ddots{\mathinner{\mkern1mu\raise6.5pt\hbox{.}\mkern2mu\raise3.6pt\hbox{.}\mkern2mu\raise0.7pt\hbox{.}\mkern1mu}}%
  \def\triangle{\mathord{\scalebox{0.9}{$\symup{\Delta}$}}}%
  \def\nabla{\mathord{\raisebox{\depth}{\scalebox{1}[-1]{$\symup{\Delta}$}}}}%
  \def\checkmark{\popcheck{popdark}}%
  \def\star{\mathbin{\ast}}\def\bigstar{\mathord{\ast}}%
  \def\diamond{\mathbin{\scalebox{0.6}{$\Diamond$}}}%
  \def\bigcup{\mathop{\mathchoice{\raisebox{-0.3em}{\scalebox{1.7}{$\cup$}}}{\scalebox{1.25}{$\cup$}}{\cup}{\cup}}\displaylimits}%
  \def\bigcap{\mathop{\mathchoice{\raisebox{-0.3em}{\scalebox{1.7}{$\cap$}}}{\scalebox{1.25}{$\cap$}}{\cap}{\cap}}\displaylimits}%
  \def\lesseqqgtr{\mathrel{\lessgtr}}\def\bigoplus{\mathop{\oplus}\displaylimits}%
}
\providecommand{\ssi}{si et seulement si} % abréviation fréquente
\AtBeginDocument{\providecommand{\wideparen}{\overparen}} % arc de cercle

%% FIGFIX-BEGIN v4 — correctifs globaux des figures (docs/onbuch-generation/tools/figfix)
\usepackage{adjustbox}\usepackage{etoolbox}
% toute figure TikZ/pgfplots plus large que la ligne est réduite à la largeur disponible
\BeforeBeginEnvironment{tikzpicture}{\begin{adjustbox}{max width=\linewidth}}
\AfterEndEnvironment{tikzpicture}{\end{adjustbox}}
% légendes pgfplots lisibles par-dessus les courbes
\pgfplotsset{every axis/.append style={legend style={fill=white,fill opacity=0.92,text opacity=1,draw=black!15,inner sep=3pt}}}
% étiquettes d'arêtes (syntaxe edge["..."]) avec fond blanc
\tikzset{every edge quotes/.append style={fill=white,inner sep=1.2pt}}
%% FIGFIX-END
\begin{document}

\pagegarde{Produire un Curriculum Vitae (CV)}{Français}{Tle Tech}{Français – classe de Terminale technique}{N06}{11 séances (fiche de progression harmonisée)}{Cette leçon outille l'élève de Terminale technique pour lire, analyser et rédiger un Curriculum Vitae adapté à une situation d'embauche réelle. À partir de modèles authentiques et d'offres d'emploi camerounaises, elle conduit progressivement de l'observation de textes iconiques à la production et à l'amélioration d'un CV personnel.
\vspace{4pt}
\begin{multicols}{2}\small
\begin{itemize}
\item À la fin de cette leçon, je sais définir le CV et identifier sa fonction dans une situation d'embauche
\item À la fin de cette leçon, je sais décrire la structure externe d'un CV (mise en page, rubriques, photo, longueur)
\item À la fin de cette leçon, je sais analyser la structure interne d'un CV (état civil, formation, expérience, compétences, centres d'intérêt)
\item À la fin de cette leçon, je sais comparer plusieurs modèles de CV (chronologique, fonctionnel, mixte) et choisir le plus adapté
\item À la fin de cette leçon, je sais analyser une offre d'emploi pour en dégager le profil recherché
\item À la fin de cette leçon, je sais rédiger un CV complet, clair et sans faute, adapté à une offre précise
\end{itemize}\end{multicols}}

\begin{sommaire}
\begin{multicols}{2}
\begin{enumerate}
\item Découverte du CV : définition et fonctions
\item La structure externe du CV
\item La structure interne du CV : les rubriques
\item Analyser la situation d'embauche
\item Rédaction et amélioration du CV
\item Activité d'intégration
\item Méthodes et exercices résolus
\item Exercices
\item Corrigés des exercices
\item Fiche bilan
\end{enumerate}
\end{multicols}
\end{sommaire}

\begin{objectifs}
\begin{itemize}
\item À la fin de cette leçon, je sais définir le CV et identifier sa fonction dans une situation d'embauche
\item À la fin de cette leçon, je sais décrire la structure externe d'un CV (mise en page, rubriques, photo, longueur)
\item À la fin de cette leçon, je sais analyser la structure interne d'un CV (état civil, formation, expérience, compétences, centres d'intérêt)
\item À la fin de cette leçon, je sais comparer plusieurs modèles de CV (chronologique, fonctionnel, mixte) et choisir le plus adapté
\item À la fin de cette leçon, je sais analyser une offre d'emploi pour en dégager le profil recherché
\item À la fin de cette leçon, je sais rédiger un CV complet, clair et sans faute, adapté à une offre précise
\item À la fin de cette leçon, je sais adapter et améliorer mon CV en fonction d'une situation d'embauche donnée
\item À la fin de cette leçon, je sais justifier mes choix de rédaction (ordre des rubriques, vocabulaire, informations retenues)
\end{itemize}
\end{objectifs}

\begin{prerequis}
\begin{itemize}
\item La lettre de motivation (vue en classe de Première)
\item Les types de textes et la communication écrite administrative
\item La correspondance administrative (demande d'emploi)
\item Le réseau internet et la recherche d'informations (usage du téléphone et de l'ordinateur)
\item La relecture et la correction orthographique et grammaticale
\end{itemize}
\end{prerequis}

\newpage

\coursec{Découverte du CV : définition et fonctions}

À Douala, une entreprise de BTP publie une offre d'emploi de technicien en génie civil : 350 candidatures arrivent en une semaine, mais le recruteur ne retient que 12 CV pour les entretiens. Pourquoi certains CV sont-ils écartés en moins de 30 secondes de lecture ? Qu'est-ce qui distingue un CV efficace d'un CV rejeté ? Cette leçon nous apprend à produire un document professionnel capable de convaincre un employeur camerounais ou international.

\courssub{Observation guidée du texte iconique N°1 : un CV authentique}

Commençons par l'analyse concrète d'un document réel. Imaginez le CV de \textbf{Jean-Pierre M.}, jeune titulaire d'un BTS en Génie Civil option Bâtiment, obtenu à l'Institut Supérieur de Technologie de Douala. Ce document, présenté sur une seule page A4, suit une mise en page aérée, structurée en blocs distincts. En haut, l'en-tête affiche son identité, ses coordonnées (téléphone, email professionnel, localisation : Douala, Akwa) et une photo d'identité récente, sobre, sur fond neutre. Viennent ensuite quatre rubriques principales : \emph{Formation}, \emph{Expériences professionnelles \& Stages}, \emph{Compétences techniques \& Transversales}, \emph{Langues \& Centres d'intérêt}.

Dans \emph{Formation}, il liste son BTS (2023), son Probatoire Série F4 (2020) et une certification « Sécurité sur chantier » (2022). Les dates sont alignées à gauche, les intitulés à droite. La rubrique \emph{Expériences} détaille deux stages de fin d'études (3 mois chacun) chez \emph{Entreprise Bâtiment Plus} et \emph{SOGEA Satom} : pour chacun, il précise le poste (Stagiaire chef de chantier adjoint), la période, puis 3 à 4 puces d'actions concrètes commençant par des verbes d'action (\emph{suivi}, \emph{contrôlé}, \emph{participé}, \emph{rédigé}) et chiffrant les résultats (\emph{suivi de 5 corps d'état}, \emph{contrôle de 200 m³ de béton}). Les \emph{Compétences} séparent les techniques (lecture de plans AutoCAD/Revit, calcul de ferraillage, gestion de planning MS Project) des transversales (esprit d'équipe, rigueur, communication). Enfin, \emph{Langues} indique Français (maternel), Anglais (niveau B2 - TOEIC 785) et Centres d'intérêt : football (capitaine d'équipe), bénévolat (encadrement jeunesse).

\begin{popfigure}
\begin{tikzpicture}[
    node distance=0.8cm and 1.2cm,
    box/.style={rectangle, draw=popblue, thick, rounded corners=4pt, fill=popblueL, text width=3.2cm, align=center, minimum height=1.6cm, font=\small},
    arrow/.style={-Stealth, thick, poporange},
    labelarrow/.style={font=\tiny, text=popdark, fill=popcream, rounded corners=2pt, inner sep=1.5pt}
]
% Main CV Block
\node[box, minimum height=10cm, text width=12cm] (cv) {};
% Header
\node[box, minimum height=1.8cm, text width=11cm, anchor=north] (header) at (cv.north) {\textbf{EN-TÊTE}\\\emph{Identité, Photo, Coordonnées, Titre visé}};
% Formation
\node[box, anchor=north west] (form) at ($(header.south west)+(0,-0.3)$) {\textbf{FORMATION}\\\emph{Diplômes, Certifications, Dates, Lieux}};
% Expériences
\node[box, anchor=north west] (exp) at ($(form.south west)+(0,-0.3)$) {\textbf{EXPÉRIENCES}\\\emph{Stages, Emplois, Missions, Résultats chiffrés}};
% Compétences
\node[box, anchor=north west] (comp) at ($(exp.south west)+(0,-0.3)$) {\textbf{COMPÉTENCES}\\\emph{Techniques (outils, méthodes) + Transversales (soft skills)}};
% Langues / Centres d'intérêt
\node[box, anchor=north west] (lang) at ($(comp.south west)+(0,-0.3)$) {\textbf{LANGUES \& INTÉRÊTS}\\\emph{Niveaux certifiés, Engagements associatifs, Sports}};

% Annotations with arrows
\draw[arrow] (header.east) -- ++(1.5,0) node[labelarrow, anchor=west, align=center]{Identification immédiate\\+ Titre du poste visé};
\draw[arrow] (form.east) -- ++(1.5,0) node[labelarrow, anchor=west, align=center]{Parcours scolaire\\du + récent au + ancien};
\draw[arrow] (exp.east) -- ++(1.5,0) node[labelarrow, anchor=west, align=center]{Preuves par l'action\\Verbes + Chiffres + Contexte};
\draw[arrow] (comp.east) -- ++(1.5,0) node[labelarrow, anchor=west, align=center]{Adéquation poste/profil\\Mots-clés de l'offre};
\draw[arrow] (lang.east) -- ++(1.5,0) node[labelarrow, anchor=west, align=center]{Ouverture, personnalité\\Différenciation};
\end{tikzpicture}
\legende{Structure annotée d'un CV modèle (BTS Génie Civil) : chaque rubrique répond à une attente précise du recruteur.}
\end{popfigure}

\noindent
\textbf{Consignes d'observation :}
\begin{enumerate}
    \item Repérez l'ordre des rubriques : pourquoi la formation précède-t-elle l'expérience ici ? (Réponse : jeune diplômé, la formation est l'atout principal).
    \item Identifiez les verbes d'action dans la rubrique Expériences. Quel temps verbal domine ? (Passé composé ou infinitif pour décrire des missions achevées).
    \item Notez la présence de chiffres (200 m³, 5 corps d'état, TOEIC 785). Quelle fonction remplissent-ils ? (Preuve, objectivation, crédibilité).
    \item Observez la mise en page : alignements, polices, espaces blancs. Quel effet sur la lisibilité ? (Repérage visuel rapide, confort de lecture).
\end{enumerate}

\begin{exoresolu}[Analyse critique d'un extrait de CV]
\textbf{Énoncé :} Voici deux façons de présenter la même expérience de stage. Laquelle est la plus efficace pour un recruteur et pourquoi ?

\textit{Version A :} « J'ai fait un stage chez Bâtiment Plus. J'ai regardé les plans. J'ai aidé le chef de chantier. C'était intéressant. »

\textit{Version B :} « Stage Chef de chantier adjoint — Entreprise Bâtiment Plus, Douala (01/2023–03/2023). Suivi quotidien de l'avancement de 3 corps d'état (gros œuvre, maçonnerie, électricité) sur un immeuble R+3. Contrôle de la conformité des ferraillages selon plans EXE (environ 15 t d'acier). Rédaction de 12 rapports journaliers d'avancement transmis au conducteur de travaux. »

\tcblower
\textbf{Solution détaillée :}
La \textbf{Version B} est professionnelle, la Version A est narrative et vague.
\begin{itemize}
    \item \textbf{En-tête contextuel :} Version B donne le titre exact, l'entreprise, le lieu, les dates (repérage temporel immédiat).
    \item \textbf{Verbes d'action :} « Suivi », « Contrôle », « Rédaction » (infinitif ou passé composé) remplacent « J'ai fait », « J'ai regardé », « J'ai aidé ». Le candidat se positionne en \cle{acteur}.
    \item \textbf{Chiffres et précision :} « 3 corps d'état », « R+3 », « 15 t d'acier », « 12 rapports ». Le recruteur mesure l'\cle{ampleur} et la \cle{nature} des responsabilités.
    \item \textbf{Mots-clés techniques :} « plans EXE », « ferraillages », « conducteur de travaux » prouvent la maîtrise du vocabulaire métier.
    \item \textbf{Résultat :} Version B permet au recruteur de cocher les cases « autonomie », « technique », « communication écrite » en 10 secondes.
\end{itemize}
\end{exoresolu}

\courssub{Définition et étymologie}

Le terme \cle{Curriculum Vitae} (souvent abrégé \cle{CV}) trouve son origine dans le latin : \emph{curriculum} (course, carrière, parcours) et \emph{vitae} (de la vie, génitif de \emph{vita}). Littéralement, c'est la « \cle{course de la vie} » ou le « \cle{parcours de vie} ». Historiquement, chez les Romains, le \emph{cursus honorum } désignait la succession des charges publiques d'un citoyen. Le sens moderne conserve cette idée de \cle{parcours structuré} et \cle{chronologique}, mais restreint au domaine professionnel et académique.

\begin{definition}[Curriculum Vitae (CV)]
Le \textbf{Curriculum Vitae (CV)} est un document écrit, bref (idéalement 1 page pour un jeune diplômé, 2 pages max pour un expérimenté) et structuré, qui présente de manière \cle{synthétique} et \cle{sélective} l'identité, la formation, l'expérience professionnelle, les compétences techniques et transversales, ainsi que les centres d'intérêt d'une personne, en vue d'obtenir un emploi, un stage, une admission en formation ou un marché.
\end{definition}

\noindent
Trois adjectifs qualifient le CV idéal :
\begin{itemize}
    \item \textbf{Synthétique :} Il condense l'essentiel sans phrases superflues. Pas de « Je suis né le... », « J'ai grandi à... ».
    \item \textbf{Sélectif :} On ne met pas toute sa vie. On choisit ce qui est \cle{pertinent} par rapport à la cible (offre d'emploi, concours).
    \item \textbf{Structuré :} Des rubriques standardisées permettent au recruteur de trouver l'information attendue au bon endroit (ex : dates à gauche, lieux à droite).
\end{itemize}

\courssub{Les fonctions du CV}

Le CV n'est pas une fin en soi ; c'est un \cle{outil de marketing personnel}. Il remplit trois fonctions majeures, articulées dans le temps du recrutement :

\begin{enumerate}
    \item \textbf{Informer (Fonction informative) :} Fournir des \cle{faits vérifiables} : diplômes obtenus, postes occupés, périodes, compétences techniques maîtrisées, langues parlées. C'est la base factuelle.
    \item \textbf{Convaincre (Fonction persuasive) :} Au-delà des faits, le CV doit \cle{démontrer l'adéquation} entre le profil du candidat et les besoins de l'employeur. C'est ici que le choix des mots, la hiérarchisation de l'information et la mise en valeur des résultats (chiffres, projets réussis) opèrent la persuasion.
    \item \textbf{Obtenir un entretien (Fonction opérationnelle) :} Le but ultime du CV n'est pas d'obtenir l'emploi, mais l'\cle{entretien d'embauche}. C'est le « sésame » qui ouvre la porte. S'il ne déclenche pas l'envie de rencontrer le candidat, il a échoué.
\end{enumerate}

\begin{aretenir}[La règle des 30--40 secondes]
Des études en psychologie du travail et en recrutement (eye-tracking) montrent qu'un recruteur consacre en \textbf{moyenne 30 à 40 secondes} à la \textbf{première lecture} (scan) d'un CV. Il ne « lit » pas, il \cle{scanne} : il cherche des mots-clés, des noms d'écoles, des titres de poste, des dates, des chiffres. \textbf{Conséquence :} L'information clé doit sauter aux yeux (mise en gras, positionnement stratégique, puces). Un CV dense, mal aligné, sans mots-clés visibles, est mis de côté avant même d'être lu.
\end{aretenir}

\courssub{Distinction entre CV et lettre de motivation}

C'est une confusion classique chez les élèves. Le CV et la lettre de motivation (ou lettre d'accompagnement) sont \cle{complémentaires} mais de \cle{nature différente}.

\begin{center}
\begin{tabularx}{\linewidth}{|l|X|X|}
\hline
\rowcolor{popblueL}
\textbf{Critère} & \textbf{Curriculum Vitae (CV)} & \textbf{Lettre de motivation} \\
\hline
\textbf{Nature} & Document \cle{factuel}, \cle{objectif}, \cle{standardisé} & Document \cle{argumentatif}, \cle{subjectif}, \cle{personnalisé} \\
\hline
\textbf{Contenu} & \emph{Quoi ?} Identité, Diplômes, Expériences, Compétences, Dates. & \emph{Pourquoi ?} Motivation, Projet pro, Adéquation profil/poste, Connaissance de l'entreprise. \\
\hline
\textbf{Structure} & Rubriques fixes, liste à puces, tableau. & Paragraphes enchaînés (Vous, Moi, Nous), style épistolaire. \\
\hline
\textbf{Temps verbal} & Principalement infinitif (rubriques) ou passé composé (missions). & Présent, futur, conditionnel (projection). \\
\hline
\textbf{Personnalisation} & Adaptation \emph{interne} (mots-clés, ordre rubriques) pour chaque offre. & Réécriture \emph{totale} pour chaque offre (adresse au destinataire, référence à l'entreprise). \\
\hline
\textbf{Question clé} & « Est-ce que ce candidat a les compétences techniques ? » & « Est-ce que ce candidat veut \emph{vraiment} ce poste \emph{ici} et \emph{maintenant} ? » \\
\hline
\end{tabularx}
\end{center}

\begin{attention}[Piège fréquent : Confondre CV et Autobiographie]
\textbf{Erreur :} Rédiger son CV comme un récit de vie chronologique détaillé : « Je suis né en 2004 à Bafoussam, j'ai fait l'école primaire à..., puis le collège à..., j'ai eu mon Probatoire en 2020, ensuite je suis allé à Douala pour mon BTS... »
\textbf{Correction :} Le CV est un \cle{document professionnel}, pas un journal intime. On ne raconte pas sa vie, on \cle{présente son profil}. Supprimez : « Je suis né le... », « J'ai fait mes classes primaires à... », « Mes parents sont... », « Je suis célibataire » (sauf exigence légale rare), « J'aime lire et voyager » (trop vague). Gardez : « 2020--2023 : BTS Génie Civil (IST Douala) », « Stage : Suivi chantier R+3 ». \textbf{Le CV est une affiche publicitaire, pas un roman.}
\end{attention}

\courssub{Contextes d'usage au Cameroun}

Au Cameroun, le CV est un passage obligé dans de nombreuses situations professionnelles et administratives. Maîtriser ses codes est une compétence de \cle{insertion professionnelle} majeure.

\begin{exemplebox}[Contextes camerounais d'utilisation du CV]
\begin{itemize}
    \item \textbf{Concours de la Fonction Publique :} ENAM (École Nationale d'Administration et de Magistrature), ENSET (École Normale Supérieure de l'Enseignement Technique), concours d'entrée aux IFORD, IRIC, etc. Le CV y est souvent obligatoire dans le dossier de candidature, accompagné des \cle{diplômes certifiés conformes}, extraits d'acte de naissance, certificat de nationalité, casier judiciaire. La rigueur administrative y est extrême : une rature, une date erronée, un diplôme non certifié entraînent le rejet pur et simple.
    \item \textbf{Emplois privés (Entreprises, PME, Grands Groupes) :} Société Générale, SABC, Brasseries, MTN, Orange, Bolloré, entreprises du BTP (Razel, MAG, etc.), cabinets d'études. Le CV est envoyé par email (PDF) ou déposé sur plateformes (LinkedIn, Emploi.cm, Jobstore.cm). La réactivité et la mise en page moderne (sobriété, lisibilité écran) sont primordiales.
    \item \textbf{Stages académiques et professionnels :} Stage de fin d'études (BTS, HND, Licence Pro, Master), stage d'insertion (programme PIAASI, PAJER-U). Le CV « étudiant » met l'accent sur la formation, les projets scolaires, les compétences techniques acquises en TP, la disponibilité.
    \item \textbf{ONG et Organisations Internationales :} ONU (PNUD, UNICEF, OMS), ONG internationales (MSF, Care, Plan International), coopération (GIZ, AFD). Exigence forte : CV en \cle{format Europass} ou format ONU (PHP - Personal History Profile), souvent en \cle{anglais} et français, avec description détaillée des responsabilités (niveau de supervision, budget géré, résultats).
    \item \textbf{Marchés publics et Consultances :} Appels d'offres pour prestations intellectuelles (études d'impact, contrôle de travaux, formation). Le CV des experts clés (chef de mission, experts sectoriels) est noté techniquement (grille de notation : diplômes, années d'expérience, références similaires). La \cle{véracité} est contractuelle (faux = fraude, exclusion des marchés).
\end{itemize}
\end{exemplebox}

\courssub{Le circuit du CV : du candidat au recruteur}

Comprendre le parcours physique et numérique du CV aide à en soigner la forme. Le CV voyage : il doit résister aux impressions, aux scans, aux transfers emails, aux logiciels de lecture automatique (ATS -- \emph{Applicant Tracking Systems}).

\begin{popfigure}
\begin{tikzpicture}[
    node distance=1.8cm and 1.5cm,
    proc/.style={rectangle, draw=popblue, thick, rounded corners=6pt, fill=popblueL, text width=2.8cm, align=center, minimum height=1.4cm, font=\small\bfseries},
    doc/.style={rectangle, draw=popgreen, thick, rounded corners=6pt, fill=popgreenL, text width=2.8cm, align=center, minimum height=1.4cm, font=\small},
    decision/.style={diamond, draw=poporange, thick, aspect=1.5, fill=poporangeL, text width=2.5cm, align=center, minimum height=1.2cm, font=\small\bfseries},
    arrow/.style={-Stealth, thick, popdark},
    label/.style={font=\tiny, text=popmuted, sloped, above, pos=0.5}
]
% Nodes
\node[proc, align=center] (candidat){CANDIDAT\\Rédaction, Mise en page,\\Conversion PDF};
\node[doc, right=of candidat, align=center] (cv){CV NUMÉRIQUE (PDF)\\\& Version Papier};
\node[proc, right=of cv, align=center] (envoi){ENVOI\\Email, Plateforme,\\Dépôt physique, Réseau};
\node[decision, below=of envoi, align=center] (ats){LOGICIEL ATS\\Filtrage automatique\\Mots-clés, Format};
\node[proc, below=of ats, align=center] (recruteur){RECRUTEUR / DRH\\Lecture rapide (30--40s)\\Scan visuel};
\node[decision, right=of recruteur, align=center] (preselect){PRÉSÉLECTION\\OUI / NON};
\node[proc, below=of preselect, align=center] (entretien){ENTRETIEN\\Échange oral,\\Vérification compétences};
\node[proc, below=of entretien, align=center] (embauche){EMBAUCHE / STAGE\\Contrat, Intégration};

% Arrows
\draw[arrow] (candidat) -- node[label] {Création} (cv);
\draw[arrow] (cv) -- node[label] {Transmission} (envoi);
\draw[arrow] (envoi) -- node[label, right] {Dépôt} (ats);
\draw[arrow] (ats) -- node[label, right] {Compatible ?} (recruteur);
\draw[arrow] (recruteur) -- node[label] {Intéressant ?} (preselect);
\draw[arrow] (preselect) -- node[label, right] {OUI} (entretien);
\draw[arrow] (entretien) -- node[label] {Convaincant ?} (embauche);

% Rejection loops (dashed)
\draw[arrow, dashed, poporange] (ats.west) -- ++(-1,0) |- node[label, left, align=center]{Rejet : Format illisible,\\Mots-clés absents} (candidat.south);
\draw[arrow, dashed, poporange] (preselect.west) -- ++(-1.5,0) |- node[label, left, align=center]{Rejet : Profil inadéquat,\\Présentation soignée ?} (recruteur.south);
\draw[arrow, dashed, poporange] (embauche.west) -- ++(-1,0) |- node[label, left] {Échec entretien} (entretien.south);
\end{tikzpicture}
\legende{Le circuit du CV : de la rédaction à l'embauche. Les boucles pointillées orange montrent les points d'élimination où la qualité du CV est décisive.}
\end{popfigure}

\noindent
Ce schéma révèle trois \cle{points de passage critiques} où le CV seul décide :
\begin{enumerate}
    \item \textbf{Le filtre ATS (Applicant Tracking System) :} Utilisé par les grandes entreprises et l'administration. Il « lit » le texte du PDF. Si votre CV est une image (scan), des colonnes complexes, ou sans mots-clés de l'offre, il est \textbf{rejeté automatiquement}. \emph{Astuce :} Utilisez du texte sélectionnable, police standard (Arial, Calibri), mots-clés exacts de l'annonce.
    \item \textbf{Le scan visuel du recruteur (30--40 s) :} L'humain ne lit que ce qui saute aux yeux : Titres de poste, Entreprises, Dates, Diplômes, Mots-clés en gras. Une mise en page \cle{respirante} (marges, interlignes) guide l'œil.
    \item \textbf{L'entretien :} Le CV devient le \cle{script} de l'entretien. Le recruteur pose des questions \emph{sur} ce qui est écrit (« Vous dites "gestion de planning", quel logiciel ? »). Toute incohérence ou exagération se paie cash.
\end{enumerate}

\begin{methode}[Check-list « Mon CV passe-t-il les 3 filtres ? »]
Avant d'envoyer votre CV, vérifiez :
\begin{enumerate}
    \item \textbf{Format :} Fichier \textbf{PDF} (pas Word, pas image), texte sélectionnable, 1 page (jeune diplômé), nom du fichier pro : \texttt{Nom\_Prenom\_PosteVisé.pdf}.
    \item \textbf{Mots-clés :} Les termes techniques de l'offre (ex: \emph{AutoCAD}, \emph{MS Project}, \emph{DAO}, \emph{Génie Civil}) apparaissent-ils dans vos rubriques Compétences / Expériences ?
    \item \textbf{Lisibilité 30s :} Imprimez-le. En 30 secondes, un tiers retrouve-t-il : Votre nom ? Le poste visé ? Votre dernier diplôme ? Votre dernière expérience ? Vos 3 compétences fortes ?
    \item \textbf{Cohérence :} Dates continues ? Pas de trous inexpliqués ? Orthographe irréprochable (relue par un tiers) ?
    \item \textbf{Véracité :} Tout est justifiable (diplômes, stages, niveaux de langue) lors de l'entretien.
\end{enumerate}
\end{methode}

\begin{savaistu}[Le CV « Europass » et le marché international]
Le format \textbf{Europass} (créé par l'UE, gratuit sur \texttt{europa.eu/europass}) est un standard reconnu dans 30+ pays. Au Cameroun, il est souvent exigé pour les stages Erasmus+, les bourses DAAD, les recrutements GIZ/UE/ONU. Il impose une structure rigide (Profil, Expérience, Formation, Compétences numériques, Linguistiques, etc.) et génère un PDF standardisé. \textbf{Avantage :} Lisible par tous les ATS européens. \textbf{Inconvénient :} Peu différenciant, visuellement lourd, pas idéal pour le secteur privé camerounais qui préfère un CV « design » sur une page. \emph{Conseil :} Ayez les deux versions prêtes : une version « Design 1 page » pour le privé local, une version « Europass » pour les appels d'offres internationaux et l'administration.
\end{savaistu}

\noindent
\emph{Transition :} Maintenant que nous avons défini l'objet, compris ses fonctions, distingué ses compagnons (lettre de motivation) et visualisé son parcours jusqu'au décideur, nous allons entrer dans l'atelier : \textbf{comment construire la structure externe (mise en page, police, format) puis la structure interne (rubriques, rédaction, verbes d'action) d'un CV qui survive aux 30 secondes ?} C'est l'objet de la section 2.

\coursec{La structure externe du CV}

Dans la section précédente, nous avons découvert ce qu'est un \cle{curriculum vitae} : un document qui résume le parcours d'un candidat et dont la fonction est de décrocher un entretien d'embauche. Mais un CV ne se juge pas seulement à son contenu : avant même de lire une seule ligne, le recruteur \emph{voit} le document. Cette première impression visuelle, c'est ce que nous appelons la \cle{structure externe} du CV : sa longueur, sa mise en page, sa typographie, sa photo, son support. C'est l'objet de cette section.

\courssub{Observation du texte iconique N°2 : deux CV, un même profil}

\begin{experience}[Observer pour comprendre : deux CV de mise en page différente]
\textbf{Matériel} : deux CV (texte iconique N°2) présentant le \emph{même} candidat — mêmes diplômes, mêmes expériences — mais mis en page différemment.

\textbf{Protocole} : distribuer les deux documents à la classe. Chronométrer 30 secondes de lecture pour chacun, puis demander aux élèves de noter ce qu'ils ont retenu.

\textbf{Observations} :
\begin{itemize}
\item Sur le CV A (texte dense, trois polices, cadres multiples, paragraphes longs), les élèves retiennent peu d'informations et peinent à repérer le diplôme principal.
\item Sur le CV B (une page aérée, rubriques titrées, puces, une seule police), les élèves repèrent immédiatement l'état civil, la formation et l'expérience.
\end{itemize}

\textbf{Interprétation} : à contenu identique, c'est la \cle{mise en page} qui détermine la quantité d'informations retenue par le lecteur. Un recruteur qui reçoit des dizaines de CV consacre en moyenne moins d'une minute à chacun : un document illisible est écarté avant même d'être lu.
\end{experience}

Cette observation \textbf{fonde} toute la section : la structure externe n'est pas un décor, c'est une \emph{condition d'efficacité} du message.

\courssub{La longueur : une page}

\begin{definition}[Longueur du CV]
Le CV d'un jeune candidat doit tenir sur \cle{une seule page}. Deux pages sont tolérées uniquement pour un cadre expérimenté, dont le parcours professionnel est long et riche.
\end{definition}

Pourquoi cette règle ? Le recruteur lit vite. Une page oblige le candidat à \emph{sélectionner} l'essentiel : c'est déjà un exercice de contraction du texte, que vous pratiquez en séquence de résumé. Un CV de trois pages pour un élève de Terminale qui postule à un stage trahit un défaut de synthèse.

\begin{attention}[Piège : vouloir tout dire]
Multiplier les pages pour « tout mettre » (toutes les classes fréquentées, tous les petits travaux) est une erreur fréquente. Le CV n'est pas une autobiographie : c'est une \emph{sélection} d'informations utiles à l'employeur. Ce qui ne sert pas la candidature doit disparaître.
\end{attention}

\courssub{La mise en page : aération, marges, alignement, hiérarchie}

La mise en page repose sur quatre principes simples.

\begin{itemize}
\item \textbf{L'aération} : le blanc n'est pas de l'espace perdu, c'est un repos pour l'œil. Un CV « tassé » fatigue ; un CV aéré invite à la lecture. Il faut laisser de l'espace entre les rubriques.
\item \textbf{Les marges} : des marges régulières (environ 2 cm de chaque côté) encadrent le texte et donnent une impression d'ordre et de sérieux.
\item \textbf{L'alignement} : le texte est aligné à gauche ; les dates sont alignées dans une même colonne. On évite de centrer les paragraphes, ce qui brouille les repères.
\item \textbf{La hiérarchie visuelle} : le nom du candidat est le plus visible ; viennent ensuite les titres de rubriques, puis le contenu. Le regard doit descendre naturellement du haut vers le bas.
\end{itemize}

\begin{popfigure}
\begin{tikzpicture}[scale=0.9]
% CV surchargé
\draw[thick, popdark] (0,0) rectangle (4.6,6.4);
\node[font=\small\bfseries] at (2.3,6.0) {CV A : surchargé};
\foreach \y in {5.5,5.2,4.9,4.6,4.3,4.0,3.7,3.4,3.1,2.8,2.5,2.2,1.9,1.6,1.3,1.0,0.7,0.4}{
  \draw[popmuted] (0.25,\y) -- (4.35,\y);
}
\draw[thick, poporange] (0.25,3.9) rectangle (4.35,4.7);
\draw[thick, poppurple] (0.25,1.0) rectangle (4.35,2.0);
\node[font=\tiny, poporange, anchor=west] at (0.3,4.3) {cadre};
\node[font=\tiny, poppurple, anchor=west] at (0.3,1.5) {cadre};
\node[font=\footnotesize, align=center, popdark] at (2.3,-0.5) {Texte dense, cadres multiples,\\ aucun repère visuel};
% CV aéré
\begin{scope}[xshift=6.4cm]
\draw[thick, popblue] (0,0) rectangle (4.6,6.4);
\node[font=\small\bfseries, popblue] at (2.3,6.0) {CV B : aéré et hiérarchisé};
\fill[popblueL] (0.25,5.35) rectangle (4.35,5.85);
\node[font=\footnotesize\bfseries, popdark] at (2.3,5.6) {NOM Prénom};
\draw[popblue, thick] (0.25,4.95) -- (2.2,4.95);
\node[font=\tiny\bfseries, popblue, anchor=west] at (0.25,5.12) {FORMATION};
\foreach \y in {4.65,4.35,4.05}{\draw[popmuted] (0.45,\y) -- (4.35,\y); \fill[popblue] (0.32,\y) circle (0.03);}
\draw[popblue, thick] (0.25,3.35) -- (3.4,3.35);
\node[font=\tiny\bfseries, popblue, anchor=west] at (0.25,3.52) {EXPÉRIENCE PROFESSIONNELLE};
\foreach \y in {3.05,2.75,2.45}{\draw[popmuted] (0.45,\y) -- (4.35,\y); \fill[popblue] (0.32,\y) circle (0.03);}
\draw[popblue, thick] (0.25,1.75) -- (2.6,1.75);
\node[font=\tiny\bfseries, popblue, anchor=west] at (0.25,1.92) {COMPÉTENCES};
\foreach \y in {1.45,1.15}{\draw[popmuted] (0.45,\y) -- (4.35,\y); \fill[popblue] (0.32,\y) circle (0.03);}
\node[font=\footnotesize, align=center, popblue] at (2.3,-0.5) {Rubriques titrées, puces,\\ espaces de respiration};
\end{scope}
% Flèche
\draw[-{Stealth[length=3mm]}, very thick, popgreen] (4.9,3.2) -- (6.1,3.2);
\node[font=\footnotesize\bfseries, popgreen] at (5.5,3.55) {à viser};
\end{tikzpicture}
\legende{Schéma comparatif de deux maquettes de CV pour un même profil : à gauche, un document surchargé (paragraphes denses, cadres multiples) ; à droite, un document aéré (nom en bandeau, rubriques soulignées, informations en puces).}
\end{popfigure}

\courssub{La typographie : une seule police lisible}

\begin{definition}[Typographie]
La \cle{typographie} désigne l'ensemble des choix relatifs aux caractères : la police (le dessin des lettres), la taille, le gras, l'italique.
\end{definition}

Les règles à respecter sont strictes :
\begin{itemize}
\item \textbf{Une seule police}, lisible et sobre (par exemple Arial, Calibri ou Times New Roman) ;
\item une \textbf{taille de 10 à 12 points} pour le texte courant, légèrement supérieure pour le nom et les titres ;
\item le \textbf{gras est réservé aux titres} de rubriques et au nom du candidat ; l'italique, avec modération, peut signaler les intitulés de diplômes ou de postes ;
\item pas de couleurs multiples : le noir suffit, éventuellement accompagné d'\emph{une} seule couleur discrète pour les titres.
\end{itemize}

\begin{attention}[Erreur fréquente : le CV « carnaval »]
Multiplier les polices, les couleurs et les cadres est la faute la plus répandue. Trois polices, du rouge, du bleu, du vert, des bordures partout : le lecteur est fatigué avant d'avoir commencé, et le document perd tout sérieux. Un CV n'est pas une affiche publicitaire : sa beauté tient à sa \emph{sobriété} et à sa \emph{régularité}.
\end{attention}

\courssub{La photo d'identité}

Au Cameroun, l'usage courant veut que le CV comporte une \cle{photo d'identité}, placée \textbf{en haut à droite ou à gauche} du document. Elle doit être :
\begin{itemize}
\item \textbf{récente} : elle doit ressembler au candidat tel qu'il se présentera à l'entretien ;
\item \textbf{professionnelle} : visage de face, fond neutre, tenue correcte, expression souriante et sérieuse — jamais de photo de fête, de selfie ou de photo découpée sur une image de groupe.
\end{itemize}

\begin{savaistu}[La photo sur le CV : une question de pays]
Dans certains pays (Royaume-Uni, États-Unis, Canada), la photo est déconseillée, voire interdite sur le CV, afin d'éviter les discriminations liées à l'apparence physique ou à l'origine. Au Cameroun et dans la plupart des pays francophones d'Afrique, elle reste d'usage courant et même attendue par de nombreux employeurs : elle doit alors être irréprochable, car elle participe de la première impression.
\end{savaistu}

\courssub{Les rubriques signalées par des titres visibles}

La structure externe fait apparaître, par des \cle{titres} clairement visibles (en gras, éventuellement soulignés ou surlignés d'une bande de couleur discrète), les grandes rubriques du CV :

\begin{center}
\begin{tabularx}{\linewidth}{|l|X|}\hline
\rowcolor{popblueL} \textbf{Rubrique} & \textbf{Contenu attendu} \\ \hline
État civil & Nom, prénom(s), date et lieu de naissance, nationalité, situation matrimoniale, contacts (téléphone, adresse électronique, quartier/ville) \\ \hline
Formation & Diplômes obtenus ou en préparation, du plus récent au plus ancien, avec les établissements et les années \\ \hline
Expérience professionnelle & Stages, emplois, travaux réalisés, du plus récent au plus ancien \\ \hline
Compétences & Savoir-faire techniques, langues parlées et écrites, maîtrise de l'outil informatique \\ \hline
Centres d'intérêt & Activités sportives, culturelles ou associatives significatives \\ \hline
\end{tabularx}
\end{center}

Ces rubriques seront étudiées en détail dans la section suivante (structure interne) ; retenons ici qu'elles doivent être \emph{repérables en un coup d'œil} grâce à leurs titres.

\courssub{La lisibilité : puces, verbes d'action, phrases courtes}

À l'intérieur de chaque rubrique, la rédaction obéit à trois règles :
\begin{itemize}
\item utiliser des \cle{puces} plutôt que des paragraphes rédigés ;
\item commencer chaque ligne par un \cle{verbe d'action} à l'infinitif ou au participe passé : \emph{installer, réparer, gérer, accueillir, réalisé, effectué} ;
\item formuler des \textbf{phrases courtes}, sans subordonnées inutiles.
\end{itemize}

\begin{exemplebox}[Du paragraphe à la puce]
\textbf{Mauvaise rédaction (paragraphe)} : « Pendant mon stage à l'entreprise SODECOTON de Garoua, j'ai eu l'occasion de participer à la maintenance des machines et j'ai aussi aidé les techniciens dans leurs tâches quotidiennes. »

\textbf{Bonne rédaction (puces et verbes d'action)} :
\begin{itemize}
\item Participé à la maintenance des machines industrielles ;
\item Assisté les techniciens dans les opérations quotidiennes.
\end{itemize}
L'information est la même, mais elle se lit en cinq secondes au lieu de vingt.
\end{exemplebox}

\courssub{Le support : papier de qualité ou fichier PDF}

Le support matériel fait aussi partie de la structure externe :
\begin{itemize}
\item pour une remise en main propre, utiliser une feuille \textbf{blanche, de bonne qualité}, non froissée, non tachée, imprimée proprement ;
\item pour un envoi électronique (courriel, WhatsApp), enregistrer le CV au format \cle{PDF}, qui conserve la mise en page sur tous les écrans, contrairement au fichier de traitement de texte.
\end{itemize}

\begin{exemplebox}[Le CV envoyé par WhatsApp]
Un technicien en froid commercial de Douala postule à une offre d'emploi en envoyant son CV par WhatsApp. S'il envoie le fichier Word, la mise en page risque d'être déformée sur le téléphone du recruteur : polices remplacées, photo déplacée, rubriques désalignées. En envoyant un fichier \textbf{PDF}, il garantit que le recruteur voit exactement le document qu'il a conçu, quel que soit l'appareil utilisé.
\end{exemplebox}

\begin{popfigure}
\begin{center}
\begin{tikzpicture}
\begin{axis}[
  ybar, bar width=0.9cm,
  width=0.85\linewidth, height=6cm,
  symbolic x coords={Surchargé, Moyen, Aéré et hiérarchisé},
  xtick=data,
  ymin=0, ymax=90,
  ylabel={Temps moyen de lecture (secondes)},
  ylabel style={font=\small},
  x tick label style={font=\small},
  ytick={0,20,40,60,80},
  nodes near coords,
  every node near coord/.append style={font=\small\bfseries},
  axis lines=left,
  grid=major, grid style={popline!40},
]
\addplot[fill=poporange, draw=popdark] coordinates {(Surchargé,18) (Moyen,35) (Aéré et hiérarchisé,65)};
\end{axis}
\end{tikzpicture}
\end{center}
\legende{Temps moyen de lecture accordé à un CV par un recruteur selon la qualité de la mise en page (données indicatives, à titre pédagogique). Plus le CV est aéré et hiérarchisé, plus le recruteur lui consacre de temps.}
\end{popfigure}

Ce diagramme confirme notre observation initiale : un CV bien mis en page est lu \emph{plus longtemps}, donc plus attentivement. La structure externe augmente directement les chances du candidat.

\courssub{Méthode et entraînement}

\begin{methode}[Pour une mise en page efficace]
\begin{enumerate}
\item \textbf{Limiter} le document à une page : ne garder que l'essentiel.
\item \textbf{Structurer} : faire apparaître les rubriques par des titres visibles (gras, taille supérieure).
\item \textbf{Alléger} : remplacer les paragraphes par des puces commençant par un verbe d'action.
\item \textbf{Unifier} : choisir une seule police lisible, taille 10 à 12, gras réservé aux titres.
\item \textbf{Encadrer} : régler des marges régulières (environ 2 cm) et aligner le texte à gauche.
\item \textbf{Finaliser} : insérer une photo professionnelle en haut, relire, puis enregistrer en PDF avant tout envoi électronique.
\end{enumerate}
\end{methode}

\begin{exoresolu}[Diagnostiquer un CV défectueux]
\textbf{Énoncé.} Un élève de Terminale technique présente un CV de deux pages, rédigé avec quatre polices différentes, des titres en rouge, vert et bleu, un long paragraphe de huit lignes pour décrire son stage, et une photo prise lors d'une sortie entre amis. Relevez cinq défauts de structure externe et proposez pour chacun une correction.

\tcblower
\textbf{Solution détaillée.}
\begin{enumerate}
\item \textbf{Longueur excessive} (deux pages pour un jeune candidat) : réduire à une page en supprimant les informations inutiles.
\item \textbf{Multiplication des polices} (quatre polices) : n'en conserver qu'une seule, sobre et lisible.
\item \textbf{Couleurs multiples} pour les titres : revenir au noir, ou à une seule couleur discrète identique pour tous les titres.
\item \textbf{Paragraphe rédigé} pour le stage : découper en puces commençant par des verbes d'action (\emph{participé, assisté, réalisé}).
\item \textbf{Photo non professionnelle} : la remplacer par une photo d'identité récente, sur fond neutre, en tenue correcte, placée en haut du document.
\end{enumerate}
\end{exoresolu}

\begin{exercice}[À vous de jouer]
Voici la description de deux CV reçus pour un poste de secrétaire comptable à Yaoundé. Le CV 1 tient sur une page, utilise la police Calibri taille 11, présente cinq rubriques titrées en gras, des puces et une photo d'identité récente. Le CV 2 tient sur trois pages, mélange trois polices, rédige chaque expérience en paragraphes et ne comporte aucune rubrique visible.
\begin{enumerate}
\item Lequel des deux CV retiendra l'attention du recruteur ? Justifiez en citant au moins quatre règles de la structure externe.
\item Rédigez, sous forme de deux puces avec verbes d'action, l'expérience suivante : « J'ai tenu la caisse de la boutique de mon oncle pendant les vacances et j'accueillais les clients. »
\end{enumerate}
\end{exercice}

\begin{corrige}[Exercice : les deux CV]
\begin{enumerate}
\item C'est le \textbf{CV 1} qui retiendra l'attention. Il respecte les règles de la structure externe : une seule page (longueur), une police unique et lisible (typographie), des rubriques signalées par des titres en gras (hiérarchie visuelle), des puces au lieu de paragraphes (lisibilité), une photo professionnelle. Le CV 2 enfreint toutes ces règles : il sera écarté ou lu en quelques secondes.
\item \begin{itemize}
\item Tenu la caisse d'une boutique commerciale pendant les congés scolaires ;
\item Accueilli et renseigné la clientèle.
\end{itemize}
\end{enumerate}
\end{corrige}

\begin{aretenir}[L'essentiel de la section]
La \cle{structure externe} du CV, c'est son apparence : \textbf{une page}, une \textbf{mise en page aérée} (marges régulières, alignement à gauche, hiérarchie des titres), une \textbf{seule police} de taille 10 à 12 avec le gras réservé aux titres, une \textbf{photo professionnelle} en haut du document, des \textbf{rubriques visibles} (État civil, Formation, Expérience professionnelle, Compétences, Centres d'intérêt), des \textbf{puces et des verbes d'action}, et un support soigné (\textbf{papier blanc de qualité} ou \textbf{fichier PDF}). Un CV beau par sa sobriété est un CV qui sera lu.
\end{aretenir}

\coursec{La structure interne du CV : les rubriques}

Dans la section précédente, nous avons étudié la \cle{structure externe} du CV : sa présentation matérielle (marges, police, longueur, aération, mise en page). Un CV beau mais vide ne décroche aucun entretien. Nous allons maintenant ouvrir ce document et regarder ce qu'il contient : c'est la \cle{structure interne}, c'est-à-dire l'ensemble des \cle{rubriques} qui organisent les informations sur le candidat. Pour cela, nous partons de l'observation du \cle{texte iconique n° 3} : un CV complet, que nous allons identifier puis décortiquer rubrique par rubrique.

\courssub{Observation du texte iconique n° 3 : un CV complet}

\begin{experience}[Observer un CV complet (texte iconique n° 3)]
\textbf{Matériel} : le CV de Mme NGO BELL Aïcha, candidate à un poste de secrétaire comptable, distribué en classe.

\textbf{Protocole d'observation} :
\begin{enumerate}
\item Lire le document en silence, sans s'arrêter sur les détails.
\item Entourer les titres de parties (les rubriques).
\item Numéroter ces rubriques dans l'ordre où elles apparaissent.
\item Pour chaque rubrique, noter en une phrase le type d'information qu'elle contient.
\end{enumerate}

\textbf{Observations} : le document comporte sept blocs clairement séparés : une identification en haut, un état civil, une formation, une expérience professionnelle, des compétences, des langues, des centres d'intérêt. Chaque bloc porte un titre visible et les informations y sont brèves, présentées en listes.

\textbf{Interprétation} : un CV n'est pas une rédaction continue. C'est un document \emph{structuré en rubriques}, chacune répondant à une question précise que se pose l'employeur : Qui es-tu ? Qu'as-tu appris ? Qu'as-tu déjà fait ? Que sais-tu faire ?
\end{experience}

\begin{aretenir}[Les rubriques d'un CV]
Un CV complet comporte généralement les rubriques suivantes :
\begin{enumerate}
\item l'\cle{en-tête} (nom, prénom, contacts) ;
\item l'\cle{état civil} ;
\item la \cle{formation} (diplômes) ;
\item l'\cle{expérience professionnelle} ;
\item les \cle{compétences} ;
\item les \cle{langues} ;
\item les \cle{centres d'intérêt} ;
\item les \cle{références} (facultatif).
\end{enumerate}
\end{aretenir}

\begin{popfigure}
\begin{center}
\begin{tikzpicture}[
  grow=right,
  level 1/.style={sibling distance=2.6cm, level distance=3.2cm},
  level 2/.style={sibling distance=0.85cm, level distance=3.6cm},
  every node/.style={draw=popblue, rounded corners=2pt, fill=popblueL, font=\small, align=center, inner sep=2pt},
  edge from parent/.style={draw=popmuted, thick}
]
\node[fill=poporangeL, draw=poporange, font=\bfseries] {CV}
  child { node {Centres d'intérêt}
    child { node {sport, lecture} }
    child { node {associations} }
  }
  child { node {Langues}
    child { node {niveau : lu, écrit, parlé} }
    child { node {français, anglais, langues nationales} }
  }
  child { node {Compétences}
    child { node {transversales} }
    child { node {techniques} }
  }
  child { node {Expérience}
    child { node {poste, structure, période, tâches} }
  }
  child { node {Formation}
    child { node {diplômes, établissements, années} }
  }
  child { node {État civil}
    child { node {naissance, nationalité} }
  }
  child { node {En-tête}
    child { node {nom, prénom, contacts} }
  };
\end{tikzpicture}
\end{center}
\legende{Arbre de la structure interne du CV : le document se décompose en rubriques principales, elles-mêmes divisées en sous-éléments.}
\end{popfigure}

\courssub{L'en-tête : nom, prénom et contacts}

L'\cle{en-tête} est la première chose que lit l'employeur. Elle doit permettre de vous identifier et de vous joindre \emph{immédiatement}. Elle contient :

\begin{itemize}
\item le \textbf{nom} (en majuscules) et le \textbf{prénom} ;
\item le \textbf{numéro de téléphone} (un numéro sur lequel on vous joint réellement) ;
\item l'\textbf{adresse électronique professionnelle} ;
\item la \textbf{ville de résidence} (inutile de donner l'adresse complète quartier par quartier).
\end{itemize}

\begin{exemplebox}[Un en-tête correct]
\begin{center}
\textbf{NGO BELL Aïcha}\\
Yaoundé (Cameroun) — Tél. : 6 90 00 00 00\\
Courriel : aicha.ngobell@mail.com
\end{center}
\end{exemplebox}

\begin{savaistu}[L'adresse électronique, votre première image]
Une adresse électronique fantaisiste du type \emph{beaugoss237@...} ou \emph{princesse\_diamant@...} donne une mauvaise image du candidat avant même la lecture du CV. Les recruteurs y voient un manque de sérieux. Créez une adresse sobre, de type \textbf{prenom.nom@...} (par exemple : \emph{aicha.ngobell@mail.com}), réservée à vos démarches professionnelles. C'est gratuit et cela prend cinq minutes.
\end{savaistu}

\courssub{L'état civil}

L'\cle{état civil} donne les informations administratives de base :

\begin{itemize}
\item la \textbf{date et le lieu de naissance} : « Née le 12 mars 2005 à Bafoussam » ;
\item la \textbf{nationalité} : « Camerounaise » ;
\item la \textbf{situation matrimoniale} : célibataire, marié(e), etc. Cette mention est \emph{facultative} : vous pouvez l'omettre si vous le souhaitez.
\end{itemize}

\begin{attention}[Sobriété de l'état civil]
Ne transformez pas l'état civil en biographie. Le nombre de frères et sœurs, la profession des parents ou le village d'origine n'ont rien à faire dans un CV. Restez sur les informations utiles à l'employeur.
\end{attention}

\courssub{La formation : les diplômes}

La rubrique \cle{formation} répond à la question : « Qu'as-tu appris ? ». Elle liste les diplômes obtenus ou en préparation, avec l'établissement et l'année.

\begin{definition}[Ordre antichronologique]
L'\cle{ordre antichronologique} consiste à présenter les diplômes et les expériences \textbf{du plus récent au plus ancien}. On commence donc par ce que l'on fait maintenant (ou ce que l'on vient d'obtenir), et l'on remonte vers le passé. C'est l'ordre attendu dans un CV, car l'employeur s'intéresse d'abord à votre situation actuelle.
\end{definition}

\begin{exemplebox}[Rubrique « Formation » bien rédigée]
\textbf{FORMATION}
\begin{itemize}
\item 2025--2026 : Baccalauréat technique F3 (en préparation), Lycée technique de Nkongsamba.
\item 2024 : Probatoire technique F3, Lycée technique de Nkongsamba.
\item 2022 : BEPC, Collège d'enseignement général de Bangangté.
\end{itemize}
Remarquez l'ordre : on commence par le diplôme en préparation (le plus récent), puis on remonte.
\end{exemplebox}

\courssub{L'expérience professionnelle}

La rubrique \cle{expérience professionnelle} répond à la question : « Qu'as-tu déjà fait ? ». Pour un élève de Terminale technique, elle comprend les \textbf{stages}, les \textbf{emplois} éventuels et les \textbf{vacations} (petits boulots, travaux pendant les vacances). Chaque expérience se présente avec quatre éléments : le \textbf{poste}, la \textbf{structure}, la \textbf{période} et les \textbf{tâches principales}.

\begin{methode}[Rédiger la rubrique « Expérience »]
Pour chaque expérience, rédigez une entrée en quatre temps :
\begin{enumerate}
\item Indiquez le \textbf{poste occupé} (stagiaire en maintenance, vendeuse, aide-électricien...).
\item Nommez la \textbf{structure} (entreprise, atelier, administration) et sa ville.
\item Précisez la \textbf{période} (mois et années : juin--août 2025).
\item Décrivez \textbf{2 ou 3 tâches} à l'aide de \cle{verbes d'action} à l'infinitif ou au participe : \emph{gérer, installer, réparer, accueillir, rédiger, contrôler, organiser}.
\end{enumerate}
Présentez les expériences, comme les diplômes, en ordre antichronologique.
\end{methode}

\begin{exemplebox}[Exemple chiffré d'entrée « Expérience »]
\textbf{EXPÉRIENCE PROFESSIONNELLE}
\begin{itemize}
\item \textbf{Stagiaire en maintenance industrielle}, SABC, Yaoundé — \emph{juin--août 2025} (3 mois) : maintenance des lignes de production ; rédaction de rapports hebdomadaires ; nettoyage et graissage des machines.
\item \textbf{Aide-comptable} (vacation), quincaillerie « Le Progrès », Bafoussam — \emph{juillet 2024} (1 mois) : accueil des clients ; tenue du cahier de caisse.
\end{itemize}
\end{exemplebox}

\begin{attention}[Ne jamais mentir dans un CV]
Erreur fréquente et grave : \textbf{mentir} sur un diplôme ou une expérience pour « gonfler » son CV. Toute fausse information peut être vérifiée (appel à l'établissement, demande d'attestation, épreuve pratique à l'embauche). Un mensonge découvert \emph{disqualifie immédiatement} le candidat, parfois après l'embauche, et détruit sa réputation. Un CV modeste mais vrai vaut mieux qu'un CV impressionnant mais faux.
\end{attention}

\courssub{Les compétences}

La rubrique \cle{compétences} répond à la question : « Que sais-tu faire ? ». On distingue deux familles de compétences :

\begin{itemize}
\item les compétences \textbf{techniques}, liées à un savoir-faire précis : informatique (traitement de texte, tableur), conduite d'engins, soudure, câblage électrique, saisie comptable ;
\item les compétences \textbf{transversales}, utiles dans tout métier : travail en équipe, ponctualité, sens de l'organisation, respect des consignes, esprit d'initiative.
\end{itemize}

\begin{exemplebox}[Rubrique « Compétences »]
\textbf{COMPÉTENCES}
\begin{itemize}
\item Techniques : traitement de texte et tableur ; lecture de schémas électriques ; conduite d'engins (permis B).
\item Transversales : travail en équipe ; ponctualité ; sens de l'organisation.
\end{itemize}
\end{exemplebox}

\courssub{Les langues}

La rubrique \cle{langues} précise les langues que vous maîtrisez, avec un \textbf{niveau} honnête. Au Cameroun, on mentionne typiquement le \textbf{français}, l'\textbf{anglais} et éventuellement des \textbf{langues nationales}. Le niveau s'exprime par les trois savoirs : \emph{lu, écrit, parlé}.

\begin{exemplebox}[Rubrique « Langues »]
\textbf{LANGUES}
\begin{itemize}
\item Français : lu, écrit, parlé (courant).
\item Anglais : lu, écrit (niveau scolaire), parlé (notions).
\item Ewondo : parlé (langue maternelle).
\end{itemize}
\end{exemplebox}

\courssub{Les centres d'intérêt}

La rubrique \cle{centres d'intérêt} (ou « loisirs ») humanise le CV et peut faire la différence entre deux candidats. Elle doit mentionner des activités \textbf{valorisantes} et \textbf{véridiques} : sport pratiqué régulièrement, lecture, engagement associatif, musique, informatique.

\begin{attention}[Centres d'intérêt : véridiques et précis]
N'écrivez jamais une activité que vous ne pratiquez pas : le recruteur peut vous interroger dessus en entretien (« Quel livre lisez-vous en ce moment ? »). Préférez des mentions précises (« football en club scolaire, capitaine de l'équipe ») à des mentions vagues (« sport »). Évitez les passe-temps sans valeur ajoutée (« regarder la télévision »).
\end{attention}

\courssub{Les références (facultatif)}

Les \cle{références} sont des personnes (chef d'atelier, enseignant, maître de stage) qui acceptent de \textbf{témoigner du sérieux} du candidat. Cette rubrique est \emph{facultative}. Si vous l'incluez, donnez le nom, la fonction et le téléphone de chaque référent — après avoir obtenu son accord. Sinon, on écrit simplement : « Références disponibles sur demande ».

\courssub{Analyser la situation d'embauche pour adapter son CV}

Avant de choisir la mise en forme finale, il faut \textbf{analyser l'offre} ou le contexte de recrutement. Cette étape détermine quelles rubriques mettre en avant et quel type de CV privilégier.

\begin{methode}[Analyser une offre d'emploi]
\begin{enumerate}
\item Identifier le \textbf{poste visé} et les \textbf{mots-clés} de l'annonce (compétences requises, diplômes exigés, qualités demandées).
\item Repérer les \textbf{compétences techniques} et \textbf{transversales} attendues.
\item Noter les \textbf{expériences} ou \textbf{formations} spécifiques valorisées.
\item Comparer avec votre propre parcours : quels points forts mettre en évidence ? Quels manques combler par une formulation adaptée ?
\end{enumerate}
\end{methode}

\begin{exemplebox}[Adapter son CV : exemple]
Pour un poste de \emph{secrétaire comptable}, on mettra en avant : la rubrique « Compétences » (maîtrise du tableur, logiciels de comptabilité), la « Formation » (Bac F3, modules gestion), l'« Expérience » (stage en comptabilité, vacation caisse). On placera ces rubriques juste après l'en-tête. Pour un poste de \emph{technicien de maintenance}, on mettra en avant « Expérience » (stage maintenance), « Compétences techniques » (lecture plans, soudure), « Formation » (spécialité F2/F3).
\end{exemplebox}

\courssub{Les trois types de CV}

Selon la manière d'organiser les rubriques « Formation » et « Expérience », on distingue trois types de CV. Le choix dépend de votre parcours et de l'offre analysée.

\begin{definition}[Les trois types de CV]
\begin{itemize}
\item Le CV \cle{chronologique} présente diplômes et expériences en ordre antichronologique : il met en valeur le parcours.
\item Le CV \cle{fonctionnel} (ou par compétences) regroupe les informations par domaines de compétences, sans suivre les dates : il met en valeur les savoir-faire.
\item Le CV \cle{mixte} combine les deux : un résumé de compétences en tête, puis le parcours en ordre antichronologique.
\end{itemize}
\end{definition}

\begin{center}
\rowcolors{1}{popblueL}{white}
\begin{tabularx}{\linewidth}{|l|X|X|X|}
\hline
\rowcolor{popblue} \textcolor{white}{\textbf{Critère}} & \textcolor{white}{\textbf{Chronologique}} & \textcolor{white}{\textbf{Fonctionnel}} & \textcolor{white}{\textbf{Mixte}} \\
\hline
Principe & Parcours du plus récent au plus ancien & Regroupement par compétences & Compétences en tête, puis parcours daté \\
Avantage & Clair, facile à lire, apprécié des recruteurs & Masque un parcours court ou interrompu & Valorise à la fois compétences et parcours \\
Limite & Révèle les trous dans le parcours & Peut sembler cacher quelque chose & Plus long à construire \\
Pour qui ? & Jeune diplômé au parcours régulier & Candidat en reconversion ou sans expérience & Candidat expérimenté ou profil polyvalent \\
\hline
\end{tabularx}
\end{center}

Pour un élève de Terminale technique, le \textbf{CV chronologique} est le plus adapté : le parcours scolaire est régulier et les stages récents sont un atout.

\courssub{Relire et améliorer son CV}

La rédaction n'est pas terminée tant que le CV n'a pas été relu et corrigé. Une coquille ou une maladresse peut coûter un entretien.

\begin{methode}[Check-list de relecture avant envoi]
\begin{enumerate}
\item \textbf{Orthographe et grammaire} : relire à l'envers (mot par mot) ; utiliser un correcteur ; faire relire par un tiers.
\item \textbf{Cohérence des dates} : pas de trous inexpliqués, pas de chevauchements impossibles.
\item \textbf{Contact} : numéro de téléphone valide, adresse mail professionnelle, nom de fichier clair (\emph{Nom\_Prenom\_CV.pdf}).
\item \textbf{Mise en page} : alignements parfaits, police unique, espacements réguliers, une page maximum.
\item \textbf{Pertinence} : chaque information répond-elle à l'offre visée ? Supprimer le superflu.
\item \textbf{Honnêteté} : tout est vérifiable, rien n'est exagéré.
\end{enumerate}
\end{methode}

\begin{exoresolu}[Identifier et compléter des rubriques]
\textbf{Énoncé.} Voici des informations extraites du CV de TCHOUPO Marc : « réparation d'installations électriques ; juillet--septembre 2025 ; entreprise ENEO, Douala ; aide-électricien ». 1) À quelle rubrique appartiennent ces informations ? 2) Rédigez l'entrée complète selon la méthode en quatre éléments. 3) Marc hésite à écrire qu'il a dirigé une équipe de dix techniciens pendant ce stage. Que lui conseillez-vous ?
\tcblower
\textbf{Solution détaillée.}
\begin{enumerate}
\item Ces informations (poste, structure, période, tâches) appartiennent à la rubrique \textbf{Expérience professionnelle}.
\item Entrée rédigée : « \textbf{Aide-électricien} (stage), entreprise ENEO, Douala — \emph{juillet--septembre 2025} (3 mois) : réparation d'installations électriques ; assistance aux techniciens ; respect des consignes de sécurité. » On retrouve les quatre éléments : poste, structure, période, tâches exprimées par des verbes d'action.
\item Marc ne doit \textbf{pas} écrire cela : c'est un mensonge. Un stagiaire aide-électricien ne dirige pas dix techniciens ; l'employeur peut vérifier auprès d'ENEO ou poser une question technique en entretien. La fausse information le disqualifierait. Il doit valoriser ce qu'il a réellement fait.
\end{enumerate}
\end{exoresolu}

\begin{exercice}[Construire les rubriques de votre CV]
À partir de votre propre parcours, rédigez sur une feuille : 1) votre en-tête complet (avec une adresse électronique sobre) ; 2) votre état civil ; 3) votre rubrique « Formation » en ordre antichronologique ; 4) une entrée « Expérience » décrivant votre dernier stage, selon la méthode en quatre éléments ; 5) vos rubriques « Compétences », « Langues » et « Centres d'intérêt ».
\end{exercice}

\begin{corrige}[Exercice : construire les rubriques de votre CV]
Corrigé type (à adapter à chaque élève) :
\begin{enumerate}
\item \textbf{En-tête} : « MBALLA Jean — Douala — Tél. : 6 99 00 00 00 — jean.mballa@mail.com ». Vérifier : nom en majuscules, téléphone exact, adresse sobre de type prenom.nom@, ville indiquée.
\item \textbf{État civil} : « Né le 5 mai 2006 à Douala. Nationalité camerounaise. » (situation matrimoniale facultative).
\item \textbf{Formation} (ordre antichronologique) : « 2025--2026 : Bac technique F4 (en préparation), Lycée technique de Douala-Bassa ; 2024 : Probatoire F4, même établissement ; 2022 : BEPC, CES de Bonabéri. »
\item \textbf{Expérience} : « Stagiaire en mécanique auto, garage Moderne, Douala — juillet--août 2025 (2 mois) : vidange et entretien des véhicules ; réparation de freins ; accueil des clients. » Poste + structure + période + verbes d'action : tous les éléments sont présents.
\item \textbf{Compétences} : « Techniques : soudure à l'arc, tableur. Transversales : ponctualité, travail en équipe. » \textbf{Langues} : « Français : lu, écrit, parlé ; anglais : notions ; douala : parlé. » \textbf{Centres d'intérêt} : « football en club, lecture de romans africains, membre du club informatique du lycée ».
\end{enumerate}
Barème indicatif : 1 point par rubrique correctement rédigée, 1 point pour l'ordre antichronologique, 1 point pour les verbes d'action, 1 point pour l'honnêteté et la sobriété des informations.
\end{corrige}

\begin{aretenir}[L'essentiel de la section]
\begin{itemize}
\item Le CV est organisé en \cle{rubriques} : en-tête, état civil, formation, expérience, compétences, langues, centres d'intérêt, références (facultatif).
\item Formation et expérience se présentent en \cle{ordre antichronologique} : du plus récent au plus ancien.
\item Chaque expérience = poste + structure + période + 2 ou 3 tâches en \cle{verbes d'action}.
\item On distingue trois types de CV : \cle{chronologique}, \cle{fonctionnel}, \cle{mixte}. Le choix dépend du parcours et de l'offre.
\item Règle d'or : un CV dit toujours la \textbf{vérité} ; toute fausse information disqualifie le candidat.
\item Avant d'envoyer : relire, corriger, vérifier la cohérence, nommer le fichier proprement.
\end{itemize}
\end{aretenir}

\coursec{Analyser la situation d'embauche}

Dans la section précédente, nous avons appris à construire les \cle{rubriques} du CV : état civil, formation, expérience professionnelle, compétences, centres d'intérêt. Mais un CV bien structuré ne suffit pas. Imaginez deux candidats qui postulent au même emploi : le premier envoie le même CV à toutes les entreprises depuis deux ans ; le second lit attentivement l'offre et adapte son CV avant de l'envoyer. À compétences égales, c'est le second qui sera convoqué à l'entretien. Pourquoi ? Parce qu'un CV n'a de valeur que par rapport à une \cle{situation d'embauche} précise. Cette section vous apprend donc à \textbf{lire une offre d'emploi comme un recruteur}, à en dégager les attentes, puis à faire correspondre votre parcours à ces attentes.

\courssub{Observation du texte iconique n° 4 : une offre d'emploi réelle}

\begin{experience}[Observer pour comprendre : l'offre d'emploi comme texte iconique]
\textbf{Matériel :} une offre d'emploi réelle, découpée dans un journal (par exemple \emph{Cameroon Tribune}), relevée sur un site d'annonces ou photographiée sur le panneau d'affichage d'une entreprise camerounaise.

\textbf{Protocole d'observation :}
\begin{enumerate}
\item Observer la \textbf{mise en page} : logo de l'entreprise, titre du poste en gros caractères, encadrés, date limite soulignée.
\item Repérer les \textbf{blocs d'information} : qui recrute ? quel poste ? quel profil ? où ? comment postuler ?
\item Noter le \textbf{vocabulaire} employé : verbes d'action (« assurer », « gérer », « maîtriser »), noms de diplômes, durées d'expérience.
\end{enumerate}

\textbf{Observations :} l'offre est un texte \emph{bref, structuré et fonctionnel}. Chaque ligne coûte cher à l'entreprise (publicité payante) : il n'y a donc \textbf{aucun mot inutile}. Tout ce qui est écrit est une exigence ou une information pratique.

\textbf{Interprétation :} si l'offre ne contient que des informations utiles, alors lire une offre, c'est recevoir la \emph{liste exacte} de ce que le recruteur cherchera dans votre CV. L'offre est la question ; le CV est la réponse. On ne peut pas répondre juste sans avoir bien lu la question.
\end{experience}

\begin{definition}[La situation d'embauche]
La \cle{situation d'embauche} est l'ensemble des éléments concrets qui entourent un recrutement : l'\textbf{offre d'emploi} (le poste, le profil exigé, les diplômes, l'expérience, les compétences, le lieu, le délai), l'\textbf{entreprise} qui recrute (son secteur d'activité, sa taille, sa localisation) et le \textbf{destinataire} du CV (le recruteur, qui lira des dizaines de candidatures en peu de temps). Analyser la situation d'embauche, c'est identifier précisément ces éléments avant de rédiger.
\end{definition}

\courssub{Lecture méthodique d'une offre d'emploi}

Une offre d'emploi se lit en repérant \textbf{sept informations obligatoires}. Prenons l'habitude de les chercher dans cet ordre :

\begin{enumerate}
\item \textbf{Le poste} : l'intitulé exact du métier (« technicien de maintenance », « assistant comptable »). C'est lui qui doit apparaître dans le titre de votre CV.
\item \textbf{Le profil exigé} : la description du candidat idéal (qualités personnelles : rigueur, sens de l'organisation, autonomie).
\item \textbf{Les diplômes} : le niveau et la spécialité demandés (BT, BTS, DUT, Licence...). Vérifiez que votre diplôme correspond ou est supérieur.
\item \textbf{L'expérience} : la durée et le domaine (« au moins 2 ans dans un poste similaire »). Un stage peut parfois être valorisé comme expérience.
\item \textbf{Les compétences} : les savoir-faire techniques (« maîtrise de l'outil informatique », « connaissance des réseaux électriques »).
\item \textbf{Le lieu} : la ville ou la région d'affectation. Mentionnez votre mobilité si nécessaire.
\item \textbf{Le délai et les modalités} : la date limite de dépôt, les pièces à fournir, l'adresse ou le courriel de destination. Un CV envoyé en retard est un CV éliminé.
\end{enumerate}

\begin{methode}[Analyser une offre d'emploi en 4 étapes]
\begin{enumerate}
\item \textbf{Souligner les exigences} : à la première lecture, surlignez ou recopiez tout ce que l'offre demande (diplôme, expérience, compétences, qualités). Distinguez les exigences \emph{obligatoires} (« exigé », « indispensable ») des simples souhaits (« serait un atout », « apprécié »).
\item \textbf{Repérer les mots-clés} : isolez les noms et verbes précis de l'offre : « maintenance préventive », « saisie comptable », « Excel », « sens du relationnel ». Ces mots sont ceux que le recruteur cherchera des yeux dans les CV.
\item \textbf{Faire correspondre son parcours} : pour chaque exigence, cherchez dans votre formation, vos stages et vos activités un élément qui y répond. Construisez un tableau de correspondance (voir plus bas).
\item \textbf{Réorganiser le CV} : placez en haut et en évidence les rubriques qui répondent le mieux à l'offre ; reprenez les mots-clés de l'annonce dans vos rubriques ; supprimez ou réduisez ce qui est hors sujet.
\end{enumerate}
\end{methode}

\begin{popfigure}
\begin{center}
\begin{tikzpicture}[
  node distance=0.55cm and 0.9cm,
  box/.style={draw=popblue, fill=popblueL, rounded corners=2pt, text width=3.1cm, align=center, font=\small, minimum height=1.1cm},
  arr/.style={-{Stealth[length=3mm]}, thick, poporange}
]
\node[box, align=center] (offre){\textbf{1. Offre d'emploi}\\ poste, diplômes, expérience, lieu, délai};
\node[box, right=of offre, align=center] (profil){\textbf{2. Profil recherché}\\ exigences obligatoires et souhaitées};
\node[box, right=of profil, align=center] (mots){\textbf{3. Mots-clés}\\ noms et verbes précis de l'annonce};
\node[box, below=of mots, align=center] (select){\textbf{4. Sélection des informations}\\ ce qui, dans mon parcours, répond à l'offre};
\node[box, left=of select, align=center] (cv){\textbf{5. CV adapté}\\ rubriques réorganisées, mots-clés repris};
\draw[arr] (offre) -- (profil);
\draw[arr] (profil) -- (mots);
\draw[arr] (mots) -- (select);
\draw[arr] (select) -- (cv);
\draw[arr, dashed, popmuted] (cv.north) to[bend left=15] node[above, font=\scriptsize, popmuted]{vérification finale} (offre.south west);
\end{tikzpicture}
\end{center}
\legende{Démarche d'analyse d'une situation d'embauche : on part toujours de l'offre, jamais de soi. Le CV final est vérifié en le confrontant une dernière fois à l'offre (flèche en pointillés).}
\end{popfigure}

\courssub{Dégager les mots-clés et faire correspondre son parcours}

Le cœur de la méthode est le \textbf{tableau de correspondance} : à gauche, ce que l'offre exige ; à droite, ce que vous possédez. Ce travail se fait \emph{sur un brouillon}, avant toute rédaction.

\begin{center}
\begin{tabularx}{\linewidth}{|l|X|X|}
\hline
\rowcolor{popblueL} \textbf{N°} & \textbf{Exigence de l'offre} & \textbf{Élément du candidat correspondant} \\
\hline
1 & BT ou BTS en électrotechnique & BT Électrotechnique, Lycée technique de Douala (2024) \\
\hline
2 & 2 ans d'expérience en maintenance & Stage de 6 mois à la SODECOTON + 18 mois comme aide-technicien \\
\hline
3 & Maîtrise de l'outil informatique & Word et Excel (bon niveau), saisie de données au stage \\
\hline
4 & Sens de l'organisation et rigueur & Responsable du stock du club technique ; tenue d'un registre \\
\hline
5 & Disponibilité immédiate, poste à Douala & Réside à Douala, disponible de suite \\
\hline
\end{tabularx}
\end{center}

\begin{aretenir}[La règle d'or des mots-clés]
Tout mot-clé de l'offre auquel votre parcours répond doit réapparaître, \textbf{avec les mêmes termes}, dans votre CV. Si l'offre dit « saisie comptable », n'écrivez pas « tenue des livres » : écrivez « saisie comptable ». Le recruteur lit vite ; il reconnaît ses propres mots.
\end{aretenir}

\begin{propriete}[Le CV efficace est toujours ciblé]
Un même candidat ne présente \textbf{jamais} le même CV pour deux offres différentes. Le CV n'est pas une biographie figée : c'est une \emph{réponse} à une offre précise. On adapte à chaque candidature : le titre du CV, l'ordre des rubriques, le vocabulaire des compétences et la sélection des expériences détaillées.
\end{propriete}

\begin{attention}[Erreur fréquente : le CV générique]
Envoyer un CV identique à toutes les entreprises est l'erreur la plus répandue — et la plus visible. Le recruteur repère immédiatement un CV non adapté : titre absent ou vague (« recherche d'emploi »), compétences sans rapport avec le poste, aucun mot-clé de l'annonce. Un tel CV part à la corbeille en moins d'une minute. Autre piège : \textbf{mentir} pour « coller » à l'offre (fausse expérience, niveau exagéré). On adapte la présentation de son parcours réel, on ne l'invente pas : tout mensonge est démasqué à l'entretien ou à l'essai.
\end{attention}

\courssub{Choisir le type de CV adapté à l'offre et à son profil}

L'analyse de l'offre détermine aussi le \textbf{type de CV} à choisir :

\begin{itemize}
\item Le \cle{CV chronologique} (rubriques « Formation » et « Expérience » présentées de la plus récente à la plus ancienne) convient quand le parcours est \textbf{continu} et riche : il montre une progression logique. C'est le choix à faire si l'offre exige « une expérience confirmée » et que vous l'avez.
\item Le \cle{CV fonctionnel} (ou CV par compétences) regroupe les informations par \textbf{domaines de compétences} (« Électricité industrielle », « Comptabilité », « Informatique ») plutôt que par dates. Il convient quand on \textbf{manque d'expérience} (jeune diplômé), quand on a des \textbf{trous} dans le parcours, ou quand on change de métier : il met en avant ce que l'on sait faire plutôt que ce que l'on a peu fait.
\end{itemize}

\begin{exemplebox}[Adapter une rubrique aux mots-clés]
L'offre exige : « maîtrise de l'outil informatique ».

\emph{Mauvaise rédaction (générique) :} « Informatique : oui. »

\emph{Bonne rédaction (ciblée et honnête) :} « \textbf{Compétences informatiques :} traitement de texte Word (bon niveau), tableur Excel (formules simples, tableaux de suivi), saisie de données (pratiquée pendant 6 mois en stage). »

La seconde version reprend le mot-clé de l'offre, précise les logiciels et donne un \textbf{niveau honnête} et vérifiable.
\end{exemplebox}

\courssub{Cas pratiques camerounais}

\textbf{Cas 1 — Technicien de maintenance à ENEO (Douala).} L'offre exige : BT/BTS en électrotechnique, 2 ans d'expérience, connaissance des réseaux électriques, disponibilité. Le candidat choisit un CV \emph{chronologique} (parcours continu), intitule son CV « Technicien de maintenance électrique », place « Expérience professionnelle » avant « Formation », et reprend les mots-clés « maintenance préventive », « réseaux électriques », « sécurité ».

\textbf{Cas 2 — Assistant comptable dans une PME de Bafoussam.} L'offre demande : BTS en comptabilité, maîtrise d'Excel, rigueur, un débutant accepté. Le jeune diplômé, sans expérience salariée, choisit un CV \emph{fonctionnel} : rubrique « Compétences » en tête (saisie comptable, Excel, classement des pièces), puis « Formation », puis les stages valorisés comme des expériences (« Stage de 3 mois au cabinet FIDUCOM : saisie des factures, rapprochement bancaire »).

\textbf{Cas 3 — Stagiaire dans une ONG de Maroua.} L'offre insiste sur : sens du relationnel, travail en équipe, maîtrise du français, connaissance du milieu rural. Le candidat met en avant ses activités associatives (membre d'un club, participation à une campagne de sensibilisation), sa pratique du français et du fulfuldé, et sa disponibilité géographique — autant d'éléments qui répondent directement à l'offre, même sans expérience professionnelle.

\begin{exoresolu}[Analyser une offre et préparer le CV]
\textbf{Énoncé.} Voici une offre publiée dans un journal : « \emph{PME agroalimentaire de Bafoussam recrute un(e) assistant(e) comptable. Profil : BTS en comptabilité ou gestion ; maîtrise d'Excel exigée ; sens de l'organisation ; débutant accepté. Poste basé à Bafoussam. Dossier (CV + lettre) à déposer avant le 30 novembre.} »
1) Relevez les sept informations de la situation d'embauche. 2) Dégagez trois mots-clés. 3) Indiquez le type de CV adapté pour un jeune diplômé sans expérience et justifiez. 4) Rédigez la rubrique « Compétences » du CV.
\tcblower
\textbf{Solution détaillée.}
1) \textbf{Les sept informations :} poste = assistant comptable ; profil = sens de l'organisation ; diplôme = BTS comptabilité ou gestion ; expérience = débutant accepté (aucune exigée) ; compétences = maîtrise d'Excel ; lieu = Bafoussam ; délai et modalités = CV + lettre à déposer avant le 30 novembre.

2) \textbf{Trois mots-clés :} « assistant comptable », « Excel », « sens de l'organisation ».

3) \textbf{Type de CV :} le CV \emph{fonctionnel} (par compétences). Justification : le candidat est débutant, il manque d'expérience professionnelle ; le CV chronologique laisserait une rubrique « Expérience » presque vide. Le CV fonctionnel met en tête ce qu'il sait faire (comptabilité, Excel), ce qui répond directement à l'offre.

4) \textbf{Rubrique « Compétences » :}
\begin{itemize}
\item Comptabilité : saisie des écritures, classement des pièces comptables (formation BTS + stage) ;
\item Informatique : Excel (tableaux de suivi, formules simples — bon niveau), Word ;
\item Qualités : organisation et rigueur (tenue du registre de caisse du club de l'école pendant un an).
\end{itemize}
Chaque ligne reprend un mot-clé de l'offre et s'appuie sur un élément réel du parcours.
\end{exoresolu}

\begin{exercice}[À vous de jouer]
Relevez dans un journal ou sur un panneau d'affichage une offre d'emploi correspondant à votre filière (par exemple : technicien de maintenance à ENEO, secrétaire dans une entreprise de Yaoundé, ouvrier qualifié dans un atelier de Bafoussam).
1) Recopiez l'offre et soulignez les exigences obligatoires.
2) Construisez le tableau de correspondance « exigences de l'offre / mes éléments ».
3) Choisissez et justifiez le type de CV (chronologique ou fonctionnel).
4) Rédigez le titre du CV et la rubrique la plus importante pour cette offre.
\end{exercice}

\begin{corrige}[Exercice]
Corrigé-type (à adapter selon l'offre choisie).
1) Les exigences obligatoires sont signalées par « exigé », « indispensable », « au moins » ; les souhaits par « serait un atout ».
2) Le tableau doit comporter au moins quatre lignes : diplôme, expérience, compétences techniques, qualités personnelles — chacune mise en regard d'un élément réel du parcours de l'élève (formation, stages, activités).
3) CV chronologique si le parcours est continu et l'expérience réelle ; CV fonctionnel si l'élève est débutant ou a un parcours interrompu. La justification doit citer l'offre (« l'offre exige 2 ans d'expérience, or... »).
4) Le titre reprend l'intitulé exact du poste ; la rubrique rédigée reprend au moins deux mots-clés de l'offre avec des niveaux honnêtes.
\end{corrige}

\begin{aretenir}[L'essentiel de la section]
\begin{itemize}
\item L'offre d'emploi est la \textbf{question}, le CV est la \textbf{réponse} : on analyse toujours l'offre avant de rédiger.
\item Sept informations à relever : poste, profil, diplômes, expérience, compétences, lieu, délai.
\item Méthode en 4 étapes : souligner les exigences, repérer les mots-clés, faire correspondre son parcours, réorganiser le CV.
\item Le CV efficace est \textbf{ciblé} : jamais deux CV identiques pour deux offres différentes.
\item CV chronologique si le parcours est continu ; CV fonctionnel si l'on manque d'expérience.
\end{itemize}
\end{aretenir}

\coursec{Rédaction et amélioration du CV}

Dans la section précédente, nous avons appris à \emph{analyser la situation d'embauche} : lire attentivement l'offre, repérer le profil recherché (diplôme, compétences, expérience) et dégager les attentes de l'employeur. Cette analyse préalable est la clé de voûte de toute la démarche. Mais une fois que l'on sait \emph{ce que l'employeur attend}, reste l'essentiel : \textbf{produire effectivement le CV}, puis l'\textbf{améliorer} jusqu'à ce qu'il soit irréprochable. C'est l'objet de cette dernière section, qui correspond directement à l'épreuve de production de texte du Baccalauréat technique : on vous donne une situation d'embauche, et vous devez rédiger le CV du candidat.

\courssub{Vue d'ensemble : les six étapes de production}

Rédiger un CV ne s'improvise pas. Comme toute production d'écrit au programme de Terminale technique, il obéit à une \cle{démarche en étapes} : on ne passe à l'étape suivante que lorsque la précédente est achevée. Voici l'organigramme complet de la production d'un CV.

\begin{popfigure}
\begin{center}
\begin{tikzpicture}[node distance=0.55cm,
etape/.style={rectangle, rounded corners=3pt, draw=popblue, fill=popblueL, text width=10.5cm, align=left, inner sep=5pt, font=\small},
fleche/.style={-{Stealth[length=2.5mm]}, thick, poporange}]
\node[etape] (e1) {\textbf{Étape 1 — Collecter} : rassembler relevés de notes, diplômes, attestations de stage, pièces d'identité.};
\node[etape, below=of e1] (e2) {\textbf{Étape 2 — Choisir le modèle} : ordre des rubriques adapté à l'offre (chronologique ou par compétences).};
\node[etape, below=of e2] (e3) {\textbf{Étape 3 — Rédiger} : formulations brèves, précises et véridiques dans chaque rubrique.};
\node[etape, below=of e3] (e4) {\textbf{Étape 4 — Soigner la langue} : orthographe, grammaire, conjugaison, vocabulaire professionnel.};
\node[etape, below=of e4] (e5) {\textbf{Étape 5 — Mettre en page} : une page, puces, titres de rubriques visibles, photo éventuelle.};
\node[etape, below=of e5] (e6) {\textbf{Étape 6 — Relire et corriger} : relecture personnelle, relecture par un tiers, vérification des dates et contacts.};
\draw[fleche] (e1) -- (e2);
\draw[fleche] (e2) -- (e3);
\draw[fleche] (e3) -- (e4);
\draw[fleche] (e4) -- (e5);
\draw[fleche] (e5) -- (e6);
\draw[fleche, dashed, popgreen] (e6.west) to[bend right=55] node[left, font=\small\itshape, text=popgreen]{amélioration} (e3.west);
\end{tikzpicture}
\end{center}
\legende{Organigramme des six étapes de production du CV. La flèche en pointillés symbolise la boucle d'amélioration : après la relecture, on retourne à la rédaction pour corriger. La boucle se répète jusqu'à obtenir zéro faute.}
\end{popfigure}

\begin{methode}[La règle des 3 C d'un bon CV]
Un bon CV respecte la règle des \cle{3 C} :
\begin{enumerate}
\item \textbf{Clair} : la mise en page est aérée, les rubriques sont repérables en un coup d'œil, l'information essentielle se lit en moins d'une minute.
\item \textbf{Concis} : une seule page, des phrases courtes ou des simples groupes nominaux, aucun détail inutile (pas de paragraphes rédigés comme dans une lettre).
\item \textbf{Convaincant} : les rubriques sont ordonnées pour mettre en avant ce qui correspond à l'offre ; chaque information apporte une preuve (diplôme, stage, résultat).
\end{enumerate}
Avant de rendre votre CV, posez-vous trois questions : \emph{Est-ce clair ? Est-ce concis ? Est-ce convaincant ?} Si l'une des réponses est non, reprenez l'étape correspondante.
\end{methode}

\courssub{Étape 1 : collecter ses informations personnelles}

On ne rédige pas un CV de mémoire : on le rédige à partir de \textbf{documents authentiques}. L'intuition est simple : un CV contient des dates, des intitulés de diplômes, des noms d'entreprises ; la moindre erreur de date ou d'orthographe dans le nom d'un établissement décrédibilise le candidat. La collecte préalable est donc une étape de vérification documentaire.

\begin{methode}[Constituer son dossier de collecte]
\begin{enumerate}
\item Rassembler les \textbf{pièces d'identité} (pour l'état civil exact : nom, prénom, date et lieu de naissance).
\item Rassembler les \textbf{diplômes et relevés de notes} (BEPC, Probatoire, Baccalauréat, CAP, BTS\ldots) pour noter les intitulés exacts et les années d'obtention.
\item Rassembler les \textbf{attestations de stage} et certificats de travail : nom exact de l'entreprise, durée, tâches effectuées.
\item Lister les \textbf{compétences vérifiables} : langues parlées, logiciels maîtrisés, permis de conduire.
\item Noter les \textbf{coordonnées fiables} : numéro de téléphone actif, adresse électronique sérieuse (prénom.nom, pas de surnom).
\end{enumerate}
\end{methode}

\begin{attention}[Véridicité obligatoire]
Toutes les informations d'un CV doivent être \cle{véridiques}. Inventer un stage ou gonfler un diplôme est une faute grave : l'employeur peut demander les attestations, et au Baccalauréat, le correcteur sanctionne les incohérences (par exemple un stage de 6 mois pendant l'année du Baccalauréat). On peut \emph{mettre en valeur} une expérience réelle, jamais en créer une fausse.
\end{attention}

\courssub{Étape 2 : choisir le modèle et l'ordre des rubriques}

Il n'existe pas un seul CV correct, mais un CV \emph{adapté à l'offre}. Deux modèles principaux s'offrent au candidat :

\begin{center}
\begin{tabularx}{\linewidth}{|l|X|X|}\hline
\rowcolor{popblueL} \textbf{Modèle} & \textbf{Principe} & \textbf{Quand l'utiliser} \\ \hline
CV \textbf{chronologique} (inversé) & Les expériences et formations sont listées de la plus récente à la plus ancienne. & Candidat ayant déjà des stages ou emplois en rapport avec l'offre. \\ \hline
CV \textbf{par compétences} (thématique) & On regroupe les savoir-faire par domaines (technique, informatique, langues). & Jeune diplômé sans expérience, ou candidat dont les compétences comptent plus que le parcours. \\ \hline
\end{tabularx}
\end{center}

L'ordre des rubriques dépend de l'analyse de la situation d'embauche faite dans la section précédente : si l'offre exige d'abord un diplôme précis, la rubrique \emph{Formation} passe avant \emph{Expérience professionnelle} ; si l'offre exige de l'expérience, c'est l'inverse. C'est ce choix raisonné — et sa justification — qui est valorisé à l'examen.

\courssub{Étape 3 : rédiger chaque rubrique}

Chaque rubrique se rédige selon trois exigences : \textbf{brièveté} (groupes nominaux ou phrases courtes), \textbf{précision} (dates, lieux, intitulés exacts) et \textbf{véridicité}. Quelques formulations modèles :

\begin{itemize}
\item \emph{Formation} : « 2024 : Baccalauréat technique, série F3 (Génie mécanique), Lycée technique de Douala-Bassa. »
\item \emph{Expérience} : « Juillet--septembre 2023 : stage d'ouvrier mécanicien, Ets Nkoulou Frères, Yaoundé — entretien et révision de moteurs diesel. »
\item \emph{Compétences} : « Soudure à l'arc ; lecture de plans ; logiciel AutoCAD (niveau intermédiaire). »
\item \emph{Langues} : « Français : courant ; Anglais : niveau scolaire ; Ewondo : langue maternelle. »
\end{itemize}

Notez la structure constante d'une ligne d'expérience : \textbf{date + fonction + entreprise + lieu + tâche principale}. Cette régularité donne au CV son aspect professionnel.

\courssub{Étape 4 : soigner la langue}

La langue du CV est une langue \cle{professionnelle} : vocabulaire précis du métier, registre soutenu, aucune abréviation familière. Trois points de vigilance :

\begin{enumerate}
\item \textbf{Orthographe} : accords (« attestations obtenues »), majuscules aux noms propres, accents (« Électricien », « Études »).
\item \textbf{Conjugaison} : pour les tâches passées, utiliser l'infinitif (« réparation des installations ») ou le participe passé, mais jamais mélanger les temps dans une même rubrique.
\item \textbf{Vocabulaire} : préférer « assurer la maintenance » à « s'occuper de », « collaborer à » à « aider à ».
\end{enumerate}

\begin{attention}[Erreur fréquente et éliminatoire]
Les fautes d'orthographe dans un CV sont \cle{éliminatoires} : un recruteur qui repère une faute dans un document d'une page conclut que le candidat sera négligent dans son travail. Au Baccalauréat, chaque faute fait perdre des points de correction de la langue. Règle d'or : \textbf{toujours faire relire son CV par une tierce personne} (enseignant, parent, camarade sérieux) — l'œil de l'auteur ne voit plus ses propres erreurs.
\end{attention}

\courssub{Étape 5 : mettre en page}

La mise en page obéit à des normes strictes :

\begin{itemize}
\item \textbf{Une seule page} : le CV tient sur le recto d'une feuille A4.
\item \textbf{Titres de rubriques visibles} : en gras, éventuellement soulignés ou colorés avec sobriété.
\item \textbf{Puces ou tirets} pour lister les éléments ; alignement rigoureux des dates.
\item \textbf{Photo} : facultative ; si elle est jointe, elle doit être d'identité, récente et sérieuse (jamais de photo de loisir).
\item \textbf{Police unique} et lisible ; pas de fantaisies typographiques.
\end{itemize}

\begin{exemplebox}[Le standard attendu au Baccalauréat technique, en chiffres]
Un CV réussi à l'examen respecte des ordres de grandeur précis : \textbf{1 page}, \textbf{5 rubriques} (état civil et coordonnées, formation, expérience professionnelle, compétences, centres d'intérêt), \textbf{moins de 20 lignes de texte par rubrique}, et \textbf{zéro faute} d'orthographe. Un CV de deux pages ou comportant des paragraphes rédigés est hors norme.
\end{exemplebox}

\courssub{Étape 6 : relire, faire relire, corriger}

La relecture est une étape à part entière, pas une formalité. Elle se fait en trois passages :

\begin{enumerate}
\item \textbf{Relecture de fond} : les rubriques répondent-elles à l'offre ? L'ordre est-il stratégique ?
\item \textbf{Relecture de forme} : orthographe, accords, homogénéité des dates et des puces.
\item \textbf{Vérification factuelle} : les dates se succèdent-elles sans contradiction ? Le numéro de téléphone et l'adresse électronique sont-ils exacts et actifs ?
\end{enumerate}

Ensuite, on \textbf{fait relire par une tierce personne}, puis on produit une \textbf{version corrigée}.

\courssub{L'amélioration : de la première version à la version corrigée}

\begin{definition}[Remédiation]
La \cle{remédiation} consiste à reprendre les erreurs identifiées dans le compte rendu (fait par l'enseignant ou par un pair) pour produire une \textbf{version améliorée} du CV. Elle ne se limite pas à corriger les fautes : elle peut réorganiser les rubriques, reformuler les expériences, supprimer les informations hors sujet. Chaque modification doit pouvoir être \emph{justifiée}.
\end{definition}

\begin{exoresolu}[Extrait de CV : avant / après correction]
\textbf{Énoncé.} Voici l'extrait de la rubrique « Expérience professionnelle » du CV d'Abena, candidate à un poste de secrétaire-comptable. Relevez les défauts et proposez une version améliorée en justifiant chaque modification.

\emph{Version 1 :} « J'ai fais un stage a la SODECOTON a garoua ou je m'occupais de plein de choses comme le classement et aussi j'aidais la comptable, c'était en 2022 je crois, pendant les grandes vacances. »

\tcblower
\textbf{Solution détaillée.}

\textbf{Défauts relevés :} (1) faute de conjugaison « j'ai fais » au lieu de « j'ai fait » ; (2) absence d'accents : « a » pour « à », « garoua » sans majuscule ; (3) formulation vague et familière : « plein de choses », « je crois » ; (4) phrase longue rédigée comme un récit, contraire à la concision du CV ; (5) date imprécise.

\textbf{Version 2 (améliorée) :}

« Juillet--août 2022 : stage de secrétaire-comptable, SODECOTON, Garoua — classement des dossiers clients ; saisie des factures ; assistance du service comptable. »

\textbf{Justifications :} la date est précise et vérifiée sur l'attestation de stage (véracité) ; la structure « date + fonction + entreprise + lieu » est conforme (clarté) ; les tâches sont exprimées par des groupes nominaux à l'infinitif sous-entendu, avec un vocabulaire professionnel (« saisie des factures » remplace « j'aidais la comptable ») ; l'orthographe est corrigée. Le passage de 2 lignes et demie à 1 ligne respecte la concision.
\end{exoresolu}

\courssub{Complément utile au Bac : le CV électronique et la lettre de motivation}

Aujourd'hui, le CV circule souvent sous forme \cle{électronique} : fichier PDF envoyé par courrier électronique ou déposé sur des \textbf{plateformes d'emploi en ligne} (sites d'annonces, réseaux professionnels). Trois précautions s'imposent : nommer le fichier clairement (\emph{CV\_Abena\_Marie.pdf} et non \emph{document1.pdf}), vérifier que la mise en page ne se déforme pas à l'ouverture, et soigner le message d'accompagnement du courriel.

Le CV est presque toujours accompagné d'une \cle{lettre de motivation} : le CV présente les faits (diplômes, expériences), la lettre explique \emph{pourquoi} l'on postule et \emph{ce que l'on peut apporter} à l'entreprise. Les deux documents doivent être cohérents : la lettre ne répète pas le CV, elle le commente et le met en valeur.

\begin{savaistu}[Le CV à l'épreuve du Baccalauréat technique]
Au Baccalauréat technique, l'épreuve de production de texte peut demander la rédaction complète d'un CV à partir d'une \textbf{situation d'embauche donnée} dans le sujet (une offre d'emploi et un profil de candidat). La grille d'évaluation valorise trois critères : la \textbf{structure} (rubriques présentes et bien ordonnées), la \textbf{pertinence} (les informations répondent à l'offre, rien d'inutile) et la \textbf{correction de la langue} (orthographe, accords, vocabulaire professionnel). S'entraîner à rédiger un CV en une heure, c'est donc préparer directement l'examen — et son insertion professionnelle.
\end{savaistu}

\begin{exercice}[À vous de produire]
À partir de la situation suivante : \emph{« L'entreprise CAMWATER recrute un technicien de maintenance. Profil exigé : Baccalauréat technique F3 ou équivalent, un stage en milieu industriel, rigueur et ponctualité. »} Vous êtes Essomba Patrice, titulaire du Baccalauréat F3 (2024, Lycée technique de Ngaoundéré), avec un stage de deux mois chez un garagiste industriel. Rédigez votre CV complet en respectant les 3 C, puis échangez votre copie avec un camarade pour la relecture et produisez une version remédiée en justifiant trois modifications.
\end{exercice}

\newpage

\coursec{Activité d'intégration}

\coursec{Produire un Curriculum Vitae (CV)}

\courssub{Objectifs de la séquence}

À l'issue de cette séquence de 11 séances, l'élève de Terminale technique sera capable de :
\begin{itemize}
\item \textbf{Analyser} un texte iconique (offre d'emploi, modèle de CV) pour en extraire les informations pertinentes (profil, mots-clés, structure) ;
\item \textbf{Maîtriser} la structure externe (mise en page, présentation, longueur) et la structure interne (rubriques obligatoires et facultatives, ordre logique) d'un CV ;
\item \textbf{Distinguer} les trois types de CV (chronologique, fonctionnel, mixte) et \textbf{choisir} le plus adapté à sa situation (jeune diplômé, expérience significative, reconversion) ;
\item \textbf{Décrypter} une situation d'embauche (offre d'emploi) pour en dégager le profil recherché et les mots-clés à réinjecter dans le CV ;
\item \textbf{Rédiger} un CV complet, tenu sur \textbf{une seule page}, sans faute d'orthographe ni de syntaxe, en phrases nominales brèves et vocabulaire précis ;
\item \textbf{Adapter} le contenu (formation, stages, compétences, centres d'intérêt) à l'offre ciblée ;
\item \textbf{S'auto-évaluer} et \textbf{corriger} sa production à l'aide d'une grille critériée, puis \textbf{justifier} oralement ses choix de rédaction.
\end{itemize}
Ces objectifs couvrent l'ensemble des \emph{Savoirs} (réception de textes iconiques, structure externe/interne, analyse d'offre, rédaction, adaptation, amélioration) et des \emph{Savoir-faire} (application à la vie courante, justification de démarche) du programme officiel MINESEC pour le Français en Terminale technique.

\courssub{Plan de la séquence (11 séances)}

Le tableau ci-dessous articule les savoirs du programme aux séances prévisionnelles :

\begin{center}
\begin{tabularx}{\linewidth}{|c|l|X|}\hline
\rowcolor{popblueL} \textbf{Séance} & \textbf{Savoir programmatique} & \textbf{Activités principales} \\ \hline
1 & Texte iconique N° 1 (Observation) & Analyse de modèles de CV (bon/mauvais), découverte de la structure externe. \\ \hline
2 & Production : Structure externe & Règles de mise en page (police, interlignes, photo, marges, une page). Exercices de mise en forme. \\ \hline
3 & Texte iconique N° 2 (Observation) & Analyse de la structure interne : rubriques, ordre chronologique inversé, phrases nominales. \\ \hline
4 & Production : Structure interne (rubriques) & Rédaction guidée de chaque rubrique : État civil, Objectif, Formation, Expérience, Compétences, Langues, Centres d'intérêt, Références. \\ \hline
5 & Texte iconique N° 3 (Observation) & Étude des 3 types de CV (chronologique, fonctionnel, mixte) : avantages/inconvénients selon le profil. \\ \hline
6 & Production : Analyse situation d'embauche & Méthode de lecture d'une offre : employeur, poste, missions, profil, mots-clés, contrat, délai. \\ \hline
7 & Production : Adapter le CV à l'offre & Exercices de réécriture : faire correspondre son profil aux mots-clés de 3 offres types (agricole, tertiaire, industriel). \\ \hline
8 & Production : Rédaction du CV & Rédaction individuelle d'un CV complet sur une offre choisie (brouillon + mise au net). \\ \hline
9-10 & Production : Rédaction et amélioration & \textbf{Activité d'intégration} (travail en groupes : analyse, choix type, rédaction, relecture croisée, amélioration, oral). \\ \hline
11 & Compte-rendu et remédiation & Restitution collective des erreurs fréquentes, remédiation individuelle, évaluation formative. \\ \hline
\end{tabularx}
\end{center}

\courssub{Rappels de cours : Outils théoriques mobilisables}

\begin{definition}[Curriculum Vitae (CV)]
Document de présentation personnelle et professionnelle, normé à \textbf{une page A4}, dont la fonction unique est de \textbf{convaincre un employeur de convoquer le candidat à un entretien d'embauche}. Ce n'est pas une biographie exhaustive, mais une \emph{publicité ciblée} pour un poste précis.
\end{definition}

\begin{propriete}[Structure externe (Forme)]
Le respect de ces normes est éliminatoire si non respecté :
\begin{itemize}
\item \textbf{Format :} Une seule page A4 (recto seul).
\item \textbf{Présentation :} Aérée, marges suffisantes (2 cm min), interlignes 1,15 ou 1,5.
\item \textbf{Police :} Une seule police professionnelle (Arial, Calibri, Times New Roman, Helvetica), corps 11 ou 12 pour le texte, 14-16 pour le nom.
\item \textbf{Rubriques :} Séparées par des titres en gras/majuscules, traits de séparation ou espaces blancs nets.
\item \textbf{Photo :} Facultative (identité, fond neutre, format passeport), en haut à droite.
\item \textbf{Propreté :} Aucune rature, surcharge, tache ; impression laser/jet d'encre qualité.
\item \textbf{Fichier :} Format PDF (nommé \texttt{Nom\_Prenom\_CV.pdf}) pour envoi par courriel.
\end{itemize}
\end{propriete}

\begin{propriete}[Structure interne (Fond) — Rubriques obligatoires et ordre recommandé]
L'ordre standard (chronologique inversé pour Formation et Expérience) :
\begin{enumerate}
\item \textbf{État civil :} Nom, Prénom, Date/Lieu de naissance, Nationalité, Situation matrimoniale, Adresse, Téléphone, Courriel \textbf{professionnel} (ex: \texttt{prenom.nom@domaine.com}).
\item \textbf{Objectif professionnel (ou Titre) :} Une phrase ciblée : \emph{« Poste visé + valeur ajoutée pour l'entreprise »}. Ne pas copier l'intitulé de l'annonce.
\item \textbf{Formation :} Du plus récent au plus ancien. Diplôme, Série/Spécialité, Établissement, Année, \textbf{Mention} (si Assez bien ou supérieur), Mémoire/Projet pertinent.
\item \textbf{Expérience professionnelle :} Du plus récent au plus ancien. Dates (Mois/Année), Intitulé du poste, Structure, Localisation, \textbf{3 à 5 missions/résultats} (verbes d'action, chiffrés si possible). Stages et jobs étudiants inclus.
\item \textbf{Compétences :} Regroupées par domaines (Techniques, Informatiques, Linguistiques, Transverses). Niveau précisé (Maîtrise, Notions, Certifié).
\item \textbf{Langues :} Français (langue de travail), Anglais (niveau CECRL : A2, B1, B2...), Langues nationales (préciser : lu, écrit, parlé).
\item \textbf{Centres d'intérêt :} 3 à 4 items variés (sport collectif, bénévolat, culture, technique) montrant des soft skills (esprit d'équipe, engagement, curiosité).
\item \textbf{Références :} 2 personnes (Noms, Fonction, Structure, Téléphone, Courriel) ayant donné leur accord. Ou mention « Références transmises sur demande ».
\end{enumerate}
\end{propriete}

\begin{exemplebox}[Les trois types de CV — Choix stratégique]
\begin{center}
\begin{tabularx}{\linewidth}{|l|X|X|X|}\hline
\rowcolor{popblueL} \textbf{Critère} & \textbf{Chronologique} & \textbf{Fonctionnel} & \textbf{Mixte (Recommandé junior)} \\ \hline
\textbf{Logique} & Parcours temps (date) & Parcours par compétences & Compétences \textbf{puis} Parcours temps \\ \hline
\textbf{Public cible} & Carrière stable, évolution logique & Reconversion, trous dans CV, freelance & \textbf{Jeunes diplômés}, peu d'expérience \\ \hline
\textbf{Atout} & Lisibilité pour recruteur classique & Masque les faiblesses temporelles & \textbf{Met en avant compétences + diplôme} \\ \hline
\textbf{Risque} & Met en évidence les périodes creuses & Décontenance le recruteur traditionnel & Exige une rédaction très rigoureuse \\ \hline
\end{tabularx}
\end{center}
\textbf{Règle pour la Terminale technique :} Le \cle{CV mixte} est quasi systématiquement le meilleur choix : il place la rubrique \textbf{Compétences} (technique, informatique, langues) juste après l'Objectif et avant l'Expérience, ce qui valorise le diplôme technique fraîchement obtenu.
\end{exemplebox}

\begin{methode}[Analyser une offre d'emploi — 5 étapes]
\begin{enumerate}
\item \textbf{Identifier l'employeur :} Secteur d'activité, taille, localisation, culture (valeurs).
\item \textbf{Extraire le poste et les missions :} Verbes d'action (gérer, encadrer, câbler, saisir, accueillir). Lister les tâches quotidiennes.
\item \textbf{Isoler le profil « dur » (Hard skills) :} Diplôme exact (série, spécialité), certifications (habilitations, permis), logiciels (Excel, logiciel de caisse, CAO), techniques (câblage, rendement, comptabilité).
\item \textbf{Repérer le profil « mou » (Soft skills) et atouts :} Rigueur, contact, discrétion, mobilité, langues locales, disponibilité.
\item \textbf{Surligner 6 à 10 \cle{mots-clés} :} Ils doivent réapparaître \emph{à l'identique ou synonymes proches} dans l'Objectif, les Compétences et la description des Expériences du CV.
\end{enumerate}
\end{methode}

\begin{attention}[Erreurs fatales à proscrire (Élimination directe)]
\begin{itemize}
\item \textbf{Dépasser une page} (sauf cadreur senior, non applicable ici).
\item \textbf{Mentir} sur un diplôme, une date, une expérience (vérifiable, fraude = exclusion).
\item \textbf{Coordonnées manquantes ou erronées} (téléphone injoignable, mail fantaisiste \texttt{kingboy@...}).
\item \textbf{Fautes d'orthographe, de grammaire, de syntaxe} (CV = test de communication écrite).
\item \textbf{CV « passe-partout »} : envoyé identique à des offres différentes (absence d'adaptation = manque de motivation).
\item \textbf{Rubriques mélangées} (ex: stages dans Formation, formation dans Expérience).
\item \textbf{Abréviations non usuelles} (écrire « Baccalauréat technique » pas « Bac Tech » ; « Microsoft Excel » pas « XL »).
\item \textbf{Phrases verbales longues} : Privilégier le \textbf{style télégraphique} (phrases nominales, participes passés, verbes d'action à l'infinitif ou 1\textsuperscript{re} personne du passé composé pour les missions passées).
\end{itemize}
\end{attention}

\courssub{Activité d'intégration : Simulation de candidature (Séances 9-10)}

\courssub{Objectifs de l'activité}

À l'issue de cette activité d'intégration, l'élève sera capable de :
\begin{itemize}
\item analyser une offre d'emploi réelle afin de dégager le \cle{profil recherché} et les \cle{mots-clés} de l'annonce ;
\item choisir le \cle{type de CV} (chronologique, fonctionnel ou mixte) adapté à la situation d'embauche ;
\item rédiger un CV complet d'\textbf{une page} respectant la structure externe (mise en page, présentation) et la structure interne (rubriques) ;
\item adapter le contenu du CV à l'offre (formation, expériences, compétences mises en avant) ;
\item relire et évaluer le CV d'un camarade à l'aide d'une grille de critères, puis améliorer sa propre production ;
\item justifier oralement ses choix de rédaction et rendre compte collectivement des erreurs fréquentes.
\end{itemize}
Cette activité mobilise l'ensemble des savoirs de l'unité : l'observation des textes iconiques, la structure externe et interne du CV, l'analyse de la situation d'embauche et les techniques de rédaction et d'amélioration.

\courssub{Situation déclenchante}

Le journal \emph{Cameroun Tribune} du mardi 14 octobre 2025, ainsi que le site \emph{emploi.cm}, publient plusieurs offres d'emploi destinées à des jeunes diplômés de l'enseignement technique. Le professeur de Français de la classe de Terminale technique du Lycée technique de Bafoussam distribue trois de ces offres authentiques, reproduites ci-dessous.

\textbf{Offre n° 1 — Coopérative des producteurs de maïs de l'Ouest (COPROMAO), Bafoussam.} La COPROMAO, qui regroupe 1\,250 exploitants agricoles répartis dans 8 villages du département de la Mifi, recrute \textbf{un(e) technicien(ne) agricole}. Missions : encadrement technique des producteurs, suivi des parcelles (environ 600 hectares), tenue des fiches de rendement, formation des paysans aux techniques de conservation. Profil : Baccalauréat technique série F3 ou BT Productions végétales ; permis B exigé ; bonne connaissance de la zone de l'Ouest ; sens du contact ; maîtrise du français, le fulfuldé ou le ghomala étant un atout. Contrat : CDD de 12 mois renouvelable, salaire mensuel de 150\,000 FCFA. Dossier : CV + lettre de motivation à déposer avant le 30 novembre 2025.

\textbf{Offre n° 2 — Mutuelle communautaire de croissance (MC\textsuperscript{2}), Douala-Bonabéri.} Cette microfinance, forte de 14\,300 clients et de 6 agences dans la région du Littoral, recrute \textbf{un(e) assistant(e) de gestion}. Missions : accueil des clients, saisie des opérations sur logiciel, classement des dossiers de crédit, élaboration des rapports hebdomadaires. Profil : Baccalauréat technique série G1/G2 ou BT Comptabilité ; maîtrise de Word et Excel exigée ; rigueur, discrétion, honnêteté ; une première expérience (stage compris) est souhaitée. Contrat : stage probatoire de 3 mois puis CDI, rémunération de départ de 100\,000 FCFA. Dossier : CV uniquement, à envoyer par courriel avant le 25 novembre 2025.

\textbf{Offre n° 3 — Entreprise ETS Ngoa Électricité, Yaoundé.} Cette entreprise de 35 salariés, spécialisée dans les installations électriques du bâtiment, recrute \textbf{deux (2) électriciens bâtiment}. Missions : câblage des immeubles neufs, pose de tableaux électriques, dépannage chez les particuliers. Profil : Baccalauréat technique série F2 ou BT Électrotechnique ; 2 ans d'expérience souhaités (chantiers ou stages longs acceptés) ; habilitations électriques appréciées ; disponibilité immédiate ; déplacements fréquents dans la ville de Yaoundé. Salaire : 120\,000 FCFA + primes de chantier. Dossier : CV + photocopie du diplôme, à déposer au siège de l'entreprise, quartier Ngoa-Ekellé, avant le 28 novembre 2025.

\textbf{Problème posé :} Tu es un jeune diplômé (ou tu incarnes un candidat fictif) qui souhaite postuler à l'une de ces trois offres. Comment produire un CV d'une seule page, correctement structuré et parfaitement adapté à l'offre choisie, de manière à maximiser tes chances d'être convoqué à un entretien ?

\courssub{Consignes de travail (Groupes de 4 élèves)}

\begin{enumerate}
\item \textbf{Analyse de l'offre (15 min).} Choisissez une des trois offres. Remplissez le tableau d'analyse : employeur, poste, missions, diplôme exigé, compétences (hard/soft), atouts, contrat, date limite. Soulignez au moins \textbf{six mots-clés} stratégiques.
\item \textbf{Choix du type de CV (5 min).} Parmi chronologique, fonctionnel, mixte, indiquez celui qui convient le mieux à votre candidat fictif. Justifiez en 2-3 phrases (profil junior, compétences techniques fortes, expérience stage).
\item \textbf{Rédaction du CV — Version 1 (30 min).} Rédigez le CV complet du candidat sur \textbf{une seule page} (format A4 fourni ou traitement de texte). Rubriques minimales obligatoires : État civil, Objectif professionnel, Formation, Expérience professionnelle (stages acceptés), Compétences, Langues, Centres d'intérêt, Références (2).
\item \textbf{Relecture croisée — Grille d'évaluation (20 min).} Échangez avec un groupe ayant choisi une \textbf{autre} offre. Évaluez le CV reçu sur 20 à l'aide de la grille ci-dessous. Annotez le document.
\item \textbf{Amélioration — Version 2 (15 min).} Récupérez votre CV et les annotations. Corrigez, améliorez, aérez. Sur une feuille annexe, listez les 3-4 modifications majeures et leur justification (lien avec l'offre, correction langue, mise en page).
\item \textbf{Justification orale (3 min/groupe).} Présentez le CV final : type choisi, rubriques valorisées, 3 mots-clés de l'offre réinjectés, difficulté rencontrée.
\item \textbf{Compte rendu collectif et remédiation (15 min).} Le professeur mène la synthèse au tableau : erreurs types, critères de réussite, remédiation ciblée (orthographe, objectif professionnel, adaptation).
\end{enumerate}

\courssub{Grille d'évaluation du CV (Relecture croisée — Sur 20)}

\begin{center}
\begin{tabularx}{\linewidth}{|l|c|X|}\hline
\rowcolor{popblueL} \textbf{Critère} & \textbf{Bareme} & \textbf{Indicateurs de réussite (Cocher / Noter)} \\ \hline
\textbf{1. Structure externe} & 5 pts & \begin{itemize}\item[$\square$] Une page unique \item[$\square$] Aération, lisibilité \item[$\square$] Police unique, titres visibles \item[$\square$] Photo pro ou absence \item[$\square$] Aucune rature/faute de frappe visuelle\end{itemize} \\ \hline
\textbf{2. Structure interne} & 5 pts & \begin{itemize}\item[$\square$] 8 rubriques présentes \item[$\square$] Ordre logique (État civil $\rightarrow$ Références) \item[$\square$] Chronologie inversée (Formation/Exp.) \item[$\square$] Phrases nominales, verbes d'action \item[$\square$] Dates alignées (Mois/Année)\end{itemize} \\ \hline
\textbf{3. Adaptation à l'offre} & 5 pts & \begin{itemize}\item[$\square$] Objectif professionnel ciblé \item[$\square$] 5+ mots-clés de l'offre réinjectés \item[$\square$] Compétences techniques matching \item[$\square$] Expériences/Stages pertinents mis en avant \item[$\square$] Centres d'intérêt en lien (soft skills)\end{itemize} \\ \hline
\textbf{4. Qualité de la langue} & 5 pts & \begin{itemize}\item[$\square$] Zéro faute orthographe/grammaire \item[$\square$] Syntaxe correcte, style télégraphique \item[$\square$] Vocabulaire précis (technique) \item[$\square$] Pas d'abréviations obscures \item[$\square$] Ponctuation cohérente\end{itemize} \\ \hline
\textbf{TOTAL} & \textbf{/20} & \textbf{Appréciation globale :} \hrulefill \\ \hline
\end{tabularx}
\end{center}

\courssub{Démarche et corrigé commenté (Exemple Offre n° 2)}

\begin{exoresolu}[Production d'un CV adapté à l'offre n° 2 — Assistant(e) de gestion MC\textsuperscript{2}]
Énoncé : Le groupe A a choisi l'offre de la microfinance MC\textsuperscript{2} de Douala. Produire l'analyse de l'offre, le choix du type de CV, le CV rédigé (V1), la grille d'évaluation remplie, les améliorations (V2) et la justification orale.

\tcblower

\textbf{Étape 1 — Analyse de l'offre et dégagement du profil recherché.}

Le groupe complète le tableau d'analyse :

\begin{center}
\begin{tabularx}{\linewidth}{|l|X|}\hline
\rowcolor{popblueL} \textbf{Élément} & \textbf{Données relevées dans l'offre} \\ \hline
Employeur & Mutuelle communautaire de croissance (MC\textsuperscript{2}), Douala-Bonabéri \\ \hline
Poste & Assistant(e) de gestion \\ \hline
Missions & Accueil des clients, saisie des opérations sur logiciel, classement des dossiers de crédit, élaboration des rapports hebdomadaires \\ \hline
Diplôme exigé & Baccalauréat technique G1/G2 ou BT Comptabilité \\ \hline
Compétences techniques (Hard) & Maîtrise de Word et Excel, saisie sur logiciel de gestion, comptabilité de base \\ \hline
Compétences transverses (Soft) & Rigueur, discrétion, honnêteté, sens de l'accueil \\ \hline
Atout différenciant & Une première expérience (stage compris) en structure financière \\ \hline
Contrat et salaire & Stage probatoire 3 mois $\rightarrow$ CDI ; 100\,000 FCFA \\ \hline
Modalités & Envoi CV par courriel uniquement, date limite 25 novembre 2025 \\ \hline
\end{tabularx}
\end{center}

\textbf{Mots-clés stratégiques soulignés (8) :} \cle{gestion}, \cle{accueil}, \cle{saisie}, \cle{logiciel}, \cle{Word}, \cle{Excel}, \cle{rigueur}, \cle{rapports}, \cle{comptabilité}, \cle{crédit}.

\textbf{Étape 2 — Choix du type de CV.}

Le groupe choisit le \cle{CV mixte}. 
\textbf{Justification :} Notre candidate est une jeune diplômée (Bac G2 2025) avec deux stages significatifs mais pas d'emploi stable. Le CV mixte permet de placer la rubrique \textbf{Compétences} (comptabilité, bureautique, logiciel caisse) \emph{avant} la rubrique Expérience, ce qui capte l'attention du recruteur sur les savoir-faire opérationnels immédiatement transférables. La Formation reste en chronologique inversé car le diplôme est l'atout principal « frais ».

\textbf{Étape 3 — Rédaction du CV (Version 1 — Brouillon corrigé pour l'exemple).}

\begin{center}
\begin{tabularx}{\linewidth}{|X|}\hline
\rowcolor{popgreenL} \textbf{MENDOUGA Aïcha Laure} \\ \hline
Née le 12 mars 2006 à Douala — Célibataire — Nationale camerounaise \\
Tél. : 6 90 45 67 89 — Courriel : \texttt{aicha.mendouga@mail.com} \\
Quartier Bonabéri, Douala (Proche agence MC\textsuperscript{2}) \\ \hline
\textbf{OBJECTIF PROFESSIONNEL} \\
Assistante de gestion rigoureuse, mettre ma maîtrise de la saisie comptable et des outils bureautiques (Excel, logiciel caisse) au service de la performance de la MC\textsuperscript{2}. \\ \hline
\textbf{FORMATION} \\
\textbf{2025} : Baccalauréat Technique Série G2 (Gestion Comptable), Mention Assez Bien. Lycée Technique de Douala-Bassa. \\
\textbf{2022} : Probatoire Technique Série G2. Lycée Technique de Douala-Bassa. \\ \hline
\textbf{COMPÉTENCES} \\
\textit{Comptabilité / Gestion :} Saisie comptable (recettes/dépenses), tenue de livres journaux, élaboration de rapports d'activité, classement dossiers crédit. \\
\textit{Informatique :} Maîtrise de la suite Microsoft Office (Word : mise en forme documents ; Excel : tableaux croisés dynamiques, formules SI/SOMME.SI, graphiques), Saisie rapide sur logiciel de gestion de caisse (stage). \\
\textit{Relationnel :} Accueil clientèle, écoute active, discrétion professionnelle, gestion des réclamations 1\textsuperscript{er} niveau. \\ \hline
\textbf{EXPÉRIENCE PROFESSIONNELLE (Stages)} \\
\textbf{Juillet – Septembre 2024 (3 mois) :} Stagiaire Assistante de Gestion — \textbf{Caisse Populaire de Bonabéri}. \\
\hspace{0.5cm}• Accueil physique et téléphonique de 30 à 40 clients/jour ; orientation vers guichets. \\
\hspace{0.5cm}• Saisie de 150+ opérations quotidiennes sur logiciel de gestion de caisse (dépôts, retraits, virements). \\
\hspace{0.5cm}• Classement alphanumérique des dossiers de crédit (archivage physique et numérique). \\
\hspace{0.5cm}• Participation à l'élaboration du rapport hebdomadaire d'activité (synthèse Excel). \\
\textbf{Juillet 2023 (1 mois) :} Stagiaire Comptable — \textbf{Établissement Ngoné Commerce de Gros (Douala)}. \\
\hspace{0.5cm}• Saisie des factures d'achat/vente sur Excel ; lettrage comptes fournisseurs. \\
\hspace{0.5cm}• Inventaire physique mensuel du stock (matériaux de construction). \\ \hline
\textbf{LANGUES} \\
Français : Langue maternelle / Courant (rédaction administrative). \\
Anglais : Niveau B1 (CECRL) — Compréhension écrite technique, oral fonctionnel. \\
Ghomala : Notions orales. \\ \hline
\textbf{CENTRES D'INTÉRÊT} \\
Football féminin (équipe lycée — Capitaine : esprit d'équipe, leadership) ; Lecture (développement personnel, gestion) ; Bénévolat : Trésorière Club Jeunes Entrepreneurs Lycée (gestion caisse, rapport moral). \\ \hline
\textbf{RÉFÉRENCES} \\
M. ETOA Jean, Chef Comptable, Caisse Populaire Bonabéri — Tél. : 6 99 00 11 22 — \texttt{jean.etoa@caisse-cm.com} \\
Mme NGO BELL Sylvie, Proviseure, Lycée Technique Douala-Bassa — Tél. : 6 77 33 44 55 \\ \hline
\end{tabularx}
\end{center}

\textbf{Étape 4 — Relecture croisée avec la grille (Simulation Groupe B évaluant Groupe A).}

\begin{center}
\begin{tabularx}{\linewidth}{|l|c|X|}\hline
\rowcolor{popblueL} \textbf{Critère} & \textbf{Note} & \textbf{Observations détaillées du groupe relecteur} \\ \hline
Structure externe (5) & 4/5 & Présentation claire, photo absente (choix assumé), police unique. \textbf{Point faible :} interlignes un peu justes entre rubriques « Compétences » et « Expérience », lecture rapide difficile. \\ \hline
Structure interne (5) & 5/5 & 8 rubriques complètes, ordre parfait (Compétences avant Expérience = Mixte respecté). Chronologie inversée respectée. Dates Mois/Année alignées. \\ \hline
Adaptation à l'offre (5) & 4/5 & Objectif ciblé « MC\textsuperscript{2} ». Mots-clés : \cle{gestion}, \cle{saisie}, \cle{Excel}, \cle{rapports}, \cle{rigueur}, \cle{accueil}, \cle{crédit}, \cle{logiciel} (8/10 repris). \textbf{Manque :} le mot « discrétion » n'apparaît pas explicitement dans Compétences (sous-entendu). \\ \hline
Langue (5) & 4/5 & Style télégraphique maîtrisé. \textbf{Erreur majeure :} « \textbf{élaborration} » (double r) dans l'expérience Caisse Populaire (rapport hebdomadaire). Accords corrects par ailleurs. \\ \hline
\textbf{TOTAL} & \textbf{17/20} & \textbf{CV convaincant, prêt pour entretien après corrections mineures.} \\ \hline
\end{tabularx}
\end{center}

\textbf{Étape 5 — Amélioration du CV (Version 2 finale).}

Trois modifications apportées et justifiées sur feuille annexe :
\begin{enumerate}
\item \textbf{Correction orthographique :} « élaborration » $\rightarrow$ « \textbf{élaboration} » (rapport hebdomadaire). \textit{Justification :} Zéro tolérance faute ; crédibilité gestion.
\item \textbf{Enrichissement mots-clés (Soft skill) :} Ajout de « \textbf{Discrétion professionnelle} » en tête de liste dans Compétences \textit{Relationnel}. \textit{Justification :} Mot-clé explicite de l'offre (honnêteté, discrétion) absent en V1.
\item \textbf{Aération mise en page :} Augmentation interligne 1,5 entre chaque rubrique majeure ; gras sur intitulés de stage. \textit{Justification :} Lisibilité en < 30 sec (test « scan » recruteur).
\item \textbf{Précision technique :} Dans stage 2024, « logiciel de gestion de caisse » remplacé par « \textbf{Logiciel Microfinance \texttt{MF-Soft} (saisie, édition reçus)} ». \textit{Justification :} Preuve concrète maîtrise « saisie sur logiciel » (mot-clé dur).
\end{enumerate}

\textbf{Étape 6 — Justification orale (Extrait synthétique 3 min).}

\begin{quote}
« Nous avons choisi le \textbf{CV mixte} car Aïcha, jeune diplômée Bac G2, doit vendre ses \textbf{compétences techniques} (comptabilité, Excel, logiciel caisse) avant son expérience encore courte. L'\textbf{Objectif} cite nommément la MC\textsuperscript{2} et les outils. Nous avons réinjecté \textbf{8 mots-clés} de l'offre (gestion, accueil, saisie, logiciel, Excel, rapports, rigueur, crédit) dans Objectif, Compétences et Expérience. La relecture a permis de corriger une faute grave (« élaborration ») et d'ajouter « discrétion » exigée par l'annonce. Le CV tient en une page, aéré, prêt à être converti en PDF pour envoi mail avant le 25/11. »
\end{quote}

\textbf{Étape 7 — Compte rendu collectif (Animé par le professeur).}

La classe dégage les \textbf{erreurs fréquentes} observées sur l'ensemble des groupes :
\begin{itemize}
\item CV de 2 pages (inclusion de détails superflus : primaire, collège, hobbies longs).
\item Coordonnées incomplètes (pas de mail, mail fantaisiste, indicatif pays manquant +237).
\item Rubrique « Objectif » absente, ou copiée mot pour mot sur l'annonce (« Je veux être assistant de gestion »).
\item Expériences : simples listes de tâches sans verbes d'action, sans résultats chiffrés.
\item Compétences : liste générique (« Informatique : Word, Excel ») sans niveau ni contexte.
\item Adaptation nulle : même CV pour COPROMAO (agricole) et MC\textsuperscript{2} (gestion).
\item Fautes d'accords (participes passés, accord adjectifs) et anglicismes (« \textit{manager} » au lieu de « gérer/encadrer »).
\end{itemize}

\textbf{Critères d'un CV réussi (Check-list finale affichée) :}
\begin{enumerate}
\item \textbf{Une page} — \textbf{Lisible en 30 sec} (scan en F / Z) — \textbf{Zéro faute}.
\item \textbf{Adapté} : Mots-clés de \textbf{cette} offre présents dans Objectif, Compétences, Expérience.
\item \textbf{Structure interne complète} : 8 rubriques, ordre logique, chronologie inversée.
\item \textbf{Structure externe pro} : Police unique, aération, PDF, nom de fichier pro.
\item \textbf{Honnête et vérifiable} : Dates, diplômes, références joignables.
\end{enumerate}
La \textbf{remédiation individuelle} (Séance 11) porte sur : reformulation de l'objectif professionnel (passer du « je veux » au « j'apporte »), correction orthographique ciblée (liste personnelle), réécriture d'une expérience en style télégraphique chiffré.

\end{exoresolu}

\courssub{Bilan de la séquence et articulation avec l'examen}

\begin{aretenir}[Acquis essentiels pour le Probatoire / Baccalauréat Technique]
\begin{itemize}
\item \textbf{Le CV est un acte de communication persuasive,} pas une fiche d'état civil. Chaque mot doit « vendre » le candidat pour \emph{ce} poste précis.
\item \textbf{L'analyse de l'offre est la boussole :} Elle dicte le type de CV, l'ordre des rubriques, le vocabulaire (mots-clés) et les expériences à valoriser. \emph{Un candidat = autant de CV différents que d'offres ciblées.}
\item \textbf{La forme conditionne la lecture :} Une page, aérée, structurée, sans faute = respect du recruteur et preuve de rigueur (compétence transverse n°1).
\item \textbf{L'écriture professionnelle} : Style télégraphique (phrases nominales, verbes d'action, chiffrage), vocabulaire technique précis, zéro abréviation obscure.
\item \textbf{La révision est constitutive de l'écriture :} Relecture croisée, grille critériée, correction, justification orale = processus complet d'un écrivain professionnel.
\item \textbf{Compétence évaluée à l'examen :} « Rédiger et justifier une démarche, une réponse ou une conclusion » (Savoir-faire officiel). L'oral de justification (Étape 6) en est l'entraînement direct.
\end{itemize}
\end{aretenir}

\begin{savaistu}[Insertion professionnelle au Cameroun — Contexte réel]
\begin{itemize}
\item \textbf{Secteurs porteurs pour Bac Tech :} Agro-industrie (F3), BTP/Électricité (F2, F1), Tertiaire/Gestion/Comptabilité (G1, G2, G3), Maintenance/Industrie (F4, F5, F6, F7).
\item \textbf{Canaux de recrutement :} \emph{Cameroun Tribune} (mercredi), \emph{emploi.cm}, \emph{cdmt.jobs}, LinkedIn (créer profil pro), réseaux alumni lycée/écoles, dépôts physiques (PME).
\item \textbf{Attentes employeurs :} Pratique immédiate (stages, projets techniques), maîtrise outils numériques (Excel avancé, logiciels métiers), mobilité (permis B), langues (Français/Anglais + locale), soft skills (ponctualité, honnêteté, travail en équipe).
\item \textbf{Stage probatoire (3-6 mois) :} Porte d'entrée quasi systématique pour primo-emploi (Offre n°2 type). Le CV doit rassurer sur la \textbf{formabilité} et la \textbf{fiabilité}.
\end{itemize}
\end{savaistu}

\begin{methode}[Check-list finale avant envoi (Auto-évaluation)]
\begin{enumerate}
\item \textbf{Contenu :} Offre relue ? Mots-clés cochés ? Objectif ciblé ? Compétences match ? Expériences pertinentes ? Références prévenues ?
\item \textbf{Forme :} 1 page ? Police unique 11/12 ? Marges 2cm ? Titres rubriques gras/majuscules ? Dates alignées ? Photo pro ou rien ? PDF généré ?
\item \textbf{Langue :} Relu à voix haute ? Correcteur ortho (LanguageTool / Antidote) passé ? 2ème relecture par tiers ? Zéro abréviation ? Verbes d'action variés ?
\item \textbf{Envoi :} Objet mail : \texttt{Candidature - [Intitulé Poste] - [Nom Prénom] - Ref [N° offre]}. Corps mail poli, 5 lignes max. Pièce jointe : \texttt{Nom\_Prenom\_CV.pdf} (+ LM si demandée). Envoi avant date limite (J-2 idéal).
\end{enumerate}
\end{methode}

\courssub{Exercices d'entraînement (Devoirs / Remédiation)}

\begin{exercice}[Analyse d'offre — Entraînement individuel]
À partir de l'\textbf{Offre n° 1 (COPROMAO)} ou de l'\textbf{Offre n° 3 (ETS Ngoa Électricité)} :
\begin{enumerate}
\item Remplir le tableau d'analyse (Étape 1).
\item Lister 8 mots-clés techniques et transverses.
\item Choisir le type de CV adapté pour un candidat Bac F3 (Offre 1) ou Bac F2 (Offre 3) avec 2 stages. Justifier.
\item Rédiger la rubrique \textbf{Compétences} (3 domaines : Technique, Informatique/Outils, Relationnel/Transverse) en reprenant les mots-clés.
\end{enumerate}
\end{exercice}

\begin{exercice}[Rédaction ciblée — Objectif professionnel]
Réécrire les objectifs professionnels suivants pour qu'ils soient \textbf{ciblés, orientés résultat et sans « je veux »} :
\begin{enumerate}
\item « Je souhaite travailler dans votre entreprise comme électricien. »
\item « Je cherche un poste de technicien agricole pour gagner de l'expérience. »
\item « Je veux être assistant de gestion dans une microfinance. »
\end{enumerate}
\textit{Indication :} Structure = [Poste visé] + [Compétence clé 1] + [Compétence clé 2] + [Valeur pour l'employeur].
\end{exercice}

\begin{exercice}[Correction de CV — Détection d'erreurs]
Le CV ci-dessous contient \textbf{7 erreurs} (forme, fond, langue, adaptation). Les identifier et proposer la correction.
\begin{quote}
\textbf{NGONO Paul} — 06 12 34 56 78 — \texttt{paul.ngono@yahoo.fr} — Yaoundé \\
\textbf{OBJECTIF} : Je veux postuler à l'offre d'électricien. \\
\textbf{FORMATION} : 2025 Bac F2. 2022 Probatoire. 2019 BEPC. École primaire 2013. \\
\textbf{EXPÉRIENCE} : 2024 Stage ENEO (câblage). 2023 Vacances : vendeur boutique. \\
\textbf{COMPÉTENCES} : Électricité, Word, Excel, Français, Anglais. \\
\textbf{LANGUES} : Français, Anglais. \\
\textbf{LOISIRS} : Foot, musique, sortir entre amis, dormir. \\
\textbf{RÉFÉRENCES} : Mon papa, mon prof.
\end{quote}
\end{exercice}

\begin{corrige}[Exercice 3 — Corrections types]
\begin{enumerate}
\item \textbf{Mail non pro} (\texttt{yahoo.fr} + prénom seul) $\rightarrow$ \texttt{paul.ngono@domaine.pro} ou Gmail/Outlook prénom.nom.
\item \textbf{Objectif} : « Je veux postuler... » $\rightarrow$ « Électricien bâtiment (Habilitation B1V), mettre ma pratique du câblage tertiaire et du dépannage au service de la qualité ETS Ngoa Électricité. »
\item \textbf{Formation} : Supprimer BEPC et Primaire (inutiles post-Bac). Garder Bac F2 (mention si $\ge$ AB) + Projet fin d'études pertinent.
\item \textbf{Expérience} : Stage ENEO : détailler 3 missions chiffrées (ex: « Câblage 20 tableaux divisionnaires », « Pose 150 prises RJ45 »). Supprimer vendeur boutique (hors cible) ou ne garder que si compétence transférable (clientèle).
\item \textbf{Compétences} : Trop vague. Décliner : \textit{Technique :} Câblage normes NF C 15-100, Lecture plans, Habilitation B1V (en cours), Dépannage. \textit{Outils :} Multimètre, Perceuse, Logiciel devis \texttt{Batiprix}. \textit{Transverse :} Respect consignes sécurité, Travail en équipe, Permis B.
\item \textbf{Loisirs} : « Dormir » proscrit. Foot $\rightarrow$ « Football club (défenseur : rigueur, endurance) ». Musique $\rightarrow$ « Chant choral (écoute, coordination) ».
\item \textbf{Références} : « Mon papa » invalide (conflit d'intérêt). Remplacer par : M. TCHOUA, Maître de stage ENEO, Tél. ; Mme BEKONO, Professeur Électrotechnique, Tél.
\end{enumerate}
\end{corrige}

\newpage

\coursec{Méthodes et exercices résolus}

\coursec{Découverte du CV : définition et fonctions}

\courssub{Méthode 1 : Appliquer les notions de l'unité à une situation concrète — analyser un texte iconique (un modèle de CV)}

\begin{methode}[Lire et analyser un CV (texte iconique n°1)]
Un CV est un \cle{texte iconique} : il combine texte et mise en page. Pour l'analyser de façon méthodique :
\begin{enumerate}
\item \textbf{Identifier la nature du document} : c'est un document de communication professionnelle, à visée persuasive (obtenir un entretien d'embauche).
\item \textbf{Décrire la structure externe} : la disposition des blocs sur la page (en-tête, colonnes, encadrés), la typographie (gras, tailles), les couleurs, la photo éventuelle, la longueur (une page maximum).
\item \textbf{Repérer les rubriques (structure interne)} : état civil / informations personnelles, objectif professionnel, formation, expériences professionnelles, compétences, langues, centres d'intérêt, références.
\item \textbf{Observer la langue} : style télégraphique (phrases nominales, sans article ni sujet), verbes d'action à l'infinitif ou au participe passé, vocabulaire technique du métier visé.
\item \textbf{Évaluer l'adéquation} : le CV correspond-il au poste visé ? Les informations sont-elles vérifiables, datées, ordonnées (ordre chronologique inverse) ?
\item \textbf{Conclure} par un jugement argumenté : points forts, points faibles, pistes d'amélioration.
\end{enumerate}
\end{methode}

\begin{exoresolu}[Analyse d'un modèle de CV — Texte iconique n°1]
\textbf{Énoncé.} On soumet à votre classe le CV suivant, rédigé par une candidate à un poste de secrétaire comptable dans une PME de Douala :

\begin{center}
\begin{tabularx}{\linewidth}{|l|X|}
\hline
\rowcolor{popblueL} \textbf{Rubrique} & \textbf{Contenu} \\
\hline
En-tête & NGOUMO Aïcha — Douala, Akwa — Tél. : 6 90 00 00 00 — aicha.ngoumo@mail.fr \\
\hline
Objectif & Obtenir un poste de secrétaire comptable \\
\hline
Formation & 2024 : Baccalauréat G2 (Lycée technique de Douala) ; 2021 : Probatoire G2 \\
\hline
Expériences & 2023-2024 : stage de 3 mois, cabinet comptable FIDEX (saisie des factures, classement) \\
\hline
Compétences & Word, Excel, logiciel Sage ; sens de l'organisation \\
\hline
Langues & Français (courant), anglais (scolaire) \\
\hline
\end{tabularx}
\end{center}

Analysez ce CV en étudiant sa structure externe, sa structure interne et sa langue, puis dites s'il est adapté à la situation d'embauche.
\tcblower
\textbf{Solution détaillée.}

\textbf{1. Nature et fonction du document.} Il s'agit d'un curriculum vitae, document iconique de recherche d'emploi. Sa fonction est persuasive : convaincre un recruteur d'accorder un entretien à la candidate.

\textbf{2. Structure externe.} Le document est présenté sous forme de tableau à deux colonnes (rubriques / contenus), ce qui rend la lecture rapide : le recruteur repère en quelques secondes chaque information. L'en-tête regroupe les coordonnées en haut du document, conformément à l'usage. On peut toutefois regretter l'absence de photo (facultative) et l'absence de titre visible « Curriculum Vitae » ou de mention du poste en grand caractère.

\textbf{3. Structure interne.} Toutes les rubriques essentielles sont présentes : état civil et coordonnées, objectif professionnel, formation, expériences, compétences, langues. La formation et l'expérience sont datées et présentées en \cle{ordre chronologique inverse} (2024 avant 2021), ce qui met en avant les acquis les plus récents. Il manque cependant les rubriques « centres d'intérêt » et « références », qui humanisent et sécurisent la candidature.

\textbf{4. La langue.} Le style est télégraphique, adapté au genre : phrases nominales (« stage de 3 mois »), pas de phrases rédigées à la première personne, vocabulaire technique du métier (« saisie des factures », « logiciel Sage »). L'objectif est formulé en une phrase brève et ciblée.

\textbf{5. Adéquation à la situation d'embauche.} Le poste visé est celui de secrétaire comptable : la formation (Bac G2, filière gestion-comptabilité), le stage en cabinet comptable et la maîtrise de Sage correspondent exactement au profil recherché. Le CV est donc bien adapté.

\textbf{Conclusion.} Ce CV est clair, complet dans ses rubriques principales et cohérent avec le poste visé. Pour l'améliorer, la candidate pourrait ajouter des références professionnelles et préciser les résultats obtenus pendant son stage (par exemple le volume de factures traitées).
\end{exoresolu}

\coursec{La structure externe du CV}

\courssub{Méthode 2 : Construire la structure externe du CV}

\begin{methode}[Organiser la présentation matérielle du CV]
\begin{enumerate}
\item \textbf{Limiter la longueur} : une seule page pour un jeune diplômé de Terminale technique.
\item \textbf{Placer l'en-tête en haut} : nom en capitales, prénom, adresse, téléphone, courriel professionnel (prénom.nom@...).
\item \textbf{Choisir une disposition lisible} : une seule colonne ou deux colonnes (coordonnées et compétences à gauche, formation et expériences à droite) ; aérer avec des blancs.
\item \textbf{Hiérarchiser par la typographie} : titres de rubriques en gras et en taille supérieure, une seule police sobre (Arial, Calibri), pas plus de deux couleurs.
\item \textbf{Soigner la finition} : alignements réguliers, puces identiques, aucune rature, impression sur papier blanc de qualité ou envoi en PDF.
\end{enumerate}
\textbf{Réflexe Bac} : au brouillon, tracer d'abord le plan de la page (schéma des blocs) avant de rédiger.
\end{methode}

\begin{exoresolu}[Schématiser la structure externe d'un CV — Texte iconique n°2]
\textbf{Énoncé.} Essomo, élève de Terminale F3 (génie civil), doit envoyer son CV à une entreprise de BTP de Yaoundé. Représentez sous forme de schéma (plan de page) la structure externe que vous lui conseillez, en justifiant chaque choix.
\tcblower
\textbf{Solution détaillée.}

\textbf{Plan de page proposé :}

\begin{center}
\begin{tabularx}{\linewidth}{|X|}
\hline
\rowcolor{popblueL} \textbf{ESSOMO Jean-Marie} — CV : technicien en génie civil \\
Adresse — Téléphone — Courriel — (photo) \\
\hline
\textbf{Objectif professionnel} (2 lignes maximum) \\
\hline
\textbf{Formation} (ordre chronologique inverse) \\
\hline
\textbf{Expériences professionnelles / stages} \\
\hline
\textbf{Compétences techniques} (logiciels de dessin, lecture de plans) \\
\hline
\textbf{Langues — Centres d'intérêt — Références} \\
\hline
\end{tabularx}
\end{center}

\textbf{Justifications.}
\begin{enumerate}
\item \textbf{En-tête en haut} : le recruteur doit identifier et contacter le candidat immédiatement ; le nom en capitales facilite le classement alphabétique.
\item \textbf{Mention du poste dans le titre} : « technicien en génie civil » annonce d'emblée l'adéquation avec l'entreprise de BTP.
\item \textbf{Objectif juste sous l'en-tête} : c'est la première phrase lue ; elle doit capter l'attention.
\item \textbf{Formation avant expériences} : pour un jeune diplômé sans longue expérience, le diplôme (Bac F3) est l'atout principal ; on le place donc en premier.
\item \textbf{Une seule page, blocs aérés} : un recruteur consacre moins d'une minute à un premier tri ; la lisibilité est décisive.
\item \textbf{Photo facultative} : si elle est mise, elle doit être d'identité, récente et professionnelle.
\end{enumerate}

\textbf{Conclusion.} Cette structure externe en une page, hiérarchisée de haut en bas (identité, objectif, formation, expériences, compétences, informations complémentaires), garantit une lecture rapide et met en valeur le profil d'Essomo pour le poste visé.
\end{exoresolu}

\coursec{La structure interne du CV : les rubriques}

\courssub{Méthode 3 : Rédiger les rubriques (structure interne)}

\begin{methode}[Remplir chaque rubrique du CV]
\begin{enumerate}
\item \textbf{État civil / coordonnées} : nom, prénom(s), date et lieu de naissance (facultatifs), adresse, téléphone, courriel. Jamais de faux renseignements.
\item \textbf{Objectif professionnel} : une phrase nominale indiquant le poste visé et l'atout principal (« Obtenir un poste de... en mettant à profit... »).
\item \textbf{Formation} : diplômes du plus récent au plus ancien, avec année, intitulé exact, établissement, mention éventuelle.
\item \textbf{Expériences} : stages et emplois, du plus récent au plus ancien ; pour chacun : dates, poste, nom de l'entreprise, 2 ou 3 tâches réalisées avec des \cle{verbes d'action} (saisir, installer, réparer, rédiger, contrôler).
\item \textbf{Compétences} : techniques (logiciels, machines, procédés) et qualités (ponctualité, esprit d'équipe) — uniquement celles utiles au poste.
\item \textbf{Langues} : langue + niveau honnête (courant, scolaire, notions).
\item \textbf{Centres d'intérêt et références} : activités révélatrices de qualités ; personnes pouvant recommander le candidat (avec leur accord).
\end{enumerate}
\textbf{Formule à retenir} : une rubrique = un titre en gras + des entrées datées, en style télégraphique, classées en ordre chronologique inverse.
\end{methode}

\begin{exoresolu}[Rédiger les rubriques « Formation » et « Expériences » — Texte iconique n°3]
\textbf{Énoncé.} Voici les informations, données en désordre, concernant Mbia Larissa, candidate à un poste d'assistante en gestion : « en 2022, BEPC au collège de Bafoussam ; de septembre 2024 à décembre 2024, stage à la SODECOTON comme aide-comptable : saisie des écritures, archivage des pièces ; en 2025, baccalauréat G1 au lycée technique de Bafoussam ; en 2023, probatoire G1 ».

Rédigez les rubriques « Formation » et « Expériences professionnelles » du CV de Larissa.
\tcblower
\textbf{Solution détaillée.}

\textbf{Étape 1 : trier les informations.} Les diplômes relèvent de la rubrique « Formation » (BEPC 2022, Probatoire 2023, Baccalauréat 2025) ; le stage relève de la rubrique « Expériences professionnelles ».

\textbf{Étape 2 : classer en ordre chronologique inverse} (le plus récent d'abord), car le recruteur veut connaître d'abord le niveau actuel de la candidate.

\textbf{Étape 3 : rédiger en style télégraphique}, avec verbes d'action pour les tâches.

\textbf{Rédaction :}

\textbf{Formation}
\begin{itemize}
\item 2025 : Baccalauréat technique G1 (gestion), Lycée technique de Bafoussam.
\item 2023 : Probatoire G1, Lycée technique de Bafoussam.
\item 2022 : BEPC, Collège de Bafoussam.
\end{itemize}

\textbf{Expériences professionnelles}
\begin{itemize}
\item Septembre -- décembre 2024 : stage d'aide-comptable, SODECOTON.
\begin{itemize}
\item Saisie des écritures comptables ;
\item Archivage et classement des pièces justificatives.
\end{itemize}
\end{itemize}

\textbf{Vérifications} : dates présentes pour chaque entrée ; ordre inverse respecté (2025, 2023, 2022) ; tâches exprimées par des noms d'action (« saisie », « archivage ») ; aucune information inventée.

\textbf{Conclusion.} Les deux rubriques sont correctement construites : classées du plus récent au plus ancien, datées, rédigées en style télégraphique, elles présentent clairement le profil de gestionnaire de Larissa.
\end{exoresolu}

\coursec{Analyser la situation d'embauche}

\courssub{Méthode 4 : Analyser la situation d'embauche avant de rédiger}

\begin{methode}[Décoder une offre d'emploi pour adapter le CV]
\begin{enumerate}
\item \textbf{Lire l'offre en entier} et souligner : l'intitulé exact du poste, l'entreprise et son secteur, le lieu de travail.
\item \textbf{Relever les exigences} : diplôme exigé, expérience demandée, compétences techniques, qualités personnelles, langues.
\item \textbf{Dresser un tableau de correspondance} : à chaque exigence de l'offre, faire correspondre un élément de son propre parcours.
\item \textbf{Sélectionner et hiérarchiser} : ne retenir dans le CV que ce qui répond à l'offre ; placer en premier l'atout le plus décisif.
\item \textbf{Reprendre le vocabulaire de l'offre} dans l'objectif professionnel et les compétences (le recruteur reconnaît ses propres mots).
\item \textbf{Vérifier la faisabilité} : ne postuler que si le profil correspond réellement ; ne jamais mentir sur un diplôme ou une expérience.
\end{enumerate}
\end{methode}

\begin{exoresolu}[Analyser une offre et préparer le CV — Texte iconique n°4]
\textbf{Énoncé.} « \emph{La société CAMTEL recrute un(e) technicien(ne) de maintenance des réseaux, poste basé à Bafoussam. Profil : baccalauréat technique F4 (électronique/électrotechnique) ou équivalent ; maîtrise des outils informatiques ; sens du travail en équipe et de la rigueur ; permis B souhaité.} »

Kamdem Rodrigue, titulaire d'un Bac F4, a fait un stage de maintenance au Cameroun Télécom Training Center, maîtrise Word et Excel, possède le permis B, joue au football en club et pratique la couture en loisir.

1) Analysez la situation d'embauche (poste, exigences). 2) Indiquez les éléments du parcours de Rodrigue à retenir, à mettre en avant ou à écarter, en justifiant.
\tcblower
\textbf{Solution détaillée.}

\textbf{1. Analyse de la situation d'embauche.}
\begin{center}
\begin{tabularx}{\linewidth}{|l|X|}
\hline
\rowcolor{popblueL} \textbf{Élément de l'offre} & \textbf{Analyse} \\
\hline
Poste & Technicien de maintenance des réseaux \\
\hline
Employeur / lieu & CAMTEL (télécommunications), Bafoussam \\
\hline
Diplôme exigé & Bac technique F4 ou équivalent \\
\hline
Compétences & Outils informatiques \\
\hline
Qualités & Travail en équipe, rigueur \\
\hline
Atout supplémentaire & Permis B souhaité \\
\hline
\end{tabularx}
\end{center}

\textbf{2. Correspondance avec le parcours de Rodrigue.}
\begin{itemize}
\item \textbf{À mettre en avant} : le Bac F4, qui correspond exactement au diplôme exigé — il ouvre la rubrique « Formation » ; le stage de maintenance, qui prouve une première expérience du métier — il ouvre la rubrique « Expériences ».
\item \textbf{À retenir} : la maîtrise de Word et Excel (répond à « outils informatiques ») et le permis B (exigence « souhaitée » : c'est un avantage décisif sur les autres candidats).
\item \textbf{À exploiter comme preuve de qualité} : le football en club, qui illustre le « sens du travail en équipe » demandé par l'offre — à placer dans « Centres d'intérêt ».
\item \textbf{À écarter} : la couture, sans rapport avec le poste ; elle encombrerait le CV sans servir la candidature.
\end{itemize}

\textbf{3. Objectif professionnel adapté} (reprise du vocabulaire de l'offre) : « Obtenir le poste de technicien de maintenance des réseaux à la CAMTEL, en mettant à profit ma formation en électrotechnique et mon expérience de la maintenance. »

\textbf{Conclusion.} Le profil de Rodrigue correspond à toutes les exigences de l'offre, y compris l'atout facultatif du permis B. En sélectionnant et hiérarchisant ses informations, il produit un CV ciblé, c'est-à-dire adapté à cette situation d'embauche précise.
\end{exoresolu}

\coursec{Rédaction et amélioration du CV}

\courssub{Méthode 5 : Rédiger et justifier une démarche — améliorer un CV}

\begin{methode}[Relire, corriger et améliorer un CV]
\begin{enumerate}
\item \textbf{Contrôler le fond} : exactitude des dates et des diplômes, cohérence avec l'offre, absence de trou chronologique injustifié.
\item \textbf{Contrôler la forme} : une page, rubriques titrées en gras, alignements, puces homogènes, photo correcte.
\item \textbf{Contrôler la langue} : style télégraphique partout (pas de « je »), orthographe (accords, majuscules aux noms propres et aux diplômes), vocabulaire professionnel précis.
\item \textbf{Quantifier quand c'est possible} : « saisie de 200 factures par semaine » vaut mieux que « saisie de factures ».
\item \textbf{Faire relire} par un tiers (enseignant, parent) et recueillir ses remarques.
\item \textbf{Justifier chaque modification} : être capable d'expliquer pourquoi telle information a été ajoutée, déplacée ou supprimée (rédiger et justifier une démarche).
\end{enumerate}
\end{methode}

\begin{exoresolu}[Corriger et améliorer un CV défectueux]
\textbf{Énoncé.} Voici un extrait du premier jet du CV de Fotso, candidat à un poste d'électromécanicien :

« \emph{Je m'appelle Fotso Eric, j'ai fait mes études au lycée technique. Je suis trés bon en mécanique et je cherche du travail n'importe ou. Mes expériences : j'ai réparé des moteurs chez mon oncle. J'aime la musique, les voyages, le cinéma, la lecture et les jeux vidéo.} »

Relevez cinq défauts de cet extrait, proposez une version corrigée et justifiez chaque correction.
\tcblower
\textbf{Solution détaillée.}

\textbf{1. Relevé des défauts.}
\begin{enumerate}
\item \textbf{Style rédigé à la première personne} (« je m'appelle », « je suis ») : un CV s'écrit en \cle{style télégraphique} — phrases nominales et verbes d'action, sans « je ».
\item \textbf{Fautes d'orthographe} : « trés » (très), « n'importe ou » (n'importe où).
\item \textbf{Objectif flou} : « je cherche du travail n'importe où » montre l'absence de cible ; l'objectif doit nommer le poste précis.
\item \textbf{Expérience non formalisée} : « chez mon oncle » n'est ni daté ni présenté professionnellement ; il faut des dates, un intitulé de poste et des verbes d'action.
\item \textbf{Centres d'intérêt en liste excessive} : cinq loisirs banals n'apportent rien ; il faut n'en garder qu'un ou deux, révélateurs de qualités utiles.
\end{enumerate}

\textbf{2. Version corrigée.}

\begin{center}
\textbf{FOTSO Eric} — Douala — Tél. : 6 99 00 00 00 — eric.fotso@mail.fr
\end{center}

\textbf{Objectif} : Obtenir un poste d'électromécanicien, en mettant à profit ma formation technique et mon expérience de la réparation de moteurs.

\textbf{Expériences}
\begin{itemize}
\item 2023-2024 : mécanicien assistant, atelier familial de réparation automobile, Douala.
\begin{itemize}
\item Diagnostic des pannes moteur ;
\item Réparation et remplacement de pièces mécaniques.
\end{itemize}
\end{itemize}

\textbf{Centres d'intérêt} : musique (chorale : sens du collectif).

\textbf{3. Justification de la démarche.} J'ai supprimé la première personne pour respecter la convention du genre CV ; corrigé les deux fautes d'orthographe qui discréditeraient immédiatement la candidature ; précisé l'objectif en reprenant l'intitulé du poste visé ; transformé l'expérience informelle en entrée professionnelle datée avec verbes d'action ; réduit les centres d'intérêt à un seul, relié à une qualité (le travail collectif).

\textbf{Conclusion.} Le CV corrigé est ciblé, correct dans sa langue et professionnel dans sa présentation : chaque modification répond à une règle du genre ou à une exigence de la situation d'embauche.
\end{exoresolu}

\courssub{Méthode 6 : Rédiger un CV complet à partir d'une situation d'embauche (synthèse type Bac)}

\begin{methode}[Produire le CV complet en situation d'examen]
\begin{enumerate}
\item \textbf{Lire le sujet deux fois} : noter le poste, l'entreprise, le profil recherché et les informations personnelles fournies par l'énoncé.
\item \textbf{Faire l'inventaire au brouillon} : classer les informations dans les six rubriques (coordonnées, objectif, formation, expériences, compétences, langues/divers).
\item \textbf{Tracer le plan de page} : en-tête, puis rubriques hiérarchisées selon l'atout principal du candidat.
\item \textbf{Rédiger} en style télégraphique, ordre chronologique inverse, verbes d'action, vocabulaire de l'offre.
\item \textbf{Mettre en forme à l'écrit} : titres de rubriques soulignés ou en capitales, puces, alignements (même à la main, la présentation compte).
\item \textbf{Relire} : orthographe, dates, cohérence avec l'offre, une page maximum.
\end{enumerate}
\end{methode}

\begin{exoresolu}[Rédaction complète d'un CV (sujet type Probatoire/Bac technique)]
\textbf{Énoncé.} La société AGRO-INDUSTRIES DU CAMEROUN recrute un(e) assistant(e) de production à Edéa. Profil : baccalauréat technique (génie mécanique ou productique), maîtrise de l'outil informatique, rigueur et ponctualité.

Vous êtes Nkoulou Sandrine, née le 12 mars 2006 à Edéa, titulaire du baccalauréat F2 (2025, lycée technique d'Edéa) et du probatoire F2 (2024). Vous avez effectué un stage de deux mois (juillet-août 2024) à la SOCACAO d'Edéa (surveillance d'une chaîne de production, contrôle qualité). Vous maîtrisez Word et Excel, parlez français couramment et anglais scolairement. Vous pratiquez le handball en club. Rédigez votre CV complet.
\tcblower
\textbf{Solution détaillée.}

\textbf{Étape 1 : analyse de l'offre.} Poste : assistant de production. Exigences : Bac technique (productique/mécanique $\rightarrow$ le F2 convient), informatique, rigueur. Lieu : Edéa — la candidate y réside, c'est un atout.

\textbf{Étape 2 : inventaire et classement.} Coordonnées : Nkoulou Sandrine, Edéa, née le 12/03/2006. Formation : Bac F2 (2025), Probatoire F2 (2024). Expérience : stage SOCACAO (juillet-août 2024). Compétences : Word, Excel. Langues : français courant, anglais scolaire. Divers : handball (esprit d'équipe, ponctualité aux entraînements).

\textbf{Étape 3 : rédaction.}

\begin{center}
\textbf{NKOULOU Sandrine}\\
Edéa — Née le 12 mars 2006 — Tél. : 6 95 00 00 00 — sandrine.nkoulou@mail.fr
\end{center}

\textbf{Objectif professionnel}\\
Obtenir le poste d'assistante de production à AGRO-INDUSTRIES DU CAMEROUN, en mettant à profit ma formation en productique et mon expérience du contrôle qualité.

\textbf{Formation}
\begin{itemize}
\item 2025 : Baccalauréat technique F2 (génie mécanique), Lycée technique d'Edéa.
\item 2024 : Probatoire F2, Lycée technique d'Edéa.
\end{itemize}

\textbf{Expériences professionnelles}
\begin{itemize}
\item Juillet -- août 2024 : stagiaire en production, SOCACAO, Edéa.
\begin{itemize}
\item Surveillance d'une chaîne de production ;
\item Contrôle qualité des produits finis.
\end{itemize}
\end{itemize}

\textbf{Compétences}
\begin{itemize}
\item Informatique : Word, Excel ;
\item Rigueur, ponctualité, sens de l'observation.
\end{itemize}

\textbf{Langues}\\
Français : courant. Anglais : scolaire.

\textbf{Centres d'intérêt}\\
Handball en club (esprit d'équipe, régularité).

\textbf{Étape 4 : vérification.} Une page ; rubriques complètes et titrées ; ordre chronologique inverse respecté ; style télégraphique sans « je » ; vocabulaire de l'offre repris (« production », « contrôle qualité ») ; aucune faute ; le handball justifie la qualité d'équipe annoncée.

\textbf{Conclusion.} Ce CV répond point par point au profil recherché par AGRO-INDUSTRIES DU CAMEROUN : il est complet, ciblé, correctement structuré et rédigé, ce qui maximise les chances d'obtenir un entretien.
\end{exoresolu}

\begin{attention}[Les 5 erreurs qui coûtent des points au Bac]
\begin{enumerate}
\item \textbf{Rédiger le CV à la première personne} (« je suis », « j'ai fait ») : le CV exige le \cle{style télégraphique} — phrases nominales et verbes d'action, sans « je » ni articles inutiles.
\item \textbf{Oublier l'analyse de l'offre} : rédiger un CV générique sans reprendre l'intitulé exact du poste ni le vocabulaire de l'annonce, c'est rater la compétence « adapter le CV à une situation d'embauche », cœur de l'épreuve.
\item \textbf{Classer les diplômes et expériences dans l'ordre chronologique simple} : le plus récent doit figurer \emph{en premier} (ordre chronologique inverse) ; c'est une règle de présentation sanctionnée à la correction.
\item \textbf{Négliger l'orthographe et la présentation} : fautes sur les noms propres, les diplômes (« baccalauréat » sans majuscule en début d'entrée, « trés », « n'importe ou »), rubriques non titrées, puces incohérentes — la forme du texte iconique fait partie de la note.
\item \textbf{Inventer ou gonfler des informations} (faux diplôme, fausse expérience) ou, à l'inverse, omettre des rubriques exigées (coordonnées, objectif, formation) : le correcteur vérifie que toutes les rubriques sont présentes, datées et cohérentes avec les données du sujet.
\end{enumerate}
\end{attention}

\coursec{Compte-rendu et remédiation}

\courssub{Bilan de la séquence : savoirs et savoir-faire mobilisés}

\begin{aretenir}[Synthèse des points clés]
\begin{itemize}
\item Le CV est un \cle{texte iconique} persuasif, normé (une page, structure externe lisible, structure interne complète).
\item La \cle{structure externe} (mise en page, typographie, en-tête) conditionne la première impression du recruteur.
\item La \cle{structure interne} repose sur des rubriques obligatoires (coordonnées, objectif, formation, expériences, compétences, langues, centres d'intérêt/références) rédigées en \cle{style télégraphique} et classées en \cle{ordre chronologique inverse}.
\item L'\cle{analyse de la situation d'embauche} (tableau de correspondance) est l'étape préalable indispensable pour \cle{adapter le CV} : sélection, hiérarchisation, reprise du vocabulaire de l'offre.
\item La \cle{rédaction finale} exige relecture croisée (fond, forme, langue) et justification de chaque choix (compétence « rédiger et justifier une démarche »).
\end{itemize}
\end{aretenir}

\courssub{Exercices de remédiation (entraînement progressif)}

\begin{exercice}[Remédiation 1 : Repérage et classification]
\textbf{Consigne.} Parmi les informations suivantes, classez chacune dans la rubrique CV appropriée (Coordonnées, Objectif, Formation, Expériences, Compétences, Langues, Centres d'intérêt) et mettez-les en ordre chronologique inverse si nécessaire.
\begin{enumerate}
\item Stage de maintenance à la SONEL (juin-juillet 2024).
\item Baccalauréat F1 (2025), Lycée technique de Ngaoundéré.
\item Maîtrise d'AutoCAD et de la lecture de plans.
\item Permis B obtenu en 2024.
\item Pratique du basketball en club (capitaine de l'équipe).
\item Probatoire F1 (2023).
\item Née le 05/01/2007 à Ngaoundéré, Tél. 6 77 88 99 00.
\item Poste visé : technicien en électricité industrielle.
\item Anglais : niveau scolaire ; Français : courant.
\end{enumerate}
\end{exercice}

\begin{corrige}[Exercice Remédiation 1]
\textbf{Coordonnées} : Née le 05/01/2007 à Ngaoundéré, Tél. 6 77 88 99 00.\\
\textbf{Objectif} : Poste visé : technicien en électricité industrielle.\\
\textbf{Formation} : 2025 : Baccalauréat F1, Lycée technique de Ngaoundéré ; 2023 : Probatoire F1.\\
\textbf{Expériences} : Juin-juillet 2024 : stage de maintenance, SONEL.\\
\textbf{Compétences} : Maîtrise d'AutoCAD et de la lecture de plans ; Permis B (2024).\\
\textbf{Langues} : Français : courant ; Anglais : scolaire.\\
\textbf{Centres d'intérêt} : Basketball en club (capitaine : leadership, esprit d'équipe).
\end{corrige}

\begin{exercice}[Remédiation 2 : Adaptation à une offre]
\textbf{Consigne.} Offre : « \emph{Hôtel 4 étoiles à Kribi recrute un(e) réceptionniste. Bac G2 ou Hôtellerie. Anglais courant impératif. Accueil, réservation, gestion des réclamations. Sens du service, présentation soignée.} »
Profil candidat : Bac G2 (2025), stage de 2 mois à la réception de l'Hôtel de la Plage (Kribi) : accueil clients, prise de réservations téléphone/mail. Anglais scolaire. Membre du club théâtre du lycée (expression orale).
\begin{enumerate}
\item Faites le tableau de correspondance offre/profil.
\item Rédigez l'Objectif professionnel adapté.
\item Sélectionnez 2 centres d'intérêt pertinents et justifiez.
\end{enumerate}
\end{exercice}

\begin{corrige}[Exercice Remédiation 2]
\textbf{1. Tableau de correspondance.}
\begin{center}
\begin{tabularx}{\linewidth}{|l|X|}
\hline
\rowcolor{popblueL} \textbf{Exigence de l'offre} & \textbf{Élément du profil} \\
\hline
Bac G2 / Hôtellerie & Bac G2 (2025) \checkmark \\
\hline
Anglais courant & Anglais scolaire (point de vigilance : mentionner « en perfectionnement ») \\
\hline
Accueil, réservation, réclamations & Stage réception : accueil, réservation téléphone/mail \\
\hline
Sens du service, présentation & Théâtre (expression orale, aisance relationnelle) \\
\hline
\end{tabularx}
\end{center}

\textbf{2. Objectif professionnel.} « Obtenir le poste de réceptionniste dans votre établissement, en mettant à profit ma formation en gestion (Bac G2) et mon expérience de l'accueil et de la réservation acquise en stage. »

\textbf{3. Centres d'intérêt.} Théâtre (club lycée) : expression orale, aisance relationnelle, gestion du stress — directement utiles à la réception. \emph{Éviter} : lecture, musique (trop génériques).
\end{corrige}

\begin{exercice}[Remédiation 3 : Correction ciblée — « Jeu des 7 erreurs »]
\textbf{Consigne.} Le CV ci-dessous contient 7 erreurs (forme, fond, langue). Relevez-les, corrigez-les et justifiez chaque correction.

\begin{center}
\begin{tabularx}{\linewidth}{|X|}
\hline
\rowcolor{popblueL} \textbf{CV DE TCHINDA Paul} \\
\hline
\textbf{Objectif} : Je veux travailler dans une grande entreprise. \\
\hline
\textbf{Expériences} \\
2022 : Stage menuiserie chez mon père. \\
2024 : Baccalauréat F3 (Lycée technique). \\
\hline
\textbf{Formation} \\
2023 : Probatoire F3. \\
\hline
\textbf{Compétences} : Je sais bien utiliser Word. Je suis courageux. \\
\hline
\textbf{Langues} : Français bien, Anglais un peu. \\
\hline
\textbf{Loisirs} : Foot, danse, cinéma, lecture, jeux, voyages, cuisine. \\
\hline
\end{tabularx}
\end{center}
\end{exercice}

\begin{corrige}[Exercice Remédiation 3]
\textbf{Erreurs et corrections.}
\begin{enumerate}
\item \textbf{En-tête incomplet} : pas de coordonnées (adresse, téléphone, mail). \textit{Correction} : ajouter en-tête complet.
\item \textbf{Objectif flou} : « Je veux travailler... » (1ère personne, imprécis). \textit{Correction} : « Obtenir un poste de technicien bois/construction, en valorisant ma formation F3 et mon stage. »
\item \textbf{Ordre chronologique inversé non respecté} : Expériences avant Formation, et stage 2022 avant Bac 2024. \textit{Correction} : Formation d'abord (Bac 2024, Probatoire 2023), puis Expériences (Stage 2022).
\item \textbf{Expérience non professionnalisée} : « chez mon père » sans intitulé ni tâches. \textit{Correction} : « Juillet 2022 : Apprenti menuisier, Atelier familial, Douala — Débits de bois, assemblage, finition. »
\item \textbf{Style « je » dans Compétences} : « Je sais..., Je suis... ». \textit{Correction} : « Word (maîtrise) ; Courage, endurance physique. »
\item \textbf{Niveaux de langues vagues} : « bien », « un peu ». \textit{Correction} : « Français : courant ; Anglais : scolaire. »
\item \textbf{Loisirs excessifs et non ciblés} : 7 loisirs. \textit{Correction} : n'en garder qu'un ou deux pertinents (ex. Foot en club : esprit d'équipe ; Bricolage : habileté manuelle).
\end{enumerate}
\end{corrige}

\newpage

\coursec{Exercices}

\courssub{Je vérifie mes connaissances}

\begin{exercice}[QCM : les rubriques du CV]
Pour chaque question, une seule réponse est correcte. Entoure la lettre correspondante.
\begin{enumerate}
\item Un CV est d'abord :
\begin{enumerate}
\item une lettre manuscrite adressée au recruteur ;
\item un document qui présente de façon ordonnée le parcours d'un candidat ;
\item une attestation de travail délivrée par l'employeur.
\end{enumerate}
\item Dans la rubrique « État civil », on trouve :
\begin{enumerate}
\item les diplômes obtenus ;
\item le nom, la date de naissance, les coordonnées ;
\item les loisirs du candidat.
\end{enumerate}
\item La rubrique « Expérience professionnelle » se présente en général :
\begin{enumerate}
\item de l'expérience la plus ancienne à la plus récente ;
\item de l'expérience la plus récente à la plus ancienne ;
\item dans un ordre quelconque.
\end{enumerate}
\item Un bon CV tient idéalement sur :
\begin{enumerate}
\item une page ;
\item trois pages ;
\item dix pages.
\end{enumerate}
\end{enumerate}
\end{exercice}

\begin{exercice}[Vrai ou faux, en justifiant]
Réponds par vrai ou faux, puis justifie ta réponse en une ou deux phrases.
\begin{enumerate}
\item Il est conseillé de mentir sur ses diplômes pour augmenter ses chances d'être recruté.
\item La photo d'identité est obligatoire sur tous les CV.
\item Un CV doit être adapté à l'offre d'emploi à laquelle on répond.
\item Les fautes d'orthographe dans un CV n'ont aucune importance car le recruteur lit vite.
\item Les « centres d'intérêt » peuvent donner au recruteur des informations sur la personnalité du candidat.
\end{enumerate}
\end{exercice}

\begin{exercice}[Définitions]
Définis en deux ou trois phrases chacun des termes suivants :
\begin{enumerate}
\item Curriculum Vitae (CV) ;
\item rubrique ;
\item situation d'embauche ;
\item CV chronologique, CV fonctionnel, CV mixte.
\end{enumerate}
\end{exercice}

\begin{exercice}[Question de cours]
\begin{enumerate}
\item Donne la définition du CV.
\item Cite les cinq rubriques principales d'un CV.
\item Donne trois qualités d'un bon CV et explique brièvement chacune d'elles.
\end{enumerate}
\end{exercice}

\courssub{Je m'entraîne}

\begin{exercice}[Classer les informations dans les rubriques]
Voici des informations relevées sur le brouillon d'un candidat. Recopie le tableau et classe chaque information dans la rubrique appropriée : \emph{État civil}, \emph{Formation}, \emph{Expérience professionnelle}, \emph{Compétences}, \emph{Centres d'intérêt}.

Né le 12 mars 2005 à Bafoussam ; Probatoire technique série F4 (2023) ; stage d'un mois à la menuiserie « Le Bon Bois » (2022) ; maîtrise du logiciel Word ; membre du club de football du quartier ; célibataire, sans enfant ; Baccalauréat technique série F4 (2024) ; aide-comptable chez SABC Douala (juillet--septembre 2024) ; lecture de romans africains ; téléphone : 6 90 00 00 00 ; bonne pratique de l'anglais ; chant choral.

\begin{center}
\begin{tabularx}{\linewidth}{|l|X|}\hline
\rowcolor{popblueL} \textbf{Rubrique} & \textbf{Informations} \\ \hline
État civil & \\ \hline
Formation & \\ \hline
Expérience professionnelle & \\ \hline
Compétences & \\ \hline
Centres d'intérêt & \\ \hline
\end{tabularx}
\end{center}
\end{exercice}

\begin{exercice}[Remettre de l'ordre dans un CV]
Un candidat a rédigé son CV dans le désordre. Voici ses rubriques telles qu'il les a présentées : \emph{Centres d'intérêt} ; \emph{État civil} ; \emph{Expérience professionnelle} ; \emph{Compétences} ; \emph{Formation}.
\begin{enumerate}
\item Remets ces rubriques dans l'ordre habituel d'un CV.
\item Justifie ton classement en deux ou trois phrases.
\end{enumerate}
\end{exercice}

\begin{exercice}[Comparer deux CV]
Deux candidats postulent le même emploi de technicien de surface à l'Hôtel Hilton de Yaoundé.

\textbf{Candidat A} : CV de trois pages, tapé avec trois polices de caractères différentes, cinq fautes d'orthographe, aucune expérience mentionnée, mais une photo récente et des centres d'intérêt détaillés (musique, cinéma, voyages).

\textbf{Candidat B} : CV d'une page, clair et sans faute, deux stages en hôtellerie mentionnés avec dates et lieux, compétences en nettoyage industriel et en anglais élémentaire, centres d'intérêt brefs.

\begin{enumerate}
\item Relève les points forts et les points faibles de chaque CV.
\item Indique, en cinq lignes, le CV que le recruteur retiendra, en justifiant ton choix.
\end{enumerate}
\end{exercice}

\begin{exercice}[Repérer les informations inutiles]
Dans le CV ci-dessous, certaines informations ne devraient pas figurer. Releve-les et explique pourquoi elles sont à supprimer ou à déplacer.

\emph{MANGA Sylvie, née le 3 juin 2004 à Ebolowa, de nationalité camerounaise, groupe sanguin O+, taille 1,65 m, célibataire. Formation : BEPC (2020), Probatoire C (2023), Bac C (2024). Expérience : vendeuse dans la boutique de ma tante (2022). Compétences : je sais faire la cuisine, je suis la meilleure de ma classe, bureautique (Word, Excel). Centres d'intérêt : football, lecture. Mon père est enseignant et ma mère est commerçante.}
\end{exercice}

\courssub{Je me prépare au Bac}

\begin{exercice}[Type Bac : rédiger un CV à partir d'une offre d'emploi]
\textbf{Offre d'emploi.} La société CAMTRANS, entreprise de transit située à Bonanjo (Douala), recrute un(e) \textbf{secrétaire bilingue}. Profil exigé : être titulaire au minimum du Baccalauréat technique (secrétariat-bureautique) ou d'un diplôme équivalent ; maîtriser le français et l'anglais ; connaître les logiciels Word et Excel ; avoir au moins une expérience (stage accepté) ; être disponible immédiatement. Les candidatures sont à déposer au siège de l'entreprise, B.P. 4521 Douala, avant le 30 novembre.

\begin{enumerate}
\item Analyse la situation d'embauche : qui recrute ? quel poste ? quelles qualités et compétences sont exigées ? (4 réponses attendues)
\item Rédige, sur une page maximum, le CV d'un candidat correspondant au profil exigé. Tu respecteras les cinq rubriques principales et tu adapteras les informations à l'offre.
\item Justifie en quelques lignes le choix du modèle de CV que tu as adopté (chronologique, fonctionnel ou mixte).
\end{enumerate}
\end{exercice}

\begin{exercice}[Type Bac : corriger un CV défectueux]
Voici le CV de M. ONANA Jean, candidat à un poste de magasinier dans une entreprise de Yaoundé. Ce document contient volontairement des défauts : fautes, désordre des rubriques, informations inutiles, mensonges.

\emph{Curiculum Vitaé. Centres d'intérêt : je joue au foot, je danse le makossa, j'aime dormir. Formation : Bac technique F4 obtenu en 2024 avec mention très bien (alors qu'il l'a obtenu avec mention passable) ; j'ai redoublé la classe de seconde. État civil : ONANA Jean, né le 15/08/2003 à Yaoundé, taille 1,78 m, poids 82 kg, religion catholique, téléphone 6 99 00 00 00. Expérience : magasinier stagiaire à la SODECOTON (3 mois, 2023) ; j'ai été chef d'entreprise pendant 5 ans. Compétences : conduite de chariot élévateur, gestion des stocks, je parle ewondo à la maison.}

\begin{enumerate}
\item Releve dix erreurs ou défauts de ce CV, en indiquant pour chacun la nature de la faute (orthographe, rubrique mal placée, information inutile, mensonge, etc.).
\item Propose la version corrigée et complète de ce CV, présentée dans l'ordre correct des rubriques.
\item Explique en cinq lignes pourquoi un recruteur peut écarter un CV contenant des mensonges.
\end{enumerate}
\end{exercice}

\begin{exercice}[Type Bac : projet personnel de CV]
Tu es élève en classe de Terminale technique. Projete-toi deux ans après l'obtention de ton Baccalauréat technique.

\begin{enumerate}
\item Imagine ta situation dans deux ans : diplômes obtenus, formation suivie ou emploi occupé, stages effectués, compétences acquises.
\item Rédige ton propre CV complet (une page maximum), en respectant les cinq rubriques principales et en veillant à l'orthographe et à la présentation.
\item Rédige un court paragraphe (cinq à huit lignes) justifiant le choix du modèle de CV que tu as adopté : chronologique, fonctionnel ou mixte.
\item Indique à quel type d'offre d'emploi ce CV pourrait répondre, en précisant le poste visé et l'entreprise ou le secteur d'activité.
\end{enumerate}
\end{exercice}

\newpage

\coursec{Corrigés des exercices}

\coursec{Leçon N06 : Produire un Curriculum Vitae (CV)}
\courssub{Corrigés des exercices d'application et de type Baccalauréat}

\begin{corrige}[Exercice 1 — QCM : les rubriques du \cle{CV}]
\begin{enumerate}
\item \textbf{Réponse b} : un \cle{CV} est un document qui présente de façon ordonnée le parcours d'un candidat (formation, expérience, compétences). Ce n'est ni une lettre manuscrite (c'est la lettre de motivation) ni une attestation de travail.
\item \textbf{Réponse b} : l'état civil regroupe le nom et les prénoms, la date et le lieu de naissance, la situation de famille et les coordonnées (adresse, téléphone, courriel). Les diplômes relèvent de la rubrique « Formation », les loisirs des « Centres d'intérêt ».
\item \textbf{Réponse b} : l'expérience professionnelle se présente de la plus récente à la plus ancienne (ordre \cle{antichronologique}), car le recruteur s'intéresse d'abord à ce que le candidat sait faire aujourd'hui.
\item \textbf{Réponse a} : un bon \cle{CV} tient idéalement sur une page : le recruteur lit très vite et doit trouver l'essentiel en un coup d'œil.
\end{enumerate}
\end{corrige}

\begin{corrige}[Exercice 2 — Vrai ou faux, en justifiant]
\begin{enumerate}
\item \textbf{Faux.} Mentir sur ses diplômes est une faute grave : le mensonge peut être découvert lors de la vérification des pièces ou à l'essai, et entraîner le refus de la candidature, voire un licenciement pour faute. L'honnêteté est une règle absolue du \cle{CV}.
\item \textbf{Faux.} La photo d'identité n'est pas obligatoire : elle est facultative et dépend de l'usage ou de la demande de l'employeur. Si on la met, elle doit être récente, nette et sérieuse.
\item \textbf{Vrai.} Un \cle{CV} efficace est adapté à l'offre d'emploi : on met en avant les diplômes, expériences et compétences qui correspondent au profil exigé par le recruteur.
\item \textbf{Faux.} Les fautes d'orthographe sont rédhibitoires : elles donnent une image de négligence et d'incompétence, et beaucoup de recruteurs écartent immédiatement un \cle{CV} fautif, même lu rapidement.
\item \textbf{Vrai.} Les centres d'intérêt renseignent sur la personnalité : la pratique d'un sport collectif suggère l'esprit d'équipe, la lecture la curiosité, une activité associative le sens de l'engagement.
\end{enumerate}
\end{corrige}

\begin{corrige}[Exercice 3 — Définitions]
\begin{enumerate}
\item \textbf{\cle{Curriculum Vitae (CV)}} : expression latine signifiant « course de la vie ». C'est un document bref (une page en général) qui présente de manière claire et ordonnée l'état civil, la formation, l'expérience professionnelle et les compétences d'un candidat à un emploi.
\item \textbf{\cle{Rubrique}} : chacune des parties titrées qui structurent le \cle{CV} (État civil, Formation, Expérience professionnelle, Compétences, Centres d'intérêt). Elle regroupe des informations de même nature sous un intitulé visible.
\item \textbf{\cle{Situation d'embauche}} : l'ensemble des circonstances concrètes d'un recrutement : l'entreprise qui recrute, le poste proposé, le profil exigé (diplômes, compétences, expérience) et les conditions de la candidature. Son analyse permet d'adapter le \cle{CV} à l'offre.
\item \textbf{\cle{CV chronologique}} : \cle{CV} qui présente la formation et l'expérience dans l'ordre du temps, de la plus récente à la plus ancienne. \textbf{\cle{CV fonctionnel}} : \cle{CV} organisé par domaines de compétences, utile quand le parcours est irrégulier. \textbf{\cle{CV mixte}} : \cle{CV} qui combine les deux : il met en avant les compétences puis retrace le parcours chronologiquement.
\end{enumerate}
\end{corrige}

\begin{corrige}[Exercice 4 — Question de cours]
\begin{enumerate}
\item \textbf{Définition du \cle{CV}} : le Curriculum Vitae est un document écrit, bref et structuré, qui présente le parcours d'un candidat (identité, formation, expérience, compétences) en vue d'obtenir un emploi ou un stage.
\item \textbf{Les cinq rubriques principales} : l'état civil (ou état civil et coordonnées) ; la formation (diplômes) ; l'expérience professionnelle ; les compétences (langues, informatique, savoir-faire techniques) ; les centres d'intérêt.
\item \textbf{Trois qualités d'un bon \cle{CV}} :
\begin{itemize}
\item \emph{La clarté} : une présentation aérée, des rubriques visibles, une seule police de caractères, pour que le recruteur trouve l'information en un coup d'œil.
\item \emph{La concision} : une page maximum, des phrases courtes, aucune information inutile, car le recruteur consacre peu de temps à chaque dossier.
\item \emph{La sincérité} : des informations exactes et vérifiables, car tout mensonge découvert discrédite définitivement le candidat.
\end{itemize}
(On pouvait aussi citer : l'orthographe irréprochable, l'adaptation à l'offre d'emploi.)
\end{enumerate}
\end{corrige}

\begin{corrige}[Exercice 5 — Classer les informations dans les rubriques]
\begin{center}
\begin{tabularx}{\linewidth}{|l|X|}\hline
\rowcolor{popblueL} \textbf{Rubrique} & \textbf{Informations} \\ \hline
État civil & Né le 12 mars 2005 à Bafoussam ; célibataire, sans enfant ; téléphone : 6 90 00 00 00 \\ \hline
Formation & Probatoire technique série F4 (2023) ; Baccalauréat technique série F4 (2024) \\ \hline
Expérience professionnelle & Aide-comptable chez SABC Douala (juillet--septembre 2024) ; stage d'un mois à la menuiserie « Le Bon Bois » (2022) \\ \hline
Compétences & Maîtrise du logiciel Word ; bonne pratique de l'anglais \\ \hline
Centres d'intérêt & Membre du club de football du quartier ; lecture de romans africains ; chant choral \\ \hline
\end{tabularx}
\end{center}
\textbf{Remarques} : dans l'expérience professionnelle, on place l'emploi le plus récent (SABC, 2024) avant le stage le plus ancien (2022) ; dans la formation, le Baccalauréat (2024) précède le Probatoire (2023). Le numéro de téléphone relève des coordonnées, donc de l'état civil.
\end{corrige}

\begin{corrige}[Exercice 6 — Remettre de l'ordre dans un \cle{CV}]
\begin{enumerate}
\item \textbf{Ordre habituel} : 1. État civil ; 2. Formation ; 3. Expérience professionnelle ; 4. Compétences ; 5. Centres d'intérêt.
\item \textbf{Justification} : l'état civil vient en tête car il identifie le candidat et permet de le contacter. Viennent ensuite les informations décisives pour le recruteur : d'abord la formation (les diplômes), puis l'expérience professionnelle, qui prouve la pratique du métier. Les compétences complètent le profil en précisant les savoir-faire (langues, informatique). Les centres d'intérêt ferment le \cle{CV} car ils sont secondaires : ils donnent seulement un aperçu de la personnalité.
\end{enumerate}
\end{corrige}

\begin{corrige}[Exercice 7 — Comparer deux \cle{CV}]
\begin{enumerate}
\item \textbf{Candidat A.} Points faibles : \cle{CV} trop long (trois pages), présentation désordonnée (trois polices), cinq fautes d'orthographe, aucune expérience professionnelle. Points forts : photo récente, centres d'intérêt détaillés — mais ceux-ci ne compensent pas l'absence d'expérience et les fautes.
\textbf{Candidat B.} Points forts : \cle{CV} d'une page, clair et sans faute ; deux stages en hôtellerie datés et localisés, donc en rapport avec le poste ; compétences utiles (nettoyage industriel, anglais élémentaire). Point faible : centres d'intérêt très brefs, mais cela reste secondaire.
\item \textbf{Choix du recruteur} : le recruteur retiendra le \textbf{candidat B}. Son \cle{CV} respecte les qualités essentielles du genre : clarté, concision, correction orthographique. Surtout, il est adapté à la \cle{situation d'embauche} : les stages en hôtellerie et les compétences en nettoyage industriel correspondent exactement au poste de technicien de surface. Le candidat A, malgré sa photo et ses loisirs, donne une image de négligence et ne prouve aucune aptitude au métier.
\end{enumerate}
\end{corrige}

\begin{corrige}[Exercice 8 — Repérer les informations inutiles]
\begin{itemize}
\item \textbf{« Groupe sanguin O+ »} : information médicale sans rapport avec l'emploi ; à supprimer.
\item \textbf{« Taille 1,65 m »} : caractéristique physique inutile (sauf métiers spécifiques) ; à supprimer.
\item \textbf{« Mon père est enseignant et ma mère est commerçante »} : la profession des parents ne concerne pas le candidat ; à supprimer.
\item \textbf{« Je sais faire la cuisine »} : compétence domestique sans lien avec un emploi de bureau ou de vente ; à supprimer (sauf candidature en restauration).
\item \textbf{« Je suis la meilleure de ma classe »} : appréciation subjective et prétentieuse, invérifiable ; à supprimer ou à remplacer par un fait objectif (ex. : « major de promotion au BEPC » si c'est exact).
\item \textbf{« Vendeuse dans la boutique de ma tante »} : la formulation est familiale ; à reformuler objectivement : « vendeuse dans un commerce de détail, Ebolowa (2022) ».
\end{itemize}
\textbf{Règle générale} : ne figurent dans un \cle{CV} que des informations objectives, vérifiables et utiles au poste visé.
\end{corrige}

\begin{corrige}[Exercice 9 — Type Bac : rédiger un \cle{CV} à partir d'une offre d'emploi]
\begin{enumerate}
\item \textbf{Analyse de la \cle{situation d'embauche}} :
\begin{itemize}
\item Qui recrute ? La société CAMTRANS, entreprise de transit, Bonanjo, Douala (B.P. 4521 Douala).
\item Quel poste ? Secrétaire bilingue.
\item Qualités et compétences exigées ? Baccalauréat technique (secrétariat-bureautique) ou équivalent ; maîtrise du français et de l'anglais ; connaissance de Word et Excel ; au moins une expérience (stage accepté) ; disponibilité immédiate.
\item Conditions ? Dossier à déposer avant le 30 novembre.
\end{itemize}
\item \textbf{Modèle de \cle{CV} attendu} (exemple) :

\textbf{ÉTAT CIVIL} — MBAPPE Laure, née le 5 mai 2005 à Douala, célibataire. Tél. : 6 94 00 00 00 ; courriel : \texttt{laure.mbappe@mail.com} ; Bonabéri, Douala.

\textbf{FORMATION} — 2024 : Baccalauréat technique, spécialité secrétariat-bureautique, Lycée technique de Douala-Bassa. 2023 : Probatoire technique, même spécialité.

\textbf{EXPÉRIENCE PROFESSIONNELLE} — Sept.--nov. 2024 : secrétaire stagiaire, cabinet FODEC, Douala (accueil, frappe de courriers, classement). Juillet 2023 : stage de découverte en entreprise de transit, Douala.

\textbf{COMPÉTENCES} — Français : langue de scolarité, excellente maîtrise ; anglais : bon niveau (lu, écrit, parlé) ; bureautique : Word, Excel ; dactylographie rapide ; disponible immédiatement.

\textbf{CENTRES D'INTÉRÊT} — Lecture, chant choral, volley-ball.

\item \textbf{Justification du modèle} : le \textbf{\cle{CV chronologique}} convient ici car le parcours de la candidate est régulier et récent : les diplômes et le stage, présentés du plus récent au plus ancien, montrent immédiatement au recruteur qu'elle possède le profil exigé.
\end{enumerate}
\textbf{Barème indicatif (sur 20)} : analyse de la situation 4 pts ; respect des cinq rubriques 5 pts ; adaptation à l'offre (bilinguisme, Word/Excel, expérience) 4 pts ; clarté et présentation 3 pts ; orthographe 2 pts ; justification du modèle 2 pts.
\textbf{Points de vigilance} : ne pas inventer de diplôme supérieur à l'offre ; mentionner explicitement le bilinguisme et les logiciels exigés ; tenir sur une page ; ordre antichronologique dans « Formation » et « Expérience ».
\end{corrige}

\begin{corrige}[Exercice 10 — Type Bac : corriger un \cle{CV} défectueux]
\begin{enumerate}
\item \textbf{Relevé de dix erreurs} :
\begin{enumerate}
\item « Curiculum Vitaé » : fautes d'orthographe (écrire \emph{Curriculum Vitae}).
\item Rubrique « Centres d'intérêt » placée en tête : désordre des rubriques (elle vient en dernier).
\item « J'aime dormir » : information inutile et dévalorisante dans les centres d'intérêt.
\item « Mention très bien » au lieu de « passable » : mensonge sur le diplôme.
\item « J'ai redoublé la classe de seconde » : information inutile et défavorable, à ne pas mentionner.
\item « Taille 1,78 m, poids 82 kg » : informations physiques inutiles dans l'état civil.
\item « Religion catholique » : information personnelle sans rapport avec l'emploi, à supprimer.
\item « J'ai été chef d'entreprise pendant 5 ans » : mensonge invraisemblable (incohérent avec l'âge et le stage de 2023).
\item « Je parle ewondo à la maison » : formulation familière et information non professionnalisée ; à reformuler (« langues : français, anglais, ewondo ») dans les compétences.
\item Style à la première personne et phrases familières (« je joue au foot, je danse le makossa ») : registre inadapté ; le \cle{CV} exige des formulations brèves et objectives.
\end{enumerate}
\item \textbf{Version corrigée} :

\textbf{CURRICULUM VITAE}

\textbf{ÉTAT CIVIL} — ONANA Jean, né le 15 août 2003 à Yaoundé. Tél. : 6 99 00 00 00.

\textbf{FORMATION} — 2024 : Baccalauréat technique série F4, mention passable.

\textbf{EXPÉRIENCE PROFESSIONNELLE} — 2023 : magasinier stagiaire à la SODECOTON (3 mois) : réception et rangement des marchandises, tenue des fiches de stock.

\textbf{COMPÉTENCES} — Conduite de chariot élévateur ; gestion des stocks ; langues : français, ewondo.

\textbf{CENTRES D'INTÉRÊT} — Football ; danse traditionnelle et musique (makossa).

\item \textbf{Pourquoi le recruteur écarte un \cle{CV} mensonger} : le mensonge détruit la confiance, qui est la base du contrat de travail. Les fausses mentions sont facilement vérifiables (relevés de notes, références d'employeurs). Un candidat qui ment sur son diplôme ou son expérience laisse supposer qu'il mentira dans le travail ; il est donc écarté, et un mensonge découvert après l'embauche peut justifier un licenciement.
\end{enumerate}
\textbf{Barème indicatif (sur 20)} : relevé des dix erreurs 8 pts ; \cle{CV} corrigé (rubriques, ordre, exactitude) 8 pts ; explication sur le mensonge 4 pts.
\textbf{Points de vigilance} : chaque erreur doit être accompagnée de sa nature ; le \cle{CV} corrigé ne doit contenir que des informations vraies, utiles et objectivement formulées.
\end{corrige}

\begin{corrige}[Exercice 11 — Type Bac : projet personnel de \cle{CV}]
\textbf{Corrigé type (exemple à adapter au parcours de chaque élève).}
\begin{enumerate}
\item \textbf{Projection} : dans deux ans, j'aurai obtenu mon Baccalauréat technique série F4 (2026), puis suivi une formation d'un an en gestion logistique au lycée technique, avec un stage de deux mois dans une entreprise de transport. J'aurai acquis des compétences en bureautique (Word, Excel) et en gestion des stocks.
\item \textbf{\cle{CV} rédigé} :

\textbf{ÉTAT CIVIL} — EKALLA Brice, né le 20 janvier 2007 à Douala, célibataire. Tél. : 6 91 00 00 00 ; courriel : \texttt{brice.ekalla@mail.com} ; Deïdo, Douala.

\textbf{FORMATION} — 2027 : attestation de formation en gestion logistique, Lycée technique de Douala. 2026 : Baccalauréat technique série F4. 2025 : Probatoire technique série F4.

\textbf{EXPÉRIENCE PROFESSIONNELLE} — Juillet--août 2027 : stage de magasinier, société de transit Bolloré, Douala (inventaires, saisie informatique). 2026 : aide-vendeur, marché central de Douala (vacances).

\textbf{COMPÉTENCES} — Gestion des stocks ; bureautique : Word, Excel ; français courant, anglais scolaire ; sens de l'organisation et ponctualité.

\textbf{CENTRES D'INTÉRÊT} — Football (club de quartier), lecture, musique.

\item \textbf{Justification du modèle} : j'ai choisi le \textbf{\cle{CV chronologique}} parce que mon parcours est récent, court et sans interruption : présenter mes diplômes et mon stage du plus récent au plus ancien permet au recruteur de voir immédiatement ma qualification la plus élevée. Le modèle fonctionnel serait utile à un candidat au parcours irrégulier, ce qui n'est pas mon cas ; le modèle mixte conviendrait si j'avais déjà de nombreuses compétences à mettre en avant avant mon expérience.
\item \textbf{Offre visée} : ce \cle{CV} peut répondre à une offre d'emploi de \textbf{magasinier} ou d'\textbf{aide-logisticien} dans une entreprise de transit ou de transport de la zone portuaire de Douala (secteur logistique), où le Bac technique F4, la formation en gestion des stocks et le stage effectué correspondent au profil demandé.
\end{enumerate}
\textbf{Barème indicatif (sur 20)} : cohérence de la projection 3 pts ; cinq rubriques complètes et ordonnées 6 pts ; adaptation et vraisemblance des informations 4 pts ; justification du modèle 4 pts ; orthographe et présentation 3 pts.
\textbf{Points de vigilance} : les dates projetées doivent être cohérentes entre elles ; aucune information inutile (taille, religion, famille) ; formulations brèves et objectives ; une page maximum.
\end{corrige}

\newpage

\coursec{Fiche bilan}

\begin{popfigure}
\begin{tikzpicture}[mindmap, grow cyclic, every node/.style={concept, font=\small}, concept color=popblue, level 1/.append style={level distance=3.6cm, sibling angle=51.4}, level 2/.append style={level distance=2.2cm}]
\node[concept, font=\small\bfseries] {Le Curriculum Vitae (CV)}
  child[concept color=poporange] { node {Définition et fonctions}
    child { node[concept, font=\scriptsize] {Document de présentation de soi} }
    child { node[concept, font=\scriptsize] {Décrocher un entretien} }
  }
  child[concept color=popgreen] { node {Structure externe}
    child { node[concept, font=\scriptsize] {Une page, claire, aérée} }
    child { node[concept, font=\scriptsize] {Mise en page soignée} }
  }
  child[concept color=poppurple] { node {Structure interne : les rubriques}
    child { node[concept, font=\scriptsize] {État civil, formation, expérience} }
    child { node[concept, font=\scriptsize] {Compétences, centres d'intérêt} }
  }
  child[concept color=popgold] { node {Analyser la situation d'embauche}
    child { node[concept, font=\scriptsize] {Lire l'offre d'emploi} }
    child { node[concept, font=\scriptsize] {Repérer le profil recherché} }
  }
  child[concept color=poppink] { node {Adapter le CV}
    child { node[concept, font=\scriptsize] {Mettre en avant l'adéquation profil/poste} }
  }
  child[concept color=popblue!70] { node {Rédiger et améliorer}
    child { node[concept, font=\scriptsize] {Relecture, correction, reformulation} }
  };
\end{tikzpicture}
\legende{Carte mentale de la leçon : les six idées forces pour produire un CV efficace.}
\end{popfigure}

\begin{bilan}
\textbf{Définitions clés}
\begin{itemize}
  \item Le \cle{CV} (Curriculum Vitae, « déroulement de la vie ») : document écrit, bref et structuré, qui présente l'identité, la formation, l'expérience et les compétences d'un candidat.
  \item La \cle{situation d'embauche} : ensemble formé par l'offre d'emploi, l'employeur, le poste et le profil recherché.
  \item La \cle{structure externe} : l'aspect matériel du CV (format, mise en page, titre, photo éventuelle).
  \item La \cle{structure interne} : le contenu organisé en rubriques.
\end{itemize}

\textbf{Les rubriques du CV (à connaître par cœur)}
\begin{center}
\begin{tabularx}{\linewidth}{|l|X|}
\hline
\rowcolor{popblueL} \textbf{Rubrique} & \textbf{Contenu attendu} \\
\hline
État civil / Identité & Nom, prénom, date et lieu de naissance, adresse, téléphone, courriel, nationalité, situation matrimoniale. \\
\hline
Objectif / Titre & Le poste visé, en une phrase courte. \\
\hline
Formation / Études & Diplômes obtenus, du plus récent au plus ancien (ordre antichronologique), avec établissements et années. \\
\hline
Expérience professionnelle & Stages et emplois : poste, entreprise, période, tâches principales. \\
\hline
Compétences & Langues, informatique, savoir-faire techniques liés au poste. \\
\hline
Centres d'intérêt & Activités valorisantes (sport, lecture, engagement associatif). \\
\hline
Références & Personnes pouvant témoigner du candidat (facultatif). \\
\hline
\end{tabularx}
\end{center}

\textbf{Méthodes express}
\begin{itemize}
  \item \textbf{Analyser l'offre} : souligner le poste, les qualités et les diplômes exigés ; dégager le profil recherché.
  \item \textbf{Adapter} : sélectionner et ordonner ses rubriques selon le poste ; mettre en premier ce qui correspond à l'offre.
  \item \textbf{Rédiger} : phrases brèves ou simples énumérations, verbes d'action, ordre antichronologique, une seule page.
  \item \textbf{Améliorer} : relire (orthographe, cohérence des dates), faire relire, simplifier, harmoniser la présentation.
\end{itemize}
\end{bilan}

\begin{aretenir}[Checklist avant le Bac]
\begin{itemize}
  \item $\square$ Je sais définir le CV et citer ses fonctions (se présenter, obtenir un entretien).
  \item $\square$ Je connais les qualités de la structure externe : une page, clarté, aération, sobriété.
  \item $\square$ Je cite sans hésiter les rubriques de la structure interne.
  \item $\square$ Je présente l'état civil de façon complète et exacte.
  \item $\square$ Je classe les diplômes et expériences du plus récent au plus ancien.
  \item $\square$ Je sais analyser une offre d'emploi : poste, profil, qualités exigées.
  \item $\square$ J'adapte mon CV à la situation d'embauche (rubriques mises en avant).
  \item $\square$ J'utilise des formulations brèves et des verbes d'action.
  \item $\square$ Je vérifie l'orthographe, la grammaire et la cohérence des dates.
  \item $\square$ Je soigne la présentation : titre, alignements, même police, pas de ratures.
  \item $\square$ Je fais relire mon CV avant de le remettre.
\end{itemize}
\end{aretenir}

\findecours

\end{document}
